\documentclass[11pt]{article}
\usepackage[margin=1.1in]{geometry}
\usepackage[T1]{fontenc}
\usepackage{lmodern}
\usepackage{amsmath,amssymb,amsthm,microtype}
\usepackage[colorlinks=true,linkcolor=blue,citecolor=blue,urlcolor=blue]{hyperref}
\newtheorem{theorem}{Theorem}
\newtheorem{lemma}{Lemma}[section]
\newtheorem{proposition}[lemma]{Proposition}
\newtheorem{corollary}[lemma]{Corollary}
\theoremstyle{definition}
\newtheorem{definition}[lemma]{Definition}
\newtheorem{remark}[lemma]{Remark}
\DeclareMathOperator{\Bin}{Bin}
\DeclareMathOperator{\Bern}{Bern}
\newcommand{\prob}{\mathbb P}
\newcommand{\ex}{\mathbb E}
\newcommand{\whp}{with high probability}
\newcommand{\step}[1]{\par\medskip\noindent\textit{#1}\enspace}
\setlength{\emergencystretch}{2em}

\hypersetup{pdftitle={Sharp thresholds for anonymous recovery in Hamming graphs},pdfauthor={Brandon Li}}
\title{Sharp thresholds for anonymous recovery in Hamming graphs}
\author{Brandon Li\\[0.35em]
 \normalsize Great Neck, New York, USA\\[0.2em]
 \normalsize \href{mailto:brandon.li.gn@gmail.com}{\nolinkurl{brandon.li.gn@gmail.com}}}
\date{}
\begin{document}
\maketitle
\begin{abstract}
Suppose a set of landmarks reports its distances to an unknown query in a Hamming graph, but the readings are unlabeled, so one observes only the distance histogram rather than which landmark produced which distance. We study how many randomly chosen landmarks are needed to identify every query in a given family simultaneously. For the binary cube, we determine the random sample size needed to distinguish all faces, both without errors and under adversarial reading errors. The leading threshold is \(e^{Cn^{1/3}}\) with \(C=3.4417465\ldots\), and an adversarial error budget of \(e^{Bn^{1/3}}\) multiplies this exponential scale by \(e^{2Bn^{1/3}}\). At the next scale, the critical polynomial factors are \(n^{1/2}\) in the noiseless model and \(n^{-5/6}\) under an exponential error budget. We determine the sharp multiplicative cutoffs in both regimes and prove that recovery fails at equality. Above the threshold, for any fixed strict excess in the leading exponent, almost every remaining failure is caused by faces that fix only one or two coordinates; a simple filter removes those configurations and doubles the logarithmic decay rate of the failure probability.

For general Hamming graphs over a fixed alphabet, we establish stretched-exponential bounds for the joint recovery of all coordinate boxes and for the classical vertex, outer-vertex, and edge multiset dimensions. For each of these three classical query families, we prove the exact sparse Bernoulli threshold, including the sharp prefactor and failure at equality. For vertices and edges, we also locate the narrower critical windows and obtain the limiting recovery probabilities over every fixed alphabet. The corresponding collision counts converge to Poisson limits, with a factor-of-two doubling in the binary case caused by complement symmetry. All of these classical thresholds and finer limits carry over to a uniformly chosen landmark set of exactly the specified size.
\end{abstract}

\noindent\textit{Keywords}: Hamming graphs, hypercube, multiset dimension, resolving sets, random landmarks, sharp thresholds, Poisson approximation

\noindent\textit{Mathematics Subject Classification}: 05C12, 05C80, 60C05

\section{Anonymous measurements}
In many localization and sensor-network problems, a collection of landmarks measures distances to an unknown target, and the observer must identify the target from those readings. When each sensor is labeled, a measurement is an ordered vector with one coordinate for each landmark. In the anonymous setting, landmark identities are stripped away, so the observer receives only the multiset, or histogram, of distances, with one coordinate for each possible distance value. For labeled measurements, the minimum number of landmarks is the metric dimension of the graph, introduced independently by Slater \cite{Slater} and by Harary and Melter \cite{HM}; the term landmark comes from Khuller, Raghavachari and Rosenfeld \cite{KRR}. Resolving sets for anonymous measurements are studied under the name of multiset dimension, introduced by Saenpholphat \cite{Saenpholphat} and independently by Simanjuntak, Siagian and Vetr\'ik \cite{SSV}; see the recent survey \cite{FKKY}. We ask how many landmarks are needed so that a single landmark set distinguishes every query in a prescribed family simultaneously.

We work in the Hamming graph \(H(n,q)\), whose vertices are words in \([q]^n\) and whose edges join pairs of words that differ in exactly one coordinate. For any nonempty set \(T\) of query vertices (such as a single vertex, an ordinary edge, or a coordinate box), the distance from a vertex to \(T\) is \(d(v,T)=\min_{t\in T}d_H(v,t)\). A landmark set \(S\subseteq[q]^n\) assigns to \(T\) the distance histogram
\[
h_S(T)_r=\#\{v\in S:d(v,T)=r\},\qquad 0\le r\le n.
\]
A coordinate box is a Cartesian product \(\prod_{i=1}^n T_i\) in which every factor \(T_i\subseteq[q]\) is nonempty. When \(q=2\), the coordinate boxes are precisely the faces of the binary cube, and the codimension of a face is the number of coordinates whose bit values it prescribes.

To illustrate how distinct faces can collide under anonymous measurements, consider the three landmarks \(000,101,111\) in the three-dimensional binary cube. The face fixing \(x_1=1,x_3=0\) and the face fixing \(x_1=0,x_3=1\) both produce the histogram \((0,3,0,0)\), because each landmark matches exactly one of the two prescribed bits of either face and therefore lies at distance one from both.

Throughout the paper, landmark sets are unweighted and contain no repeated vertices, and the empty set does not resolve. In our random model, each vertex of the Hamming graph is included independently in \(S\) with probability \(m/q^n\), so that \(m=\ex|S|\) is the specified expected size. The alphabet size \(q\ge2\) is held fixed as the dimension grows, and all logarithms are natural. An event holds \emph{with high probability} if its probability tends to one as \(n\to\infty\).

\paragraph{Main contributions.}
This paper establishes four main results. First, we determine the sharp random threshold for distinguishing all \(3^n\) faces of the binary cube, both in the noiseless model and under adversarial reading errors, including the critical polynomial multipliers and exact multiplicative cutoffs at the second scale. Second, we identify which faces drive the rare failures that persist above the threshold, showing that they differ from the faces that govern the threshold itself and that a simple filter doubles the logarithmic decay rate of the failure probability. Third, for every fixed alphabet, we prove sharp sparse thresholds for vertices, outer vertices, and edges, and for vertices and edges we also find the finer critical windows and the limiting recovery probabilities. These thresholds, critical windows and limiting probabilities also hold when the landmarks form a uniformly chosen set of a fixed size; the binary face results in the first two parts are proved for independent selection of landmarks. Fourth, we determine the deterministic order \(\exp(\Theta(n^{1/3}))\) for the classical vertex, outer-vertex, and edge multiset dimensions of Hamming graphs, as well as for joint resolution of all coordinate boxes.

\section{Binary faces: thresholds and rare failures}
To model adversarial noise, we call two distance histograms confusable at error budget \(E\) if an adversary can make them coincide by replacing at most \(E\) readings in each histogram. Because every query produces a histogram whose total mass equals the number of landmarks, two histograms of the same total mass are confusable at budget \(E\) if and only if half their \(\ell^1\) distance is at most \(2E\).

Permitting deletions alongside replacements within the same total budget does not change which pairs are confusable. Indeed, for two corrupted outputs to coincide they must have the same total mass, which forces the adversary to delete the same number \(d\) of readings from each side. Appending \(d\) identical common readings to both outputs converts those deletions into replacements without exceeding the budget.

\begin{theorem}[Binary box first moment]\label{thm:box}
Suppose \(q=2\) and \(m=e^{A n^{1/3}}\), where \(A>0\). Let \(E=\lfloor e^{B n^{1/3}}\rfloor\) for a fixed \(B\ge0\), or set \(E=0\) and \(B=0\) in the noiseless case. Put \(D=A-2B\) and
\[
 h_{\rm face}(\rho)=-\rho\log\rho-(1-\rho)\log(1-\rho)+\rho\log2,
 \qquad
 a_\rho(D)=\tfrac23\sqrt{\rho/2}\,\max(D,0)^{3/2}.
\]
Use the continuous convention \(0\log 0=0\). For \(r/n\to\rho\in(0,1]\), let \(Z_r\) count unordered confusable pairs of codimension-\(r\) faces. Then
\[
\lim_{n\to\infty} n^{-1}\log\ex Z_r
=\max\{2h_{\rm face}(\rho)-a_\rho(D),\ h_{\rm face}(\rho)-\tfrac12a_\rho(D)\}.
\]
The expected number of confusable pairs among all faces tends to zero for
\[
D>\max_{0<\rho\le1}\Bigl(3h_{\rm face}(\rho)/\sqrt{\rho/2}\Bigr)^{2/3}=3.4417465\ldots
\]
and grows exponentially for \(D\) below this constant. In particular all faces are distinguished \whp{} above the cutoff. Theorems~\ref{thm:face-wHP} and~\ref{thm:noisy-face} give matching probability cutoffs, with distinct equality conclusions for noiseless and exponentially noisy sampling.
\end{theorem}

\subsection{Recovery probabilities}
Let \(C\) denote the exact threshold constant from Theorem~\ref{thm:box}. While Theorem~\ref{thm:box} locates the first-moment transition, Theorems~\ref{thm:face-wHP} and~\ref{thm:noisy-face} establish the actual probability thresholds for recovery. For a fixed parameter \(A>0\) and expected landmark count \(m=e^{An^{1/3}}\), they yield the following phase diagram:
\[
\begin{array}{c|c|c}
\text{Error budget} & \text{Recovery with high probability} &
                        \text{Failure with high probability} \\
\hline
0\ \text{or fixed} & A>C & A\le C \\
\lfloor e^{Bn^{1/3}}\rfloor,\ B>0 & A-2B\ge C & A-2B<C .
\end{array}
\]

At equality in the leading exponent, the behavior depends on the exact polynomial factor multiplying the exponential scale. Theorem~\ref{thm:face-window} resolves this second scale when the expected sample size is written as \(\kappa n^\beta\) times the exponential threshold. The critical polynomial exponent is \(\beta=1/2\) for noiseless measurements and \(\beta=-5/6\) for an exponentially growing error budget: recovery fails with high probability below the separating exponent and succeeds above it. At the critical exponent itself, the sharp multiplicative cutoff is \(\kappa_*=0.0602978769\ldots\) in the noiseless case, and Theorem~\ref{thm:noiseless-exact} shows that recovery still fails at equality. In the noisy case, the sharp multiplicative cutoff is \(\kappa_{\rm noise}=11.333347057\ldots\), and Theorem~\ref{thm:noisy-exact} proves failure at equality there as well. Notice that equality behaves differently in the two rows of the table, failing in the noiseless regime and succeeding in the noisy regime. This reversal is explained by counting the admissible error vectors under a single total budget across the histogram, rather than under separate budgets in each bin.

\subsection{Which failures remain above the cutoff?}
When the leading exponent strictly exceeds the binary box cutoff by a fixed amount, two different geometric regimes govern the threshold transition and the residual failure probability. The expectation switch at the threshold is controlled by faces that fix about \(36.55\%\) of the coordinates. Above the cutoff, however, those intermediate faces contribute negligibly, and the remaining failure probability is asymptotic to \(n(n-1)(4E+1)/\sqrt{\pi m}\), which comes almost entirely from faces fixing only one or two coordinates. Applying a simple signed-coordinate filter eliminates those low-codimension failures and doubles the logarithmic decay rate of the failure probability; see Corollary~\ref{cor:rare-faces}.

\section{General Hamming graphs}
We now turn from binary faces to general Hamming graphs over a fixed alphabet of size \(q\ge2\) and to the three classical query families: all vertices, outer vertices (vertices lying outside the landmark set), and ordinary edges. Write \(D_n^V,D_n^O,D_n^E\) for the minimum number of landmarks required to distinguish all queries in these three families, respectively, with the convention that the minimum over an empty collection of resolving sets is infinite. In the notation of \cite{FKKY}, \(D_n^V\), \(D_n^O\) and \(D_n^E\) are the multiset dimension \(\dim_m\), the outer multiset dimension \(\dim_{om}\) and the edge multiset dimension \(\operatorname{edim}_m\) of \(H(n,q)\). For a random landmark set, write \(Z^V,Z^O,Z^E\) for the corresponding numbers of unordered colliding pairs in each family.

To state the sharp thresholds, set \(L=\log q\) and \(\sigma^2=(q-1)/q^2\), and define
\[
a(A)=\frac23\sqrt{2\sigma^2}\,A^{3/2},
\qquad C_q=\left(\frac{3\log q}{\sqrt{2\sigma^2}}\right)^{2/3}.
\]
By construction, the constant satisfies \(a(C_q)=2L\). The thresholds stated below govern sparse Bernoulli landmark sets sampled at these prescribed scales. Note that multiset resolution is not monotone under adding landmarks: taking the entire vertex set as landmarks, for instance, assigns the exact same distance histogram to every vertex, and likewise assigns the exact same histogram to every ordinary edge.

\begin{theorem}[Deterministic order and joint box resolution]\label{thm:order}
For each \(D_n\in\{D^V_n,D^O_n,D^E_n\}\),
\[
\frac{C_q}{4^{2/3}}
 \le \liminf_{n\to\infty}\frac{\log D_n}{n^{1/3}}
 \le \limsup_{n\to\infty}\frac{\log D_n}{n^{1/3}}
 \le C_q.
\]
For each \(A>C_q\), independent inclusion with expected landmark count \(e^{A n^{1/3}}\) resolves the three families within each type \whp. If \(A\) is sufficiently large in terms of \(q\), the same construction jointly resolves all \((2^q-1)^n\) coordinate boxes \whp. The minimum size for that stronger problem is also \(\exp(\Theta_q(n^{1/3}))\). For binary boxes the sharp random scale was given above.
\end{theorem}

\begin{theorem}[Exponential failure rate]\label{thm:failure}
Under independent inclusion with mean \(e^{A n^{1/3}}\), \(A>0\), let \(F\) denote failure to distinguish one of the three original query types. Then
\[
\lim_{n\to\infty}\frac1n\log\prob(F)=
\begin{cases}
\min\{0,2L-a(A)\},&q\ge3,\\
\min\{0,L-a(A)/2\},&q=2.
\end{cases}
\]
This applies to each type separately and to their union. For \(A\le C_q\), all three types fail simultaneously \whp. For \(A<C_q\), each collision-count logarithm divided by \(n\) converges in probability to \(2L-a(A)\). The exact probability threshold at \(A=C_q\) is refined below.
\end{theorem}

\begin{theorem}[Sharp classical probability window]\label{thm:window}
Fix \(A=C_q\), \(\kappa>0\), and take mean landmark count \(\kappa\sqrt n\,e^{A n^{1/3}}\). Let
\[
R=\sqrt{2\sigma^2 A},\quad \kappa_1=\frac{\sigma}{2\sqrt{2\pi}},
\quad B=R\log(\kappa_1/\kappa).
\]
For each original query type, let \(Z\) be its number of unordered collisions and \(F=\{Z>0\}\). Then
\[
\begin{aligned}
\lim_{n\to\infty}n^{-2/3}\log\ex Z&=
\begin{cases}B,&q\ge3,\\ \max(B,B/2),&q=2,\end{cases}\\
\lim_{n\to\infty}n^{-2/3}\log\prob(F)&=
\begin{cases}\min(0,B),&q\ge3,\\ \min(0,B/2),&q=2.\end{cases}
\end{aligned}
\]
For \(\kappa>\kappa_1\), one random set resolves all three types \whp. For \(\kappa\le\kappa_1\), all three types fail simultaneously \whp. More generally, if the prefactor \(\sqrt n\) is replaced by \(n^\beta\) with fixed \(\beta\ne1/2\), success holds for \(\beta>1/2\) and failure holds for \(\beta<1/2\). At equality the normalized log expectations are zero, but each expectation diverges.
\end{theorem}


\subsection{Relation to prior work}
Metric dimension has a large literature; Tillquist, Frongillo and Lladser \cite{TFL} survey it together with its applications. For hypercubes the labeled problem is closely related to the coin-weighing problem of Erd\H{o}s and R\'enyi \cite{ER}, solved asymptotically by Lindstr\"om \cite{Lindstrom} and by Cantor and Mills \cite{CM}, and for general Hamming graphs to Chv\'atal's analysis of static Mastermind \cite{Chvatal}; see also \cite{Laird,KMSY}. Bollob\'as, Mitsche and Pra{\l}at \cite{BMP} studied the metric dimension of binomial random graphs. On the multiset side, the outer multiset dimension was introduced by Gil-Pons, Ram\'irez-Cruz, Trujillo-Rasua and Yero \cite{GRTY} and studied further in \cite{KKY}. Edge versions go back to the edge metric dimension of Kelenc, Tratnik and Yero \cite{KTY}, and the edge multiset dimension was introduced by Ikhlaq, Ismail, Siddiqui and Nadeem \cite{IISN}.

Farhan, Klav\v{z}ar, Kuziak and Yero \cite{FKKY} pose the problem of determining the multiset dimension parameters of binary hypercubes and record values for small dimensions. Allikvere \cite{Allikvere} classifies when the binary edge multiset dimension is finite and uses distance-level forests to build random resolving sets. In Section 10, Questions 2 and 4 of \cite{Allikvere}, Allikvere asks for the growth rate of the edge multiset dimension of hypercubes and whether sparse landmark sets can improve upon density one half. Theorem~\ref{thm:order} answers both questions: the order is \(\exp(\Theta(n^{1/3}))\), and sparse random landmark sets of that size already resolve with high probability. For binary edges, the same two constants \((3\sqrt2\log2)^{2/3}\) and \((3\sqrt2\log2)^{2/3}/4^{2/3}\) were obtained independently in a public repository posted in October 2026 \cite{GH}. Theorem~\ref{thm:order} extends these bounds to vertices, outer vertices and edges over every fixed alphabet, and Theorem~\ref{thm:window} sharpens the random upper bound to an exact threshold, including the prefactor \(\kappa_1\sqrt n\).

Several adjacent lines of work use related ideas or address different measurement models. Ordinary vertex multiset resolution is equivalent to an ID-coloring, as shown by Hakanen and Yero \cite[Theorem 2.1]{ID}. Eide and Pra{\l}at \cite{EP} analyze sparse landmarks on random graphs via products of clipped binomial atoms; here we apply related probabilistic tools to the different geometry of moderate-deviation distance shells in deterministic Hamming graphs, extending them to control joint collisions and all coordinate boxes simultaneously. In coding theory, the distance histogram of a set appears as a row of the outer distribution matrix \cite{CCM}. By contrast, labeled Hamming resolvability \cite{Laird} and the resolution of sets by landmark-indexed distances \cite{Several} preserve landmark identities, while counting codewords contained in a face \cite{Abdurakhmanov} observes only the distance-zero entry of the histogram; none of these coincides with the anonymous recovery problem studied here.

\subsection{Proof mechanism and organization}
The stretched-exponential scale emerges from a balance between the entropy of the query family and the information carried by the distance histogram. When the expected number of landmarks \(m\) is exponential in \(n^{1/3}\), the shell counts in the histogram are appreciably populated only within a moderate-deviation band of width on the order of \(n^{2/3}\) around the typical distance. The logarithms of the shell occupancies sum across this band to produce a negative log-probability of order \(n\), which matches the exponential cardinality of the query families.

To prove the upper bounds on the failure probability, we examine vertices at prescribed distances from a pair of queries, which contribute signed increments to the difference of their histograms. Organizing these increments through a forest of distance-level pairs or a descending ordering yields enough counts that can be tested independently. For the binary converse, we restrict attention to collisions between pairs of faces that share the same coordinate support but assign opposite bits on it. We show that their typical joint collision probabilities factor asymptotically as products: the neutral shells localize the relevant Fourier integral, while Walsh orthogonality \cite{ODonnell} controls the covariance between pairs. In the noisy model, we introduce a positive Gaussian weight to carry out the relative second-moment comparison.

The rest of the paper is organized as follows. We first develop the coordinate-box geometry and prove the binary face thresholds and rare-failure results; the proofs for the classical vertex, outer-vertex, and edge problems are collected in the appendices.

\subsection{Open questions}
Three questions remain. First, determine the exact deterministic leading constant in Theorem~\ref{thm:order}; our bounds place it between \(C_q/4^{2/3}\) and \(C_q\).

Second, determine the sharp random threshold for jointly resolving all coordinate boxes when \(q\ge3\), analogous to the binary face thresholds established here for \(q=2\).

Third, transfer the binary face thresholds, with and without adversarial errors, to a uniformly chosen landmark set of fixed size, as Section~\ref{sec:fixed-cardinality} does for the classical query families.
\section{Joint resolution of coordinate boxes}
A coordinate box is a Cartesian product $X=\prod_i S_i$ of nonempty subsets of $[q]$, one for each of the $n$ coordinates. The distance from a word to $X$ is the number of coordinates in which the word falls outside the prescribed subset $S_i$.

In this section we prove that the minimum number of landmarks required to resolve all coordinate boxes simultaneously is $\exp(\Theta_q(n^{1/3}))$. To establish the upper bound, we sample landmarks at random: write $s=n^{1/3}$ and independently retain each word with probability $p=K/q^n$, where $K=e^{As}$ for a constant $A$ to be chosen large in terms of $q$. Write $N$ for the total number of retained landmarks and $H_X$ for the resulting distance histogram of a box $X$.

The argument proceeds in two stages. First, we show that the empirical moments of the distance histograms distinguish any two boxes whose coordinate subset sizes occur with different multiplicities. Second, for two distinct boxes with the same multiplicities of subset sizes, we extract a family of disjoint distance-pair cells forming a forest and use their independence to bound the probability of a histogram collision.

\subsection{Different types}\label{sec:different-box-types}
The type of a coordinate box records the multiplicities of the subset sizes among its $n$ coordinates. We show that the first $q-1$ cumulants of the distance from a uniform random word determine the type with a uniform quantitative gap, and that the empirical cumulants computed from the sampled landmarks approximate these population cumulants within less than half that gap for all boxes simultaneously.

Under a uniformly random word $U$, the coordinates are independent, so the distance from $U$ to a box $X$ is a sum of independent Bernoulli random variables with parameters $p_i=1-|S_i|/q$. The $k$th cumulant of a Bernoulli random variable with parameter $p$, obtained by differentiating its cumulant generating function $\log(1+p(e^z-1))$ $k$ times at zero, is an integer polynomial in $p$ with zero constant term and leading coefficient $(-1)^{k-1}(k-1)!$. Because cumulants are additive over independent summands, each of the first $q-1$ cumulants of the distance is a linear combination of the power sums of the $p_i$ up to that order, with nonzero leading coefficient. By additivity and triangular elimination, the first $q-1$ cumulants therefore determine the power sums $\sum_i p_i^k$ for $1\le k\le q-1$.

Since the nonzero values of $p_i$ are distinct fractions corresponding to the subset sizes smaller than $q$, the linear system recovering their multiplicities from these power sums has coefficient matrix equal to a Vandermonde matrix times an invertible diagonal matrix. This system is invertible, and subtracting the sum of the nonzero multiplicities from $n$ then gives the zero multiplicity (the number of unconstrained coordinates). Because each $p_i$ is an integer multiple of $1/q$ and the $k$th cumulant is an integer polynomial in the $p_i$ of degree $k$, any two boxes of different types must differ in some $k$th cumulant by at least $q^{-k}$.

To compare the empirical moments with their population values, let $I_u$ be the indicator that word $u$ is retained as a landmark. For each pair $X,k$, consider the centered sum $\sum_u (I_u-p)(d(u,X)^k-\mathbb E d(U,X)^k)$. This sum has mean zero, variance at most $Kn^{2k}$, and independent centered summands each bounded in magnitude by $n^k$. By Bernstein's inequality, the probability that this sum exceeds $Kn^{-J}$ in absolute value is at most $2\exp(-cKn^{-2J-2k})$.

We now take a union bound over all $(2^q-1)^n$ boxes and all $q-1$ moment orders, and intersect with the high-probability event $N\ge K/2$. Dividing the centered sum by $N$ on this event shows that all empirical raw moment errors are at most $2n^{-J}$. Since the first $q-1$ cumulants are fixed polynomials in the raw moments, and each raw moment of order $k$ has magnitude at most $n^k$, choosing $J$ sufficiently large in terms of $q$ ensures that every empirical cumulant differs from its population value by less than $1/(4q^{q-1})$. Because this error is less than half the minimum separation between distinct types, the distance histograms distinguish all box types with high probability.

\subsection{Equal types: constructing independent cell counts}\label{sec:box-forest}
We now show that with high probability no two distinct boxes of the same type share the same distance histogram. For any fixed pair of distinct boxes of the same type, our strategy is to construct a large collection of disjoint subsets of words (cells) whose pairs of distances to the two boxes form the edges of a directed forest. Because the landmark counts in disjoint cells are independent and the oriented incidence matrix of a forest has full column rank, a histogram collision forces these counts to take at most one specific vector of values, which we bound using the binomial maximal-atom estimate.

\step{Isolating a homogeneous block and a central window.}
Put $\Lambda=As$, so that the expected number of landmarks is $K=e^{\Lambda}$. Two distinct boxes of the same type have the same number $r\ge1$ of constrained positions (coordinates where the prescribed subset is proper). Among the $r$ constrained positions of the first box, the pigeonhole principle yields a block of $m\ge r/(2^q-1)^2$ positions on which the two boxes share a common ordered subset pair $(S,T)$.

Under a uniform word, the distance to the first box on this block has distribution $j\sim\operatorname{Bin}(m,p)$, where $p=1-|S|/q\in[1/q,1-1/q]$. Outside this block, the pair of distances to the two boxes takes at most $(n+1)^2$ possible values; we fix a most probable pair of off-block contributions $(c_1,c_2)$, which restricts us to an event of probability mass at least $1/(n+1)^2$, at cost at most a factor $(n+1)^2$ in mass. When $S=T$, the two boxes agree on the chosen block, so we must ensure their total distances differ; we reserve one discrepant coordinate outside the block before fixing the remaining off-block contributions, and then fix the letter at that coordinate so that $c_1\ne c_2$. This costs at most another factor $q$ in mass, because for two distinct nonempty subsets the difference of their membership indicators cannot be constant over $[q]$.

Within the block, we restrict attention to the consecutive integers $j\in[0,m]$ satisfying $|j-mp|\le1+\beta\min(m,\sqrt{m\Lambda})$, where $\beta$ is a small positive constant depending only on $q$. The number $W$ of such values of $j$ is at least two and at least $c_q\min(m,\sqrt{m\Lambda})$. Each such $j$ has binomial probability mass $\ge\exp(-\Lambda/4-O_q(\log n))$. Indeed, the atom at $j$ under parameter $j/m$ is modal and therefore has probability at least $1/(m+1)$, and changing the binomial parameter from $j/m$ to $p$ multiplies this probability by $\exp(-mD(j/m\Vert p))$, where the relative entropy satisfies $D(x\Vert p)\le (x-p)^2/[p(1-p)]$. Choosing $\beta$ sufficiently small ensures that $mD(j/m\Vert p)$ is at most $\Lambda/4$.

\step{Constructing a forest of distance-pair cells.}
We next associate to each $j$ in the window an oriented edge from the distance to the first box to the distance to the second box, corresponding to a cell of words with that distance pair. When $S=T$, the two boxes have identical distances on the block, so the level pairs $(j+c_1,j+c_2)$ with $c_1\ne c_2$ directly form a forest of fixed-shift paths.

Otherwise, when $S$ and $T$ are distinct, set $p_1=|S^c\setminus T|/|S^c|$, $p_0=|S\setminus T|/|S|$, and $b_k(u)=\min(k,\lfloor(k+1)u\rfloor)$, which gives a mode of a binomial distribution with $k$ trials and success probability $u$. Conditional on the first block distance $j$, each of the $j$ block coordinates outside $S$ falls outside $T$ with probability $p_1$, and each of the $m-j$ block coordinates inside $S$ falls outside $T$ with probability $p_0$. Hence the second block distance is a sum of two independent binomial random variables with parameters $(j,p_1),(m-j,p_0)$. Fixing each of these two binomials to its mode restricts us to a cell at only a polynomial cost in mass (at most a factor $(n+1)^2$), and defines a map sending the first box distance $j+c_1$ to the second box distance $b_j(p_1)+b_{m-j}(p_0)+c_2$.

This map is 1-Lipschitz in $j$, because incrementing the trial count of a binomial distribution by one changes its mode by either zero or one. On any periodic orbit, a 1-Lipschitz map (a nonexpansion) must be an isometry, as one sees by iterating the map around the orbit; on the integer line, the only isometries of a finite set are the identity and a reversal, so every cycle has length at most two. Moreover, the number of loops (fixed points of the map) is $O_q(1)$, because the fixed-point equation is affine in $j$ up to bounded rounding error, with slope $p_1-p_0-1$ bounded away from zero: having $p_1=1,p_0=0$ would imply $S=T$, which is excluded in this case.

We now delete all loops and, for each 2-cycle (doubled edge), retain only one of its two orientations. The resulting graph is a forest: every vertex has outdegree at most one, so any undirected cycle of length at least three would have to be a directed cycle, which does not exist. At least $(W-O_q(1))/2$ edges remain. If $W$ is bounded, then $m$, and hence $r$, is bounded for large $n$; in that case we simply reserve any single discrepant coordinate between the two boxes and fix all other coordinate contributions to a modal pair, after which a suitable choice of letter at the discrepant coordinate yields at least one nonloop cell of mass at least $1/[q(n+1)^2]$.

\step{Collision probability and union bound.}
The cells corresponding to the retained forest edges are disjoint subsets of words, each of uniform probability mass at least $\exp(-\Lambda/4-O_q(\log n))$. Because landmarks are sampled independently with expected total count $K=e^{\Lambda}$, the number of landmarks falling in each retained cell is an independent binomial random variable of mean $\ge\exp(3\Lambda/4-O_q(\log n))$. By the binomial maximal-atom bound, the probability that any single cell count equals a prescribed integer is $\le e^{-c_q \Lambda}$.

Now condition on all landmark selections outside the retained cells. The difference between the two distance histograms equals a fixed vector (determined by the exterior landmarks) plus the oriented incidence matrix of the forest applied to the vector of independent cell counts. Eliminating leaves iteratively shows that the oriented incidence matrix of any forest has full column rank, so equality of the two histograms requires the vector of cell counts to equal at most one particular vector. Multiplying the maximal-atom bounds over all retained edges gives a collision probability of at most
\[
 \exp\{-c_q \Lambda\min(r,\sqrt{r\Lambda})\}.
\]

It remains to sum this bound over all ordered pairs of distinct boxes of the same type with $r$ constrained positions. There are at most $(e(2^q-2)n/r)^{2r}$ ordered pairs at level $r$. For $r\le \Lambda$, the exponent $c_q \Lambda r$ dominates the logarithm of this pair count. For $r>\Lambda$, the exponent is $c_q A^{3/2}\sqrt{nr}$, and the ratio of the log pair count to $\sqrt{nr}$ is controlled by $\sqrt{r/n}\log(e(2^q-2)n/r)$, which is uniformly bounded. Taking $A$ sufficiently large in terms of $q$ makes the union bound over all levels $r\ge1$ equal to $o(1)$. Level zero consists of the single unconstrained box, so it produces no pairs. Finally, $N\le2K$ holds with high probability, giving the upper bound on the number of landmarks. Singleton boxes give the matching lower bound order by the vertex lower bound in Appendix~\ref{app:entropy}.

\section{Robust histograms and lower bounds from types}
We now extend the analysis to noisy observations in which an adversary may corrupt up to $E$ entries of the distance histogram. Put $E=\lfloor e^{Bs}\rfloor$ with $B\ge0$ (recalling $s=n^{1/3}$), and separately allow the noiseless case $E=B=0$.

For two histograms of equal total mass, the replacement distance is $d(H,H')=\|H-H'\|_1/2$, since each single replacement moves one unit from a surplus bin to a deficient bin. Consequently, the radius-$E$ balls around two histograms intersect if and only if $d\le2E$, as one sees by splitting a shortest path of replacements at its midpoint. Moreover, if some readings are deleted rather than replaced, padding the deleted readings back up to the known original mass with arbitrary distances costs at most one replacement per deletion. Thus a recovery guarantee against replacements automatically covers combined deletions and replacements as well.

\subsection{Uniform robust resolution}\label{sec:robust-forest}
In this subsection we show that the random landmark sample separates all coordinate boxes up to $E$ adversarial replacements whenever $D=A-2B$ is sufficiently large in terms of $q$. The proof adapts the forest construction by summing the histogram difference across forest cuts and verifies that the moment gap between different types survives the perturbation.

Write $D=A-2B$ and take $\Lambda=Ds$. Running the forest construction of Section~\ref{sec:box-forest} with this parameter $\Lambda$, the actual expected landmark counts in the retained cells are now at least $\exp(As-\Lambda/4-O_q(\log n))$, because the landmark sampling rate still has exponent $As$.

Consider any edge of the forest. Deleting that edge splits its connected tree into two parts. Summing the entries of the histogram difference over all vertices on one side of the cut cancels every other forest edge (each edge wholly inside that side contributes $+1$ and $-1$, while edges on the other side contribute $0$), leaving only the signed landmark count on the deleted edge plus a fixed term determined by the exterior landmarks. If $d\le2E$, then the $\ell_1$ norm of the histogram difference is at most $4E$, and because the entries of the histogram difference sum to zero, the absolute value of the sum over any subset of vertices is at most $2E$.

Conditional on the exterior landmarks, each independent cell count on a forest edge is therefore confined to an interval of at most $4E+1$ integers. By the binomial maximal-atom bound, the probability of landing in this interval is at most $\exp(-3\Lambda/8+O_q(\log n))$. Multiplying these probabilities over all edges of the forest and applying the same pair enumeration as in the noiseless case proves uniform separation within each type for $D>D_0(q)$ sufficiently large.

For boxes of different types, a perturbation with $d\le2E$ changes the empirical raw moment of order $j$ by at most $2En^j/N$. On the high-probability event $N\ge K/2$, this perturbation is exponentially small in $s$, so the fixed population cumulant gap still precludes confusion between different types. In particular, one may take $B=(1/2-\varepsilon)A$ for any $0<\varepsilon<1/2$ and sufficiently large $A$.

\subsection{Type entropy and typical histograms}
We next establish lower bounds on the sampling exponent $A$ for both deterministic and random landmark sets. The idea is to fix a box type with large combinatorial entropy and bound how many distinguishable quantized histograms it can generate.

\step{Distance distribution for a uniform box of fixed type.}
For a type fraction vector $\alpha=(\alpha_1,\ldots,\alpha_q)$, specifying the limiting proportion of coordinates whose prescribed subset has each size $a$, set
\[
 p_a=1-a/q,\quad h=-\sum_a\alpha_a\log\alpha_a+
 \sum_a\alpha_a\log\binom qa,
\]\[
 v=\sum_a\alpha_a p_a(1-p_a),
 \qquad C_{\rm box}(q)=\sup_{v>0}\left(\frac{3h}{4\sqrt{2v}}\right)^{2/3}.
\]
Round the type counts to integers summing to $n$; by Stirling's approximation, there are $\exp(hn+O_q(\log n))$ boxes of this type.

For a uniformly random box of this type, the distance from any fixed word has the same distribution $W$: conditional on which coordinates are assigned each subset size, the random subsets at those coordinates are chosen independently, and the resulting Bernoulli sum law does not depend on the coordinate assignment. Write $\pi_j=\Pr(W=j)$ and $\mu=\mathbb EW$. The variance of $W$ is $nv+O(1)$. Inside the central window $|j-\mu|\le Mn^{2/3}$, applying the Chernoff bound with exponential tilt parameter $t=(j-\mu)/\mathrm{Var}(W)$ and cubic remainder $O(n|t|^3)=O(1)$ yields $\log\pi_j\le-(j-\mu)^2/(2nv)+O(1)$. Outside this window, Hoeffding's inequality gives $\pi_j\le2e^{-2M^2s}$, where we take $M$ large relative to $A$.

\step{Deterministic lower bound via entropy.}
To account for adversarial replacements, set the quantization width $w=1+\lfloor4E/(n+1)\rfloor$. If two histograms have the same coordinatewise floors $\lfloor H/w\rfloor$, then each of their $n+1$ bins differs by at most $w-1$, which implies $\|H-H'\|_1\le4E$ and hence intersecting radius-$E$ balls.

Consider any deterministic set of $k=\exp((A+o(1))s)$ landmarks, and feed a uniformly random box of the fixed type into it. In bin $j$, the floored count has mean at most $u=k\pi_j/w$, and therefore Shannon entropy at most $g(u)=(1+u)\log(1+u)-u\log u$, by comparison to the geometric distribution, which maximizes entropy among nonnegative integer random variables of a given mean. We bound $g(u)\le\log_+u+2$ inside the central window and $g(u)=O(u(1+|\log u|))$ for small $u$ outside the window. Since $\log w=Bs+o(s)$, subadditivity of entropy over the bins and passing to the Riemann sum in the rescaled variable $x=(j-\mu)/n^{2/3}$ bound the normalized entropy of the floored histogram by
\[
 \int_{\mathbb R}[A-B-x^2/(2v)]_+\,dx+o(1)
 =\tfrac43\sqrt{2v}[A-B]_+^{3/2}+o(1).
\]
For the floored histogram map to be injective on the type, this normalized entropy must be at least $h+o(1)$. Taking the supremum over types yields $A-B\ge C_{\rm box}(q)$.

\step{Random converse via typical histograms.}
For independent random landmark sampling with expected count $K=e^{As}$, let $\lambda_j=K\pi_j$ be the expected count in bin $j$ for a fixed box of the chosen type. Since each bin count is a binomial random variable, its variance is at most $\lambda_j$. Applying Markov's inequality to the total count outside the central window and Chebyshev's inequality to each bin inside the window shows that, for any given box, with probability $1-o(1)$ all outside counts vanish and each inside count differs from its mean by at most $n^2\max(1,\sqrt{\lambda_j})$.

By Markov's inequality over the random landmark sample, at least half of all boxes of the chosen type satisfy this typicality event with high probability. For such typical boxes, the floored histograms have $O(1+n^2\max(1,\sqrt{\lambda_j})/w)$ possible values in each inside bin. The polynomial factors contribute only $O(n^{2/3}\log n)=o(n)$ to the logarithm of the total number of typical floored histograms. Evaluating the remaining sum as a Riemann integral bounds the normalized log count above by
\[
 \int_{\mathbb R}[A/2-B-x^2/(4v)]_+\,dx+o(1)
 =\tfrac23\sqrt{2v}[A-2B]_+^{3/2}+o(1).
\]
Whenever this quantity is strictly below $h$, the pigeonhole principle forces collisions among typical boxes and precludes recovery. Thus $A-2B<2^{2/3}C_{\rm box}(q)$ implies failure with high probability.

\step{The binary case and separation from vertices and edges.}
In the binary case ($q=2$), write $\rho$ for the fraction of constrained coordinates. Then $v=\rho/4$ and the type entropy is $h_{\rm face}(\rho)=-\rho\log\rho-(1-\rho)\log(1-\rho)+\rho\log2$. The sign of the derivative $(h_{\rm face}/\sqrt\rho)'$ matches that of $2\rho h_{\rm face}'-h_{\rm face}$. Differentiating the latter gives $\log(2(1-\rho)/\rho)-2/(1-\rho)$, which decreases strictly from $+\infty$ to $-\infty$ on $(0,1)$. Since $2\rho h_{\rm face}'-h_{\rm face}$ starts at zero at $\rho=0$ and tends to $-\infty$ as $\rho\to1$, the ratio $h_{\rm face}/\sqrt\rho$ attains its unique maximum at $\rho\simeq0.3655138962$.

At this maximizer we obtain $C_{\rm box}(2)\simeq1.3658580$, giving the random converse threshold $2.1681644\ldots$. For every alphabet size $q$, mixing a small fraction $\eta$ of unconstrained positions into the vertex type strictly increases $h/\sqrt v$, because the entropy gains a term of order $\eta\log(1/\eta)$ while the variance changes only linearly in $\eta$. Already at $\rho=1/2$, the binary converse constant is $(9\log2/2)^{2/3}>2.1$, whereas the vertex and edge success constant is $(3\sqrt2\log2)^{2/3}<2.1$. At $A=2.1$, the same random landmark sample therefore distinguishes all vertices and edges jointly (their population mean distances differ by $1/2$ and concentrate uniformly), yet fails to resolve all boxes with high probability, even without adversarial errors.

\section{The leading first-moment cutoff for binary boxes}\label{sec:face-first}
In the binary case, coordinate boxes are faces of the binary cube. We now compute the exact exponential growth rate of the expected number of colliding pairs of faces at each codimension, and deduce the resulting first-moment threshold.

Let $Z_r$ count unordered pairs of distinct faces with $r$ fixed positions whose radius-$E$ balls intersect, and let $Z$ count all such pairs across all codimensions, including cross-codimension pairs. For $r/n\to\rho\in(0,1]$ set $a_\rho(D)=\tfrac23\sqrt{\rho/2}[D]_+^{3/2}$. We show that
\[
 \lim_{n\to\infty}n^{-1}\log\mathbb E Z_r
 =\max\{2h_{\rm face}(\rho)-a_\rho(D),\ h_{\rm face}(\rho)-a_\rho(D)/2\}.
\]
The two terms inside the maximum reflect the competition between generic pairs of faces and pairs with nearly opposite prescriptions on the same support. Consequently, $\mathbb E Z\to0$ above, and grows exponentially below,
\[
 C=\max_{0<\rho\le1}\left(\frac{3h_{\rm face}(\rho)}{\sqrt{\rho/2}}\right)^{2/3}
 =3.4417465206\ldots
\]
as a cutoff in $D=A-2B$. Above $C$, Markov's inequality immediately gives recovery with high probability. This first-moment computation alone gives no converse. Theorems~\ref{thm:face-wHP} and~\ref{thm:noisy-face} prove the converse and determine the equality behavior at their specified sampling scales.

\subsection{Upper rates at positive limiting codimension}\label{sec:face-modal}
To prove the upper bound on $n^{-1}\log\mathbb E Z_r$, we decompose any pair of distinct faces of codimension $r$ according to how their supports overlap. Depending on whether the overlap is balanced, nearly identical, or nearly opposite, we construct a forest of disjoint distance-pair cells and evaluate the resulting Riemann sum.

\step{Balanced overlap.}
Given two distinct faces with $r$ fixed positions each, split their supports into a block of $a$ common coordinates where their prescribed bits agree, a block of $b$ common coordinates where their prescribed bits are opposite, and two exclusive blocks (each constrained in only one face) of size $c$ each, so that $r=a+b+c$.

Condition on the distance $j$ from a random word to the first face; the mismatches to the first face then form a uniformly random $j$-subset of its $r$ constrained coordinates. Fix the mismatch count $x$ on the agreeing block to a hypergeometric mode, fix the mismatch count $y$ on the opposite block to a conditional hypergeometric mode given $x$, and fix the mismatch count on the second face's exclusive block to $\lfloor c/2\rfloor$. Each of these three restrictions retains a modal outcome out of at most $n+1$ possibilities, so together they cost at most a factor $(n+1)^3$ in conditional probability mass.

Recall that the hypergeometric mode for population size $P$, marked count $M$, and $k$ draws differs from the mean $kM/P$ by $O(1)$ (specifically, one may use $\lfloor(k+1)(M+1)/(P+2)\rfloor$). Since $a=r$ would make the two faces identical and $b/(r-a)\le1$, substituting these modal values into the second face's distance gives
\[
 j'=x+b-y+\lfloor c/2\rfloor,\qquad
 j'-r/2=\frac{a-b}{r}(j-r/2)+O(1).
\]
Fix a small parameter $\delta>0$. When $a,b\le(1-\delta)r$, the factor $(a-b)/r$ has absolute value at most $1-\delta$, so this map strictly decreases $|j-r/2|$ outside a bounded central window around $r/2$. The resulting oriented edges $(j,j')$ therefore form a forest: in any undirected cycle, a radially outermost vertex (one maximizing $|j-r/2|$) would need two edges going inward, hence outdegree at least two, which is impossible because each vertex has outdegree at most one.

For bounded rescaled coordinate $z=(j-r/2)/n^{2/3}$, Stirling's approximation gives cell means at least $\exp((A-z^2/(2v_n))s-O(\log n))$, where $v_n=r/(4n)$, uniformly for $r/n$ bounded away from zero. As in the robust forest argument above, each cut of the forest constrains its edge count to an interval of length at most $4E+1$, costing at most the smaller of one and $4E+1$ times the maximal atom (or a single atom when $E=0$). Summing the log-probabilities over the central window gives a Riemann sum of $[D/2-z^2/(4v_n)]_+$, yielding a collision probability $\le\exp(-a_\rho(D)n+o(n))$. For $D\le0$, we simply use the trivial probability bound of one.

\step{Nearly identical and nearly opposite pairs.}
When $a>(1-\delta)r$, we instead treat the agreeing block as our main block and fix the off-block contributions at their classwise binomial modes, reserving one discrepant coordinate outside the agreeing block. Toggling the bit at that discrepant coordinate changes the difference of the two distances by one (if the coordinate is exclusive) or by two (if it is opposite), so one choice of bit yields unequal offsets. Varying the mismatch count on the agreeing block then gives a forest of fixed-shift paths, and the Riemann sum yields the full integral with $a$ in place of $r$.

When $b>(1-\delta)r$, we treat the opposite block as our main block and fix the agreeing count and both exclusive counts at their lower central modes. The two off-block offsets $u$ then coincide, and pairing mismatch counts $k$ and $b-k$ on the opposite block yields edges $(k+u,b-k+u)$ for $k<b/2$. These edges form a matching (hence a forest) of half the window size, giving half the integral with $b$ in place of $r$. In both extreme regimes, all fixed off-block contributions cost only polynomial factors in $n$.

To sum over all pairs, count the near-opposite choices first. Given the first face and the number $t=r-b\le\delta r$ of coordinates where the second face does not prescribe the opposite bit, there are at most $\binom rt\binom tc\binom{n-r}c2^c$ choices for the second face: mark the $t$ exceptions to opposite bits among the $r$ constrained coordinates, choose $c$ of them to omit from the second support, and prescribe $c$ replacement positions and bits outside the first support. Summing over all $t\le\delta r$ and $c\le t$ costs at most $\exp(O(\delta\log(1/\delta))n+o(n))$ per first face. All other pairs (balanced or near-identical) number at most the square of the total face count. Combining the pair counts with the collision probabilities and letting $\delta\downarrow0$ yields the asserted maximum as an upper bound.

\subsection{Matching lower rates}
We now establish both matching lower bounds on $n^{-1}\log\mathbb E Z_r$. For general pairs of faces of codimension $r$, we construct a cube isometry that expresses the histogram difference as the difference of two independent identically distributed random vectors and apply Cauchy--Schwarz; for exactly opposite pairs on the same support, we exploit reflection symmetry and independence across distance shells.

\step{General pairs.}
Any two faces with the same number $r$ of fixed positions are exchanged by an involutive isometry of the binary cube: flip the coordinates where both faces have conflicting prescriptions, pair up the two exclusive blocks and swap the coordinates in each pair, and flip both bits in a swapped pair whenever the two prescribed bits differ. Under this involution, every word either is neutral (having equal distance to both faces) or belongs to a nonneutral two-word orbit.

On each nonneutral two-word orbit, choose one representative with histogram contribution $e_j-e_{j'}$; its partner then contributes the negative. Because landmarks are sampled independently, the histogram difference has the same law as $X-X'$, where $X$ and $X'$ are independent and identically distributed sums of signed vectors. Each coordinate $j$ of $X$ has variance at most $\lambda_j=K\Pr(\operatorname{Bin}(r,1/2)=j)$, because each orbit contributing to coordinate $j$ contains a distinct word at distance $j$ from the first face.

Summing the shell means outside the central window shows that the probability of any nonzero outside-window contribution is $o(1)$. Inside the window, Chebyshev's inequality simultaneously confines every coordinate of $X$ to its mean plus or minus $n^2\max(1,\sqrt{\lambda_j})$ with probability $1-o(1)$. Quantize each coordinate by $w=1+\lfloor2E/(r+1)\rfloor$; if $X$ and $X'$ have equal coordinatewise floors, then $\|X-X'\|_1\le2E$, which is sufficient for the radius-$E$ balls around the two histograms to intersect.

By the preceding Riemann sum calculation, the logarithm of the number of typical quantized vectors is at most $a_\rho(D)n+o(n)$. By Cauchy--Schwarz, the probability that $X$ and $X'$ land in the same typical quantization bin is at least the squared probability of the typical event divided by the number of typical quantizations. Because this lower bound is uniform over all pairs of faces of codimension $r$, multiplying by the total number $\binom{\binom nr2^r}{2}$ of unordered pairs gives the lower rate $2h_{\rm face}-a_\rho$.

\step{Opposite pairs.}
Now consider a pair of faces with exactly opposite prescriptions on the same support of size $r$. Their distance histograms are exact reflections of each other: a word at distance $j$ from the first face is at distance $r-j$ from the second. Because the distance shells are disjoint, shell counts at distinct indices are independent, and reflected pairs of shell counts are identically distributed.

To force a collision, require both tails outside the central window to be empty (which occurs with probability $1-o(1)$), and for each reflected pair of inside indices require their quantized counts (with width $w$) to be equal. By Chebyshev's inequality, each shell count lies within $2\max(1,\sqrt{\lambda_j})$ of its mean with probability at least $3/4$, so by Cauchy--Schwarz the probability that two independent identically distributed counts fall in the same quantization bin is at least $(3/4)^2$ divided by the number of quantization bins intersecting that interval. Multiplying over half the inside indices costs $\exp(-a_\rho(D)n/2-o(n))$, while the middle shell (when $r$ is even) is invariant under reflection and therefore neutral. Since there are exactly $\binom nr2^r/2$ opposite pairs, this yields the second lower rate.

\subsection{Summing codimensions}\label{sec:sum-codim}
We now sum the collision expectations over all codimensions and include cross-codimension pairs to establish the cutoff at $C$ for $\mathbb E Z$.

First suppose $D>C$, and fix a small constant $\varepsilon>0$. For large codimensions $r\ge\varepsilon n$, the upper-rate argument gives uniform exponential decay: the limiting rate has a strict negative margin on the compact interval $[\varepsilon,1]$, and the loss from passing to a shorter block vanishes uniformly as $\delta\to0$.

For small codimensions $r<\varepsilon n$, we apply the robust forest estimate $\exp(-c\Lambda\min(r,\sqrt{r\Lambda}))$ with $\Lambda=Ds$. There are at most $(2en/r)^{2r}$ pairs at codimension $r$. When $r\le \Lambda$, the exponent $c\Lambda r$ beats the logarithm of $(2en/r)^{2r}$. When $\Lambda<r<\varepsilon n$, the exponent $cD^{3/2}\sqrt{nr}$ dominates the log pair count for small enough $\varepsilon$, since $\sup_{0<x<\varepsilon}\sqrt x\log(2e/x)\to0$ as $\varepsilon\to0$. Both sums over $r$ therefore vanish as $n\to\infty$.

For pairs of faces with different codimensions, the population mean distances differ by at least $1/2$. The uniform Bernstein bound in Section~\ref{sec:different-box-types} fails with probability at most $\exp(-K/\mathrm{poly}(n))$, and on its complementary good event, adversarial replacements shift the empirical mean by at most $2En/N=o(1)$, which cannot close the gap of $1/2$. Even after multiplying by the total number of face pairs, this exceptional probability vanishes, so $\mathbb E Z\to0$.

Finally, when $D<C$, we simply choose a fixed limiting fraction $\rho$ at which the lower rate is strictly positive; the corresponding codimension alone makes $\mathbb E Z$ grow exponentially. Notice that the maximizing fraction in the definition of $C$ is the same unique maximizer of $h_{\rm face}/\sqrt\rho$ determined in the binary type-entropy analysis above.

\section{The sharp noiseless threshold for binary faces}\label{sec:noiseless-face-proof}
\begin{theorem}\label{thm:face-wHP}
In the noiseless binary model of Theorem~\ref{thm:box}, the probability
that all faces are resolved tends to one for fixed
\(A>C=3.4417465206\ldots\), and tends to zero for fixed \(0<A\le C\).
\end{theorem}

\begin{proof}
The recovery statement for \(A>C\) was already established in Theorem~\ref{thm:box}, so only the converse for \(0<A\le C\) remains. We prove it by applying the second moment method to the number of colliding pairs of opposite faces of a fixed codimension, showing that the expected number of collisions diverges while the variance is negligible compared to the squared mean.

Choose \(r\) to be the nearest even integer to \(\rho n\), where \(\rho\) is the maximizing fraction from Theorem~\ref{thm:box}. Each candidate pair consists of two opposite assignments on the same set of \(r\) prescribed coordinates, giving \(N=\binom nr2^{r-1}\) candidates in total. Orient each pair arbitrarily as a first and second assignment. For a cube word at distance \(d\) from the first assignment, its distance to the opposite assignment is \(r-d\); we define its signed jump in row \(j<r/2\) by \(g_j=\mathbf1_{d=j}-\mathbf1_{d=r-j}\). The two faces produce the same distance histogram, and hence collide, if and only if the sum of signed jumps over all sampled landmarks is zero in every row \(j<r/2\). Throughout the proof, write \(s=n^{1/3}\), \(W=n^{2/3}\), \(K=e^{As}\), \(v=\rho/4\), \(R=\sqrt{2vA}\), \(a=\tfrac23\sqrt{2v}A^{3/2}\), and \(\lambda_j=K\prob(\Bin(r,1/2)=j)\) for the expected landmark count in shell \(j\).

\step{Single candidate.}
We first estimate the common collision probability \(P_1\) of a single candidate and verify that the expected number of collisions \(NP_1\) tends to infinity even at \(A=C\). Because distinct shells are disjoint sets of cube words, \(P_1\) factors as a product over the opposing shell pairs \((j,r-j)\) for \(j<r/2\). Each opposing pair receives independent binomial landmark counts with mean \(\lambda_j\), and the probability that their difference is zero has negative logarithm \(\tfrac12(\log\lambda_j)_++O(1)\). Indeed, when \(\lambda_j\ge1\), the upper bound follows from the maximum atom bound for a binomial distribution and the lower bound follows from Cauchy--Schwarz on a Chebyshev interval around the mean; when \(\lambda_j<1\), requiring both shell counts to be zero yields the same estimate.

Restrict attention to the fixed window \(|j-r/2|\le MW\), choosing \(M\) large enough that the combined intensity of all shells outside this window is \(o(1)\). Applying the zero-count estimate to the shells outside the window contributes only \(o(1)\) to \(\log P_1\). Inside the window, Stirling's formula gives
\[
 \log\lambda_j=s\left(A-\frac{((j-r/2)/W)^2}{2v}\right)
                 -\tfrac12\log n+O(1).
\]
Rounding \(\rho n\) to the nearest even integer \(r\) changes the quadratic term by \(o(1)\) uniformly across the window. Summing the positive part of \(\log\lambda_j\) over the half band \(j<r/2\) yields
\[
 \log P_1=-an/2+(R/4)W\log n+O(W).
\]
To justify the error term, observe that the \(-\tfrac12\log n\) contribution shifts the boundary where \(\log\lambda_j\) becomes positive by \(O(s\log n)\) rows, which alters the sum approximation by at most \(O(s\log^2n)\). At the threshold \(A=C\) we have \(a=2h_{\rm face}(\rho)\), so the leading linear exponents in \(N\) and \(P_1\) cancel and the positive subleading term ensures that \(NP_1\) exceeds \(e^{c W\log n}\) for some \(c>0\). Below \(C\), the linear exponent is strictly positive, so \(NP_1\) grows exponentially in \(n\).

\step{Typical joint atoms.}
We next show that for almost all pairs of candidates, the joint collision probability is \((1+o(1))P_1^2\). For two candidates, let \(a_0,b_0\) count respectively the shared prescribed coordinates on which their first assignments agree and conflict, and let \(c\) count the prescribed coordinates exclusive to each candidate. Call the pair typical when \(c\ge\eta n\) and \(|a_0-b_0|/r\le n^{-1/12}\), for a sufficiently small fixed \(\eta>0\). Two independent uniform candidates fail to be typical with probability at most \(e^{-\Omega(n^{5/6})}\): the overlap size follows a hypergeometric distribution, which controls the event \(c\ge\eta n\), and conditional on the overlap, the signed agreement count obeys a Hoeffding bound when the oriented prescriptions are chosen independently and uniformly. Because neither condition depends on the orientations of the two pairs, the same probability bound holds under any arbitrary orientation.

To evaluate the joint collision probability of a typical pair, we pass to a Fourier integral, replace the Bernoulli weights by Poisson weights, and truncate to the retained window. Pairing each cube word with its bitwise complement negates both jump vectors simultaneously, which makes the exact joint characteristic function nonnegative: each complementary pair contributes a Bernoulli factor of the form \(1-2p(1-p)(1-\cos(\theta\cdot g))\) with \(p=K2^{-n}\). Uniformly over all frequencies, replacing these Bernoulli factors by the corresponding Poisson factors changes the integrand by a multiplicative factor \(\exp(O(2^np^2))=1+o(1)\), and the same comparison applies to the two marginal integrals. Next, projecting the jump vectors onto the retained window of rows changes the Poisson integrand by a uniform factor \(1+o(1)\), because the combined intensity of all discarded rows is \(o(1)\) and each summand in the log characteristic function has absolute value at most twice its intensity.

\step{Joint intensities and core localization.}
We now bound the cross-row intensities of a typical pair and localize the Fourier integrals to a small frequency core around the origin. For a typical pair, let \(I_{jk}\) denote the intensity of cube words that produce jumps in absolute folded rows \(j,k\) of the two blocks simultaneously. Uniformly over all retained rows,
\[
 I_{jk}\le \frac{\lambda_j\lambda_k}{K}\exp(o(s)).
\]
To prove this, observe that for a uniform word, the covariance matrix of its distances to the two first assignments has diagonal entries \(r/4\) and off-diagonal entry \((a_0-b_0)/4\). For each of the four sign choices corresponding to the folded rows, tilt each distance exponentially by its centered deviation divided by \(r/4\). Expanding the joint log moment generating function by Taylor's theorem over its independent coordinate blocks produces the sum of the two separate quadratic rate exponents plus a cross-covariance error bounded by \(O(n^{-1/12}s+1)\). The third-order remainder is \(O(n\,n^{-1})\), and comparing the resulting rate against the marginal probabilities via Stirling's formula adds \(O(\log n)\).

Call a retained row large if \(\lambda_j\ge e^{n^{1/4}}\). We claim that for each large row, there is a singleton jump affecting only that row in the first block (and no row in the second block) with intensity at least \(\lambda_j\exp[-O(n^{1/6}+\log n)]\). Indeed, conditional on the word's distance to the first assignment, freeze the three mismatch allocations on the shared agreeing, shared conflicting, and first-exclusive coordinates at their sequential hypergeometric modes; each step retains a polynomial fraction of the probability. Since the conditional slope of the second distance with respect to the first is \((a_0-b_0)/r\), the resulting conditional mean of the second query's distance differs from \(r/2\) by \(O(n^{7/12})\). The independent \(\Bin(c,1/2)\) mismatches on the second candidate's exclusive block then hit the neutral value \(r/2\) of the second query with the stated probability cost. This target value is an integer because \(r\) is even, and the required deviation on the exclusive block is \(o(c)\). Swapping the two blocks gives the analogous singleton construction for the second block.

These singleton jumps localize the Fourier integrals. Define the frequency core by the condition \(|\theta_j|\le e^{n^{1/5}}/\sqrt{\lambda_j}\) for every large row in both blocks. Outside the core, the joint integral is negligible relative to \(P_1^2\): if any large row violates the core bound, its singleton term alone bounds the integrand by \(\exp[-\exp(\Omega(n^{1/5}))]\), using the inequality \(1-\cos\theta\ge2\theta^2/\pi^2\) on the fundamental interval. The exact same localization holds for the two separate marginal integrals.

\step{Quadratic approximation and Gaussian comparison.}
Inside the core, we decouple the small rows from the large rows and replace the large-row exponents by their quadratic parts. The difference between the joint exponent and the sum of the two marginal exponents arising from a two-row jump that involves at least one small row is bounded by a constant times \(I_{jk}|\theta_j\theta_k|\). Summing over all such terms yields \(o(1)\) uniformly on the core: within each block, the sum of \(\lambda_j|\theta_j|\) over small rows is \(\le n e^{o(s)}\), whereas the sum over large rows is \(\le n\sqrt K e^{o(s)}\); applying the intensity bound for \(I_{jk}\) and dividing by \(K\) makes the total cross-contribution vanish. For the remaining large-row terms in both the marginal and joint exponents, we may replace the cosine terms by their quadratic approximations with a uniform \(o(1)\) error, because the sum of all quartic errors is at most \(O(n e^{4n^{1/5}-n^{1/4}})\). Here we bound the fourth power of a sum of two phases by eight times the sum of their fourth powers, and use that the total intensity touching row \(j\) is at most \(2\lambda_j\). As a consequence, the small-row integrals factor out identically in both the joint calculation and the product of marginals.

It remains to compare the large-row Gaussian integrals. For a given candidate, define the normalized radial Walsh vector of degree \(d\) as the sum of the degree-\(d\) coordinate characters on its prescribed support divided by \(\sqrt{\binom rd}\), weighted by the signs of the assignment. Vectors of distinct degrees are orthogonal. For two supports, the cross-Gram entry in degree \(d\) equals the signed sum over all \(d\)-subsets of the shared support, divided by \(\binom rd\). Its absolute value is at most \((1-\eta)^d\), and for each fixed \(d\) it equals \(((a_0-b_0)/r)^d+O_d(1/r)\) by comparing sampling with and without replacement. Because folded rows are antisymmetric under complementation, only odd degrees appear, so the entire cross-Gram matrix has Frobenius norm \(o(1)\) uniformly over typical pairs. Restricting to the large rows is a compression between orthonormal subspaces, which preserves this norm bound. Whitening each large-row coordinate by \(\sqrt{2\lambda_j}\) produces identity within-block covariances and a cross block \(T\) of Frobenius norm \(o(1)\), so the Gaussian determinant factor is \(\det(I-T^{\mathsf T}T)=1+o(1)\). In these whitened coordinates, all marginal variances of the Gaussian Fourier density are bounded, so the tails omitted outside the core are negligible. Combining these comparisons shows that every typical pair has joint collision probability \((1+o(1))P_1^2\) uniformly.

\step{Atypical contributions.}
Finally, we bound the contribution of atypical pairs by splitting them into three regimes according to how closely the two candidates align. First, suppose \(|a_0-b_0|/r\le1-\delta\) for a fixed small \(\delta\). In each block, choose one modal cell per folded row outside \(O_\delta(1)\) central rows; by the conditional mean contraction and the boundedness of rounding errors, any word in the modal cell of a row in one block has a strictly smaller folded radius in the other block. Consequently, the two collections of modal cells are disjoint, and ordering them by decreasing folded radius gives distinct leading coordinates, so the corresponding columns are independent. By the sequential-mode argument, each modal cell captures at least a polynomial fraction of its shell mass. Applying the conditional binomial atom bounds therefore yields a joint collision probability of at most \(P_1^2\exp(O_\delta(W\log n))\), which is overwhelmed by the atypical pair fraction \(\exp(-\Omega(n^{5/6}))\).

Second, suppose \(\sqrt n<r-|a_0-b_0|<\delta r\). The same modal-cell construction omits only \(O(\sqrt n)\) central rows, yielding the joint probability bound \(P_1^2e^{o(n)}\). Counting such pairs by the coordinate changes relative to a given assignment or its opposite shows that there are at most \(N e^{O(\delta\log(1/\delta))n}\) such configurations. For sufficiently small fixed \(\delta\), their total contribution is \(o(N^2P_1^2)\). Third, if \(r-|a_0-b_0|\le\sqrt n\), there are at most \(N e^{O(\sqrt n\log n)}\) such configurations, and the joint probability is trivially bounded by \(P_1\). Dividing their aggregate contribution by \(N^2P_1^2\) yields a quantity that vanishes even at \(A=C\), thanks to the positive subleading correction in \(NP_1\).

Combining the typical and atypical estimates shows that the variance of the number of opposite-face collisions is \(o((NP_1)^2)\). Chebyshev's inequality therefore establishes the converse.
\end{proof}

\section{The sharp threshold with an exponential error budget}\label{sec:noisy-face-proof}
\begin{theorem}\label{thm:noisy-face}
In the binary model of Theorem~\ref{thm:box}, fix \(A>0\) and \(B>0\).
Let \(K=e^{An^{1/3}}\) and \(E=\lfloor e^{Bn^{1/3}}\rfloor\).
The probability of recovering every face under the error budget \(E\)
tends to one for \(A-2B\ge C\) and to zero for \(A-2B<C\), where
\(C=3.4417465206\ldots\) is the binary first-moment constant.
For noiseless recovery or any fixed nonnegative error budget, failure
instead occurs at equality \(A=C\).
\end{theorem}

\begin{proof}
The recovery assertion strictly above the cutoff was proved in Theorem~\ref{thm:box}, and the failure assertion at equality \(A=C\) for any fixed nonnegative error budget follows immediately from Theorem~\ref{thm:face-wHP}. Two claims remain: failure below the cutoff \(A-2B<C\), which we establish by a smoothed second moment method on opposite-face pairs, and recovery at equality \(A-2B=C\) when \(B>0\), which we prove by refining the union bound with a triangular pivot-row exposure and an \(\ell^1\)-ball count. Write \(D=A-2B\), \(s=n^{1/3}\), and \(W=n^{2/3}\).

\step{Failure below the cutoff.}
It suffices to treat the range \(0<D<C\): if \(D\le0\), we may choose a smaller positive \(B'\) at fixed \(A\) such that \(0<A-2B'<C\) and invoke monotonicity in the error budget. Use the same even codimension \(r\) nearest \(\rho n\) at the maximizing fraction \(\rho\), and the same \(N=\binom nr2^{r-1}\) opposite-face candidates, as in the noiseless proof. For a given candidate, let \(S_j=H_j-H_{r-j}\) denote its folded shell differences for \(j<r/2\), which have variance \(2(1-p)\lambda_j\), where \(\lambda_j=K\prob(\Bin(r,1/2)=j)\). Half of the candidate's \(\ell^1\) histogram discrepancy equals \(\sum_j|S_j|\), so the two opposite faces are confusable whenever this sum is at most \(2E\). Set \(M_0=E/n^2\) and assign the candidate the smoothed weight
\[
 Y=\mathbf1_{\sum_j|S_j|\le2E}
       \exp\left(-\frac{\sum_j S_j^2}{2M_0^2}\right).
\]
This weight takes values in \([0,1]\) and vanishes unless the candidate is confusable. Dropping the indicator gives an untruncated weight that differs pointwise from \(Y\) by at most \(e^{-\Omega(n^3)}\), by applying Cauchy--Schwarz on the event \(\sum_j|S_j|>2E\). Consequently, the pointwise difference for a product of two weights is at most twice as large.

We now evaluate the first and second moments of the untruncated weight using Fourier analysis. The untruncated single and joint expectations are Fourier integrals against products of the wrapped centered Gaussian probability measure of frequency variance \(M_0^{-2}\), whose integer Fourier coefficients are \(e^{-k^2/(2M_0^2)}\). By complement symmetry, the characteristic functions of the signed jumps are nonnegative, and the uniform Bernoulli-to-Poisson approximation and window projection estimates from the noiseless proof still apply with multiplicative error \(1+o(1)\). Define the row core half-width by \(T_j=e^{n^{1/5}}/\sqrt{M_0^2+2\lambda_j}\). Whenever \(\lambda_j\le M_0^2\), the wrapped Gaussian measure assigns mass at most \(e^{-\exp(\Omega(n^{1/5}))}\) to the tail \(|\theta_j|>T_j\). Whenever \(\lambda_j>M_0^2\), the typical pair's singleton intensity bound from the noiseless proof gives the same upper bound directly for the characteristic function outside that row core. For the single-candidate integral, either the wrapped Gaussian kernel tail bound or the direct marginal singleton bound applies. Since all remaining coordinate measures have total mass one, every tail outside the core is negligible on the \(e^{-O(n)}\) scale.

Inside the core, the total quartic error in the Poisson characteristic exponent is at most
\[
 O\left( e^{4n^{1/5}}
          \sum_j\frac{\lambda_j}{(M_0^2+2\lambda_j)^2}\right)
 =O\left(\frac{n e^{4n^{1/5}}}{M_0^2}\right)=o(1).
\]
Furthermore, on the core the wrapped Gaussian density equals its main Gaussian term up to a relative factor \(1+o(1)\), uniformly after taking the product over all rows. In the resulting Gaussian integral, the whitened cross block is simply the noiseless cross block attenuated on both sides by the diagonal factors \(\sqrt{2\lambda_j/(M_0^2+2\lambda_j)}\). The signed Walsh estimate from the noiseless proof therefore keeps its Frobenius norm \(o(1)\), so the determinant ratio tends uniformly to one. After normalizing each coordinate by \(\sqrt{M_0^2+2\lambda_j}\), the Gaussian tails outside the core are again negligible. Writing \(p_*=\ex Y\), we conclude that
\[
 \log p_*=-\tfrac12\sum_{j<r/2}\log(1+2\lambda_j/M_0^2)+o(1)
          =-a_\rho(D)n/2+O(W\log n),
 \qquad \ex(YY')=(1+o(1))p_*^2
\]
uniformly over typical pairs. The asymptotic evaluation of the sum follows from the shell expansion of \(\log\lambda_j\) together with \(\log M_0=Bs-2\log n+o(1)\), using that replacing all rows by the retained row window changes the sum by only \(o(1)\).

It remains to bound the atypical pairs. For any atypical pair, the event \(YY'>0\) requires every signed coordinate in both blocks to lie in \([-2E,2E]\). We apply the noiseless two-block modal construction, retaining only pivot rows with \(\lambda_j/E^2>1\). For any prescribed vector of values on these pivot rows, triangularity determines the conditioned cell counts, and each pivot coordinate allows at most \(4E+1\) choices. On pairs with fixed contraction, the combined log loss from the polynomial cell-mass factors, from replacing \(M_0\) by \(E\), and from omitting a bounded number of central rows is \(O(W\log n)\), giving a joint contribution per pair of at most \(p_*^2 e^{O(W\log n)}\). For pairs with contraction slack greater than \(n^{-1/2}\), at most \(O(\sqrt n)\) rows are omitted and the total loss is \(o(n)\), which is absorbed by the entropy gap of near partners. Both regimes use the atypical pair counts from Section~\ref{sec:noiseless-face-proof}. Finally, for micro-close partners we use the trivial bound \(\ex(YY')\le p_*\); there are at most \(Ne^{O(\sqrt n\log n)}\) such pairs, whereas \(Np_*\) grows exponentially in \(n\). Combining all regimes shows that the variance of the sum of weights is \(o((Np_*)^2)\), so by Chebyshev's inequality a confusable pair exists with probability tending to one.

\step{Success at equality.}
Take \(D=C\). Fix \(\varepsilon>0\) small as in the small codimension bound of Section~\ref{sec:sum-codim}, and consider any codimension \(r\ge\varepsilon n\). Write \(\rho=r/n\), \(v=\rho/4\) and \(R_\rho=\sqrt{2vD}\), and focus on the band \(|j-r/2|\le R_\rho W\). Throughout this band, the shell means are exponential in \(s\) because \(B>0\), and uniformly
\[
 \log\lambda_j=As-\frac{(j-r/2)^2}{2vn}-\tfrac12\log n+O(1).
\]

We first establish a general triangular pivot-row upper bound. Suppose we choose \(m\) pivot rows inside this band, and in each pivot row we select a subset of cube words of intensity at least \(\beta_j\lambda_j\) that contributes \(+1\) at that row and places any remaining contribution either outside the pivots or in a pivot row of strictly smaller folded radius. Condition on all landmarks outside these subsets, and expose the subsets in order of decreasing folded radius. For any fixed target output vector on the pivot rows, by the time a pivot row is tested, all contributions from earlier subsets are known and no unexposed group other than its own can affect it. Its own subset contributes an independent binomial count, so the probability of hitting the fixed output vector is at most \(\prod_j O((\beta_j\lambda_j)^{-1/2})\). The subsets of words are pairwise disjoint because each word has a unique leading pivot row. For the pair to be confusable, the projected integer vector on the pivot rows must have \(\ell^1\) norm at most \(4E\), which allows at most \(2^m\binom{4E+m}{m}\) possibilities; since \(E\) is exponential, the logarithm of this count is at most \(m(\log E-\log m+O(1))\). Combining the atom bound with the count of target vectors shows that a full band of \(m=2R_\rho W+O(1)\) pivots has leading log cost \(-a_\rho(D)n\), and a half band of \(m=R_\rho W+O(1)\) pivots has leading log cost \(-a_\rho(D)n/2\), in both cases with the subleading correction
\[
 -\tfrac5{12}m\log n+O(W)-\tfrac12\sum_j\log\beta_j.
\]
Here Stirling's formula in \(\log\lambda_j\) contributes \(+\tfrac14 m\log n\), whereas the \(-m\log m\) term from the total budget count contributes \(-\tfrac23 m\log n\). Whenever the number of omitted rows is \(o(W)\), each omitted row changes the estimate by at most \(O(s)\). Moreover, because \(D=C\), we have \(a_\rho(D)\ge2h_{\rm face}(\rho)\) for all \(\rho\).

\step{Case analysis by alignment.}
We now partition all pairs of faces of codimension \(r\) into three regimes and apply the pivot bound or the forest bound to each.

First, consider a regular pair, meaning \(|a_0-b_0|\le(1-\delta)r\) for a fixed \(\delta>0\). For outer rows at folded radius at least \(\sqrt n\log n\), condition on a word lying in the first shell. Its distance to the second face has variance \(O(n)\): the contribution from shared coordinates is a sum of weights in \(\{-1,0,1\}\) over a uniform subset of specified size, and the contribution from the exclusive coordinates is an independent binomial variable. Because the conditional radius of the mean is at most \((1-\delta)\) times the first radius, Chebyshev's inequality shows that at least half of the shell mass strictly descends in folded radius, and we use all of this descending mass. For inner rows, we instead use the polynomial-mass modal cell, omitting \(O_\delta(1)\) central rows; the extra log penalty from these inner rows is only \(O(\sqrt n\log^2 n)=o(W\log n)\). The number of pairs at codimension \(r\) is at most \(e^{2h_{\rm face}(\rho)n+O(\log n)}\), so the negative subleading correction yields an aggregate \(e^{-\Omega(W\log n)}\) upper bound uniformly over all such codimensions.

Second, consider near agreement, where \(a_0-b_0>(1-\delta)r\). Here the fixed-shift forest bound from the first-moment proof gives
\[
 \log\prob(\text{confusable})\le -\sqrt{a_0/r}\,a_\rho(D)n+o(n).
\]
Counting pairs by choosing edits from one face to the other shows that there are at most \(e^{h_{\rm face}(\rho)n+O(\delta\log(1/\delta))n+O(\log n)}\) such pairs. Since \(a_0>(1-\delta)r\) and \(\rho\ge\varepsilon\), this leaves an order-\(n\) gap in the exponent for sufficiently small fixed \(\delta\).

Third, consider near opposition, and set \(k=r-(b_0-a_0)\). If \(k>s\sqrt{\log n}\), the rounded modal contraction with slack \(k/r\) omits at most \(O(n/k)=o(W)\) central rows. The full-band leading cost \(-a_\rho(D)n+o(n)\) already suffices in this subcase, even with the coarser per-cut interval bound, because counting pairs by edits from the opposite face yields the same entropy gap as in near agreement. If instead \(k\le s\sqrt{\log n}\), then every word satisfies \(|d_1+d_2-r|\le a_0+c\le k\), where \(c\) is the size of each exclusive support. Consequently, every word in a first-distance shell \(j<r/2-k\) in the lower half band contributes \(+1\) in its own pivot row and nothing in any other lower pivot row. We therefore use the full shell mass on the lower half band. Omitting \(O(k+1)\) central rows costs \(O((k+1)s)=o(W\log n)\), and the near-partner count multiplier is at most \(e^{O((k+1)\log n)}=e^{o(W\log n)}\). The half-band leading cost cancels the face count exponent at worst, again leaving \(-\Omega(W\log n)\) from the projection correction.

Summing over all codimensions \(r\ge\varepsilon n\) yields a negligible failure probability. For smaller codimensions, the robust forest bound from the first-moment theorem applies at this fixed positive \(D\), and its constants permit the same sufficiently small \(\varepsilon\). Pairs of faces with different codimensions are separated by their empirical means. This proves the endpoint assertion.
\end{proof}

\section{The binary-face threshold on the polynomial scale}
Let \(\rho_*\) be the maximizing fraction in Theorem~\ref{thm:box}, and define the associated variance parameter, scaled band half-width, and threshold constant by
\[
 v_* = \rho_*/4,\qquad R_* = \sqrt{2v_* C},\qquad
 \kappa_* = \frac{\sqrt{v_*}}{2\sqrt{2\pi}}=0.0602978769\ldots.
\]
The next theorem pins down the polynomial-scale transition between the failure and recovery regimes of Theorems~\ref{thm:face-wHP} and~\ref{thm:noisy-face}.

\begin{theorem}\label{thm:face-window}
Suppose the noiseless landmark mean is \(K=\kappa n^\beta e^{Cn^{1/3}}\). Then the probability of resolving every face tends to zero for \(\beta<1/2\) and to one for \(\beta>1/2\). At the critical exponent \(\beta=1/2\), this probability tends to zero for \(\kappa<\kappa_*\) and to one for \(\kappa>\kappa_*\).

For fixed \(B>0\) and error budget \(E=\lfloor e^{Bn^{1/3}}\rfloor\), suppose \(K=\kappa n^\beta e^{(C+2B)n^{1/3}}\). Then the probability of recovering every face under the error budget tends to zero for \(\beta<-5/6\) and to one for \(\beta>-5/6\). At the critical exponent \(\beta=-5/6\), this probability tends to zero for \(\kappa<\kappa_{\rm noise}\) and to one for \(\kappa>\kappa_{\rm noise}\), where
\[
 \kappa_{\rm noise}=\kappa_*\left(\frac{4e}{R_*}\right)^2
 =11.333347057\ldots .
\]
Theorems~\ref{thm:noiseless-exact} and~\ref{thm:noisy-exact} establish failure at the respective exact multipliers.
\end{theorem}

To prove Theorem~\ref{thm:face-window}, we track the polynomial prefactors and leading constants in both the pivot union bounds and the second-moment lower bounds. For the upper bounds we refine the pivot groups to use both orientations of each pair simultaneously, while for the noisy lower bound we introduce a hybrid weight that combines exact dense-row counts with a Gaussian tilt on the remaining rows.

\begin{proof}
Write \(s=n^{1/3}\) for the exponential scale and \(W=n^{2/3}\) for the bandwidth scale. For any codimension \(r=\rho n\ge\varepsilon n\), set \(v=\rho/4\) and \(R=\sqrt{2vC}\). Within the band \(|j-r/2|\le RW\), the shell-mean expansion takes the form
\[
\begin{split}
 \log\lambda_j={}&
 \left(C+2B-\frac{((j-r/2)/W)^2}{2v}\right)s
   + (\beta-\tfrac12)\log n
   +\log\kappa -\tfrac12\log(2\pi v)+o(1).
\end{split}
\]
We set \(B=0\) in the noiseless case. In this expansion, the cubic correction vanishes by the symmetry of the binomial distribution around \(r/2\), and the quartic remainder is \(o(1)\) uniformly for fixed \(\varepsilon\). For the noiseless calculation, we retain only rows whose leading log intensity is at least \(M\log n\), where \(M\) is a fixed large constant depending on \(\beta\). Compared with the full band, discarding these transition rows changes the total log bound by only \(O(s\log^2 n)=o(W)\).

\step{Upper bounds.}
We begin by tightening the union-bound estimates for regular and near-agreeing pairs so that the constant \(\kappa_0(v)\) emerges cleanly. Refine the pivot groups of Section~\ref{sec:noisy-face-proof} so that each row uses both orientations of the pair: landmarks at distance \(j\) from the first face with strictly smaller distance to the second face, together with landmarks at distance \(j\) from the second face with strictly smaller distance to the first face. Within row \(j\), these two subsets are disjoint and contribute with opposite signs; across different rows, the groups are also disjoint because of the strict inequality on the opposite radius.

For regular pairs at radius at least \(\sqrt n\log n\), conditional Chebyshev's inequality shows that each group retains a \((1-o(1))\) fraction of its shell mass. The binomial Fourier bound for the signed sum of both subsets is
\[
 (1+o(1))(4\pi\lambda_j)^{-1/2}.
\]
Indeed, the modulus of the characteristic function sees only the combined pool of independent Bernoulli variables, whose variance is \((2-o(1))\lambda_j\). Fourier inversion and Gaussian rescaling therefore yield this bound uniformly whenever the shell means diverge. We expose both full sets together in decreasing order of radius. On the inner pivot rows we fall back to the earlier polynomial-mass modal groups, whose combined log loss is \(O(\sqrt n\log^2 n)=o(W)\). Only \(O(1)\) central rows are omitted, and near-agreeing pairs retain the same order-\(n\) entropy gap as before.

\step{Near-opposite pairs.}
Next we handle pairs whose prescribed assignments nearly oppose each other, where the full-band pivot structure must be replaced by a half-band argument. Put \(k=r-(b_0-a_0)\), where \(a_0,b_0\) denote the numbers of agreeing and conflicting coordinates between the two faces. When \(k>Hs\) for a sufficiently large fixed constant \(H\), the full ordered modal pivots omit at most \(O(n/k)\) rows, and the resulting \(O(n/H)\) loss is absorbed into the extensive entropy gap.

When \(k\le Hs\), we use both orientations in the lower half band, stopping below \(r/2-\sqrt{k+1}\log n\). Given that the first distance equals \(j\), the conditional mean of \(d_1+d_2-r\) is \((k/r)(j-r/2)\). Its conditional variance is \(O(k)\), because the conditional sample from the first face's support assigns nonzero weights only to agreeing and exclusive coordinates, while the other exclusive binomial has size at most \(k\). Consequently, the conditional mean of \(d_2\) exceeds \(r/2\) by \((1-k/r)(r/2-j)\). Chebyshev's inequality then shows that \(d_2\ge r/2\) holds with conditional mass \(1-o(1)\), and the reverse orientation behaves symmetrically. Each of these groups touches only its own lower pivot. Omitting the central rows costs \(O(\sqrt{k+1}\log n\cdot s)=o(W)\), and the log count of such neighbor pairs is \(O((k+1)\log n)=o(W)\).

\step{Combining the pivot bounds.}
Let \(m\) be the number of chosen pivot rows, and set \(R'=2R\) for a full band or \(R'=R\) for a half band, so that \(m=R'W+o(W/\log n)\). Without errors, only the zero discrepancy vector is permitted. Applying Stirling's approximation and the signed atom bound across the \(m\) pivots gives the log probability bound
\[
 -\frac{R'}{2R} a_\rho(C)n
 -\frac m2\big[(\beta-\tfrac12)\log n+
                      \log(\kappa/\kappa_0(v))\big]+o(W),
 \qquad \kappa_0(v)=\frac{\sqrt v}{2\sqrt{2\pi}} .
\]

In the noisy case \(B>0\), any confusable pair has total discrepancy at most \(4E\) in \(\ell^1\). On a full band, we apply this budget directly to the projected pivot vector. For near-opposite pairs, at least one of the lower and upper half bands must have discrepancy at most \(2E\), without requiring any assignment of the central rows. We therefore apply either the lower half-band groups constructed above or their reflection to the upper half band, and sum the two bounds; the reflection is valid by bitwise complementation.

Thus, with \(m=R'W+o(W/\log n)\) pivots, the discrepancy budget \(Q\) is \(4E\) on a full band and \(2E\) on a half band. In both cases, counting allowed pivot vectors by \(2^m\binom{Q+m}{m}\) shows that their logarithm is at most \(m\log(4eE/(RW))+o(W)\). Relative to the same leading exponential term, the log correction per pivot becomes
\[
 -\left(\frac\beta2+\frac5{12}\right)\log n
   -\tfrac12\log(\kappa/\kappa_0(v))
   +\log(4e/R)+o(1).
\]
Because \(a_\rho(C)\ge 2h_{\rm face}(\rho)\), a negative log correction immediately yields the strict success regimes on the polynomial scale. At the critical powers of \(n\) and the stated strict values of \(\kappa\), continuity provides a strictly negative correction on a fixed neighborhood of \(\rho_*\), while uniqueness of the maximizer provides an extensive exponential margin away from that neighborhood. As in the earlier proofs, smaller codimensions and cross-codimension pairs remain negligible, since polynomial factors in the landmark mean cannot overcome their fixed exponential slack.

\step{Lower bounds.}
We now establish the matching failure regimes using the second moment method on opposite-assignment face pairs at the optimal codimension. Take the even codimension nearest \(\rho_* n\) and consider the same family of opposite-assignment candidates.

In the noiseless case, the folded rows of a single candidate are independent, and each dense row has zero-atom probability \((1+o(1))(4\pi\lambda_j)^{-1/2}\). The transition rows cost \(o(W)\) as in the upper bound, and the farther tails are negligible because their summed intensities tend to zero. Combining these estimates gives the log mean number of collisions:
\[
 -\frac{R_*}{2}W\big[(\beta-\tfrac12)\log n+
                               \log(\kappa/\kappa_*)\big]+o(W).
\]
This logarithm grows at least as a positive multiple of \(W\) throughout the stated failure regimes. Polynomial prefactors in the mean do not change the typical relative Fourier estimates or the atypical macro and meso bounds, and micro-close pairs cost only \(O(\sqrt n\log n)=o(W)\). The same relative second moment calculation therefore proves failure.

For the noisy lower bound, we replace the hard global \(\ell^1\) indicator by a hybrid weight that treats dense rows exactly and tilts the remaining rows smoothly. Set \(M_0=E/n^2\), \(\gamma=1/9\), and \(\eta_n=n^{-1/2}\), and declare a folded row dense when \(\lambda_j\ge E^2\exp(4n^\gamma)\). Assign each candidate the weight
\[
 \mathbf1_{\sum_{\rm dense}|S_j|\le (2-\eta_n)E}
 \mathbf1_{\sum_{\rm other}|S_j|\le \eta_n E}
 \exp\left(-\sum_{\rm other} S_j^2/(2M_0^2)\right).
\]
Every pair in the support of this weight is confusable. For a single candidate, the folded rows are independent, and the tilted first absolute moment on each nondense row is \(O(M_0)\). When \(\lambda_j\le M_0^2\), this bound follows directly from the untilted variance. When \(\lambda_j>M_0^2\), it follows by combining the local central lower bound \(\gtrsim M_0/\sqrt{\lambda_j}\) for the normalizer with the maximum-atom bound inside the Gaussian moment sum. Applying Markov's inequality to these tilted nondense rows shows that the second indicator discards only a \(O(n M_0/(\eta_n E))=o(1)\) fraction of the mean.

For any prescribed dense vector of \(\ell^1\) norm at most \(2E\), the marginal local atom probability is uniformly \((1+o(1))\prod_{\rm dense}(4\pi\lambda_j)^{-1/2}\); the hybrid Fourier calculation below also yields this single-block estimate. On the nondense rows, the Gaussian product weight has log mean \(-\tfrac12\sum_{\rm other}\log(1+2\lambda_j/M_0^2)+o(1)\), by the same calculation as in the preceding noisy proof restricted to those rows. This sum is \(o(W)\). Indeed, the shoulder between \(\lambda_j/M_0^2=n^{-3}\) and the dense cutoff contains \(O(s(n^\gamma+\log n))\) rows and contributes \(O(s(n^\gamma+\log n)^2)=o(W)\), while below the shoulder the intensity ratios sum to \(o(1)\).

The number of dense rows is \(m=R_*W+O(s(n^\gamma+\log n))\), and discarding the shoulder changes the leading shell sum, after subtracting the extensive \(B\) terms, by \(o(W)\). By counting positive and nonnegative compositions, the number of integer points in an \(\ell^1\) ball of radius \(Q\) lies between \(2^m\binom Qm\) and \(2^m\binom{Q+m}{m}\). Setting \(Q=\lfloor(2-\eta_n)E\rfloor\), its logarithm is \(m\log(4eE/m)+o(W)\). Multiplying by the dense atom probability and cancelling the extensive rate against the log candidate count gives the log mean total weight:
\[
 R_*W\left[
 \log(4e/R_*)-\tfrac12\log(\kappa/\kappa_*)
                -\tfrac12(\beta+\tfrac56)\log n
 \right]+o(W).
\]

\step{Relative hybrid Fourier estimate.}
To control the second moment of the weighted count on typical candidate pairs, we show that the joint hybrid Fourier integral factorizes up to a \(1+o(1)\) factor uniformly over prescribed dense vectors. Use the typicality conditions of Section~\ref{sec:noiseless-face-proof}, but tighten the signed-overlap cutoff to \(n^{-1/8}\); the excluded atypical fraction is still \(\exp(-\Omega(n^{3/4}))\).

For two prescribed dense vectors \(z,z'\) in their respective \(\ell^1\) balls, integrate with Haar measure on the dense frequencies and with the wrapped Gaussian of frequency variance \(M_0^{-2}\) on all other coordinates, weighted by the target phase \(\exp(-i\theta\cdot z-i\theta'\cdot z')\). Drop distant rows of total intensity \(o(1)\) and pass to Poisson characteristic functions as before. Under the tightened overlap cutoff, the conditional neutral singleton bound improves to \(\lambda_j e^{-O(n^{1/12}+\log n)}\). This holds because an off-anchor slope of at most \(n^{-1/8}\) shifts the conditional mean by \(O(Wn^{-1/8})\), and the other exclusive binomial pays the square of that shift divided by order \(n\).

Write \(L_j=0\) on dense rows and \(L_j=M_0\) on nondense rows, and localize both frequency blocks to the box
\[
 |\theta_j|\le e^{n^\gamma}/\sqrt{2\lambda_j+L_j^2} .
\]
On rows with \(\lambda_j\ge L_j^2\), the singleton bounds control the integral outside this box; on the remaining rows, the wrapped Gaussian tails provide the localization. Using the normalized measures and the condition \(\gamma>1/12\), the omitted tail integral has absolute value at most \(\exp(-\exp(\Omega(n^\gamma)))\). This error is negligible compared to the product of the main Gaussian integrals, whose log magnitudes are \(-O(n)\) up to the dense normalizations, which themselves incur a log cost of \(O(n)\).

On the localized core, the dense phase is uniformly \(1+o(1)\), because \(\|z\|_1\le2E\) and every dense frequency is at most \(O(E^{-1}e^{-n^\gamma})\). The uniform relative Poisson and projection errors therefore transfer directly to the core integrals. The quartic error is \(O(ne^{4n^\gamma}/M_0^2)=o(1)\). After diagonal scaling, the covariance matrix is the Walsh cross-Gram matrix with coordinate contractions in both blocks; its cross part has Frobenius norm \(o(1)\), so the Gaussian determinant ratio is \(1+o(1)\). Truncating these local Gaussian integrals back to the box also has negligible relative cost, as before.

Consequently, the hybrid joint integral is at most \(1+o(1)\) times the product of the corresponding single-candidate integrals, uniformly over \(z,z'\). Summing over both dense balls, upper bounding the tail-truncated joint weights by the untruncated ones, and applying the single-candidate truncation lower bound completes the estimate for typical pairs. The atypical estimates continue to apply to the supported weights: the fixed-contraction atypical count loses \(\exp(-\Omega(n^{3/4}))\) against a shell penalty of \(\exp(O(W\log n))\), and the intermediate close-pair count retains its extensive gap. In the microscopic regime, the total mean weight grows faster than \(\exp(O(\sqrt n\log n))\) on the asserted failure sides. Chebyshev's inequality now proves failure for \(\beta<-5/6\), and for \(\beta=-5/6\) with \(\kappa<\kappa_{\rm noise}\).
\end{proof}

\section{The exact noiseless endpoint}\label{sec:noiseless-exact}
\begin{theorem}\label{thm:noiseless-exact}
For noiseless recovery at the critical landmark mean \(K=\kappa_*\sqrt n\,e^{C n^{1/3}}\), the probability of resolving all binary faces tends to zero.
\end{theorem}

At the exact multiplier \(\kappa_*\), both the extensive order-\(n\) terms and the order-\(W\) terms cancel in the first moment of opposite-face collisions. We therefore track the order-\(s\) contribution to show that the mean number of collisions still diverges exponentially in \(s\), and then establish a sharper joint collision bound for distinct close candidate pairs so that the second moment method succeeds.

\begin{proof}
\step{First-moment lower bound.}
Take the even codimension \(r\) nearest \(\rho_*n\) and the same family of \(N=\binom nr2^{r-1}\) opposite-face candidate pairs. Set \(v=r/(4n)\), \(R=\sqrt{2vC}\), \(s=n^{1/3}\), and \(W=n^{2/3}\), and index each row in the lower half band by its offset \(u=r/2-j\in\{1,\ldots,r/2\}\) from the center. Define \(t_u=Cs-u^2/(2vn)\). In any fixed moderate window \(u\le MW\), Stirling's approximation gives
\[
 \lambda_j\le (4\pi)^{-1}e^{t_u+\epsilon_n},
 \qquad \epsilon_n=O(n^{-2/3}).
\]
Indeed, the symmetric relative entropy is at least its quadratic term, the local prefactor changes by \(O(u^2/r^2)=O(n^{-2/3})\) in the logarithm, and rounding \(r\) changes \(\kappa_*/\kappa_0(v)\) by \(1+O(n^{-1})\). Farther out beyond this window, the total intensity is \(o(1)\) once \(M\) is chosen large and fixed.

Consider the negative log zero-atom probability for the difference of two independent Poisson variables of mean \(\lambda\),
\[
 g(\lambda)=-\log\left(\frac{1}{2\pi}
       \int_{-\pi}^{\pi}e^{-2\lambda(1-\cos\theta)}\,d\theta\right).
\]
This function is increasing and nonnegative. As in Section~\ref{sec:noiseless-face-proof}, passing from Bernoulli variables to Poisson variables costs only \(o(1)\) in the logarithm of the full collision probability \(P_1\), so it suffices to bound the sum of \(g(\lambda_j)\) over the half band.

We now bound the excess of \(g(e^t/(4\pi))\) over \(\tfrac12t_+\). For \(t\le0\), requiring both Poisson counts to vanish gives \(g(e^t/(4\pi))\le e^t/(2\pi)\). For \(t>0\), substituting \(1-\cos\theta\le\theta^2/2\) and integrating over the fundamental interval yields
\[
 g(e^t/(4\pi))\le
 \tfrac12 t-\log\operatorname{erf}\bigl(\sqrt{\pi}e^{t/2}/2\bigr).
\]
Using the bound \(1-\operatorname{erf}x\le e^{-x^2}\) for \(x\ge0\), an integrable upper bound for the correction to \(\tfrac12t_+\) is
\[
 b(t)=
 \begin{cases}
 e^t/(2\pi),&t\le0,\\
 \displaystyle
 \frac{e^{-(\pi/4)e^t}}{1-e^{-\pi/4}},&t>0.
 \end{cases}
\]
Applying \(e^t\ge1+t\) on the positive half-line gives the explicit integral bound
\[
 \int_{\mathbb R} b(t)\,dt
 \le \frac1{2\pi}+
   \frac{e^{-\pi/4}}{(\pi/4)(1-e^{-\pi/4})}<1.23 .
\]

Next we evaluate the two sums over \(u\ge1\). The half-band shell sum, which omits the central row, satisfies
\[
 \tfrac12\sum_{u\ge1}(t_u+\epsilon_n)_+
      =\tfrac12 a_{r/n}(C)n-\tfrac14Cs+O(1).
\]
This identity follows either by summing the quadratic directly up to its positive boundary or by the trapezoid estimate, since the last boundary interval costs \(O(1)\); the shift \(\epsilon_n\) likewise contributes only \(O(1)\). Near the boundary \(u=RW\), the spacing between consecutive values of \(t_u\) is asymptotic to \(R/(vs)\), which yields the Riemann-sum asymptotic
\[
 \sum_{u\ge1} b(t_u+\epsilon_n)
       =\frac{vs}{R}\int_{\mathbb R}b(t)\,dt+o(s).
\]
To justify this Riemann sum, restrict first to bounded values of \(t\). For \(u\ge RW/2\), the spacing is bounded below by a constant times \(s^{-1}\), so the two exponential tails of \(b\) contribute an arbitrarily small \(O(s)\) remainder. For \(u<RW/2\), the positive argument is of order \(s\), making that portion of the sum negligible. Both here and in the previous display, we may extend the sum in \(u\) past the moderate window because the artificial far tails decay exponentially.

At the optimizing fraction, \(a_{\rho_*}(C)/2=h_{\rm face}(\rho_*)\), and rounding \(r\) alters this extensive cancellation by only \(O(1)\). Combining Stirling's formula for \(N\) with the preceding bounds gives
\[
 \log(NP_1)\ge
 \left(\tfrac14 C-1.23\frac vR\right)s+o(s)\ge c s
\]
for a positive fixed constant \(c\). To check positivity, note that evaluating the variational expression at \(\rho=1/2\) yields \(C>3\), while \(v\le1/4\) implies \(v/R\le 1/\sqrt{8C}\).

\step{Close-pair bound and second moment.}
We now strengthen the joint collision bound for distinct close candidate pairs, because using only a single candidate's atom would lose too much at this exact endpoint. Orient each pair of opposite-assignment candidates so that their prescribed assignments almost agree whenever \(|a_0-b_0|>(1-\delta)r\); their common agreeing block then has size \(a_0>(1-\delta)r\).

Because the two pairs are distinct, they must disagree on at least one shared or exclusive coordinate. Reserve one such coordinate. On all other off-block positions, namely the conflict block and the two exclusive blocks, fix the mismatch counts at their respective binomial modes. The two possible choices for the reserved bit define two disjoint classes of landmarks, each having polynomial relative mass. In each query, the distance offsets coming from these fixed off-block coordinates lie within \(O(1)\) of \((r-a_0)/2\). Moreover, the resulting relative shifts \(\Delta_1,\Delta_2\) between the two queries differ: flipping the reserved bit changes the offset difference by one if the reserved coordinate is exclusive, and by two if it is a conflicting coordinate.

Within each class, vary the mismatch count on the common agreeing block across the lower moderate band, staying a bounded number of rows below the center so that both folded signs are positive. In the two blocks of discrepancy coordinates, each cell then produces a column \(e_j+e_{j+\Delta_i}\), with the two summands lying in separate blocks. Negate the second block to view these columns as oriented incidence vectors in a bipartite graph. The columns form the edges of two fixed-shift matchings. Any cycle in this graph would have to alternate between the two matchings, but every two-step move applies a fixed nonzero translation to the row index, so no cycle can exist.

Because these disjoint cells carry independent counts and have an acyclic incidence graph, taking the maximal atom product across the two half bands yields the joint collision bound
\[
 \prob(\text{both collide})
       \le \exp\left[-\sqrt{a_0/r}\,a_{r/n}(C)n+o(n)\right]
\]
uniformly over such distinct close neighbors. Here the shell-volume calculation is the usual half-band integral for the Binomial\((a_0,1/2)\) common block applied twice, and the polynomial mass factors and bounded offsets contribute only \(o(n)\).

Finally, we assemble the second moment. The number of close candidate pairs is at most \(N\exp(O(\delta\log(1/\delta))n)\). Since \(P_1=\exp[-a_{r/n}(C)n/2+o(n)]\), their total contribution relative to \((NP_1)^2\) is exponentially small for small fixed \(\delta\). Identical pairs contribute \(NP_1\), which is negligible relative to \((NP_1)^2\) by the positive order-\(s\) lower bound on the log mean. Typical pairs still have relative joint atom \(1+o(1)\), and the remaining atypical fraction is controlled by the fixed-contraction estimate of Section~\ref{sec:noiseless-face-proof}. The relative variance therefore tends to zero, proving that an opposite collision occurs with probability tending to one.
\end{proof}

\section{The exact noisy endpoint}\label{sec:noisy-exact}
\begin{theorem}\label{thm:noisy-exact}
Fix \(B>0\) and \(E=\lfloor e^{Bn^{1/3}}\rfloor\). At
\(K=\kappa_{\rm noise}n^{-5/6}e^{(C+2B)n^{1/3}}\), the probability
of recovering all binary faces under budget \(E\) tends to zero.
\end{theorem}

\begin{proof}
We prove failure at the exact noisy endpoint by applying the second moment
method to a weighted count of confusable opposite-face pairs of codimension
near the optimal value. The main idea is to split the distance-histogram rows
into dense rows, where we enforce an exact \(\ell^1\) constraint, and nondense
rows, where we exponentially tilt the counts and truncate to a narrow window
around the tilted mean so that every positively weighted pair is guaranteed to
be confusable.

Let \(r\) be the even codimension nearest to \(\rho_*n\), and consider its
\(N=\binom nr2^{r-1}\) opposite-face candidates.
Put \(s=n^{1/3}\), \(W=n^{2/3}\), \(v=r/(4n)\),
\(R=\sqrt{2vC}\), \(\tau=RW/(2E)\), and \(\gamma=1/10\).
As before, for each \(j<r/2\), write \(S_j=H_j-H_{r-j}\) for the signed
difference between the landmark counts at distance \(j\) and distance \(r-j\),
and let \(\lambda_j=K\prob(\Bin(r,1/2)=j)\) be their common expected count.
We call a row dense if \(\lambda_j\ge E^2e^{4n^\gamma}\).
There are \(m=RW+O(sn^\gamma)\) dense rows. Let \(T\) denote the
sum of \(|S_j|\) over all nondense rows.

\step{A weight supported at the exact budget.}
To control the nondense rows, we tilt each nondense row by the factor
\(e^{-\tau |S_j|}\) while preserving independence across rows. Under this
tilted law, we have \(\ex T=O(q_0/\tau)\) and
\(\operatorname{Var}T=O(q_0/\tau^2)\), where \(q_0=O(sn^\gamma)\) bounds the
number of nondense rows satisfying \(\lambda_j\tau^2\ge n^{-3}\). The bound on
the size of this shoulder of nondense rows follows from Stirling's
approximation.

To justify these moment estimates, we partition the nondense rows into three
groups. For shoulder rows with \(\lambda_j\le\tau^{-2}\), the weight is a
decreasing function of \(|S_j|\), so monotone weighting bounds the tilted
moments by the unweighted ones, whose second moment is at most \(2\lambda_j\).
For the remaining shoulder rows, the local lower bound on point probabilities
over \(|S_j|\le\tau^{-1}\) shows that the normalizing constant is
\(\gtrsim1/(\tau\sqrt{\lambda_j})\); combining this with the maximum point
probability bounds the weighted first and second absolute moments by
\(O(\tau^{-1})\) and \(O(\tau^{-2})\) after normalization. Finally, the
remaining light rows have total tilted first and second absolute moments
\(o(\tau^{-1})\) and \(o(\tau^{-2})\) by Cauchy--Schwarz and their
\(n^{-3}\) intensity cutoff.

Now let \(T_+,T_-\) be the tilted mean of \(T\) plus and minus \(\log n\)
times the uniform standard-deviation bound \(O(\sqrt{q_0}/\tau)\), clipping
\(T_-\) below at zero. Assign each opposite-face candidate the weight
\(Y=e^{\tau(T_+-T)}\) if \(T\in[T_-,T_+]\) and its dense \(\ell^1\) norm
is at most \(2E-T_+\), and assign it weight zero otherwise. Whenever a
candidate has positive weight, its total \(\ell^1\) norm across all rows is at
most \(2E\), so positive weight implies confusability. By Chebyshev's
inequality, restricting to this window preserves a \(1-o(1)\) fraction of the
tilted tail mean. Moreover, because \(T\) lies within the window below \(T_+\)
on the support of \(Y\), we have the uniform bound
\(\sup Y\le e^{O(\sqrt{q_0}\log n)}=e^{o(s)}\).

\step{The mean at the order-\(s\) scale.}
We next compute \(\log\ex Y\) up to \(o(s)\) and verify that the expected
total weight \(N\ex Y\) grows exponentially on the \(s\) scale. Uniformly over
all target vectors in the dense \(\ell^1\) ball, the point probabilities on the
dense rows equal \((1+o(1))\prod_{\rm dense}(4\pi\lambda_j)^{-1/2}\); the
Fourier estimate proved below also establishes this uniformity. By standard
composition bounds and Stirling's formula, the integer \(\ell^1\) ball of
radius \(Q\) on \(m\) rows has log volume
\(m\log(2eQ/m)+O(\log m)+o(1)\) whenever \(Q\gg m^2\). Taking
\(Q=\lfloor 2E-T_+\rfloor\), the log volume loses
\(mT_+/(2E)+O(q_0^2/W)\) relative to a ball of radius \(2E\). Because
\(|m-RW|=O(q_0)\), the factor \(e^{\tau T_+}\) in the weight cancels this
volume loss up to \(O(q_0^2/W)=o(s)\).

Write \(M_j=\ex e^{-\tau|S_j|}\). On every dense row, we have
\(M_j=(1+o(n^{-1}))/(\tau\sqrt{\pi\lambda_j})\), which follows from the
normal comparison proved below together with the fact that
\(\tau\sqrt{\lambda_j}\) is exponentially large in \(n^\gamma\).
Substituting this identity into the product of the dense atom probability, the
ball volume, and the nondense tilted expectation, and using
\(m+m\log(RW/m)=RW+O(q_0^2/W)\), gives
\[
 \log\ex Y=2E\tau+\sum_{j<r/2}\log M_j+o(s).
\]

To evaluate the sum on the right, let \(G\) be a standard normal random
variable and define \(\phi(y)=\ex e^{-y|G|}\). We show that the sum of
\(\log M_j\) differs by \(o(1)\) from the sum of
\(\log\phi(\tau\sqrt{2\lambda_j})\). For rows with
\(\lambda_j<e^{n^{1/4}}\), both absolute logarithms are
\(O(\tau\sqrt{\lambda_j})\), which is negligible even after summing over all
such rows. On the remaining rows, we use the Poisson kernel with parameter
\(e^{-\tau}\) as a probability measure on frequencies, since its integer
Fourier coefficients are \(e^{-\tau|k|}\). Complement pairing allows us to
replace the Bernoulli counts by Poisson counts uniformly across rows, with a
relative error that is exponentially small in \(n\). The marginal
characteristic function localizes to the interval
\(|\theta|\le e^{n^\gamma}/\sqrt{\lambda_j}\). On this interval, the
characteristic function has a Gaussian quadratic exponent with uniform error
\(O(e^{4n^\gamma-n^{1/4}})\), and the kernel density equals
\(\tau/[\pi(\tau^2+\theta^2)]\) up to a relative factor of
\(1+O(\tau+\theta^2)\). Extending the quadratic integral from this interval to
the whole real line costs a superexponentially small amount, which is
negligible relative to its \(e^{-O(s)}\) main term. The extended Cauchy kernel
integral evaluates to exactly \(\phi(\tau\sqrt{2\lambda_j})\). Because all
errors are \(o(n^{-1})\) per row, their total across all rows is \(o(1)\).
The dense asymptotic formula for \(M_j\) also follows from substituting
\(t=y|G|\) in the normal integral for \(\phi\).

Now reindex the rows by \(u=r/2-j\ge1\) and set \(t_u=Cs-u^2/(2vn)\).
As in Section~\ref{sec:noiseless-exact}, uniformly in the moderate window of
rows, the fair-binomial relative entropy and Stirling prefactor give
\[
 \tau\sqrt{2\lambda_j}
 \le e\sqrt{2/\pi}\,e^{(t_u+\epsilon_n)/2},
 \qquad \epsilon_n=O(n^{-2/3}),
\]
with the constant determined by \(\kappa_{\rm noise}\); outside the moderate
window, Hoeffding's inequality gives negligible total intensity.
Write \(z=t_u+\epsilon_n\). Applying Jensen's inequality gives
\[
 \log\phi(e\sqrt{2/\pi}\,e^{z/2})
 \ge -(1+z/2)\mathbf1_{z>0}-b(z),
 \quad
 b(z)=
 \begin{cases}
 (2e/\pi)e^{z/2},&z\le0,\\
 (\pi/(2e^2))e^{-z},&z>0.
 \end{cases}
\]
The first branch follows directly from Jensen's inequality and
\(\ex |G|=\sqrt{2/\pi}\). For the second branch, we rewrite the normal
integral as \(\phi(y)=\sqrt{2/\pi}\,y^{-1}\ex e^{-X^2/(2y^2)}\) with \(X\) a
mean-one exponential random variable (whose second moment is two) and then
apply Jensen's inequality. Integrating \(b\) yields
\(\int b=4e/\pi+\pi/(2e^2)<3.68\). There are \(RW+O(1)\) levels with positive
\(z\), and the sum of the half-quadratic terms over those levels is
\(a_{r/n}(C)n/2-Cs/4+O(1)\). The shoulder Riemann-sum argument from
Section~\ref{sec:noiseless-exact} gives
\(\sum b(z)=(v/R)s\int b+o(s)\). Combining these estimates with \(2E\tau=RW\)
and Stirling's approximation for \(N\), we obtain
\[
 \log(N\ex Y)\ge (C/4-3.68v/R)s+o(s)\ge cs
\]
for a positive fixed constant \(c\). Indeed, evaluating at \(\rho=1/2\) shows
that \(C>3.3\), and since \(v/R\le1/\sqrt{8C}\), the coefficient of \(s\) is
strictly positive. Rounding the maximizing codimension to the nearest even
integer changes the extensive cancellation by only \(O(1)\).

\step{The relative joint bound.}
To finish the second moment proof, we bound the joint second moment across all
pairs of opposite-face candidates and show that the relative variance tends to
zero. We start with the typical pairs of
Section~\ref{sec:noiseless-face-proof}, defined using the signed-overlap cutoff
\(n^{-1/8}\). For prescribed dense targets on both candidates, we represent the
joint weight in frequency space using Haar Fourier measure on dense rows and
the Poisson-kernel measure on nondense rows. Complement pairing, uniform
Poissonization, and projecting out distant rows of total intensity \(o(1)\)
work just as in the previous proofs. On every retained row with
\(\lambda_j\ge e^{n^{1/4}}\), we apply the neutral singleton bound
\(\lambda_j e^{-O(n^{1/12}+\log n)}\) from the strict noisy case. This
localizes the integral to the core region
\(|\theta_j|\le e^{n^\gamma}/\sqrt{\lambda_j}\) on all such rows, because the
off-core contribution is negligible even compared to the \(e^{-O(n)}\) product
of single-core integrals.

On the remaining light rows, we keep the full fundamental frequency interval.
Cross terms involving any light row are \(o(1)\) uniformly by the shell-pair
intensity bound from Section~\ref{sec:noiseless-face-proof}: within a block,
the sum of the \(\lambda_j|\theta_j|\) terms over light rows is at most
\(\pi n e^{n^{1/4}}\), while the corresponding sum on the localized rows is at
most \(n\sqrt K e^{n^\gamma}\). Dividing their product by \(K\) and allowing
for the \(e^{o(s)}\) intensity correction shows that these cross terms vanish.
The remaining light-row marginal weights have product \(1+o(1)\). On the core,
the dense target phases are also \(1+o(1)\) uniformly, because the dense
targets satisfy \(\ell^1\le2E\) and the dense frequencies obey the core cutoff.

After absorbing a uniform \(o(1)\) quartic error on the core, we scale every
localized frequency by \(\sqrt{2\lambda_j}\). Under this rescaling, each
marginal probability density on the core is proportional to a standard Gaussian
times an even decreasing function of absolute value, coming from the core
cutoff and the Poisson kernel. The normalized coordinates are therefore
independent within each marginal and centered 1-subgaussian: by monotone
weighting, their \(\cosh(tx)\) expectations are bounded above by those of the
standard Gaussian. By the signed Walsh cross-Gram estimate, the normalized
cross block \(T_0\) has Frobenius norm \(o(1)\). For independent marginal
vectors \(U,V\) drawn from the rescaled core distributions, Jensen's
inequality and Gaussian linearization give
\[
 1\le\ex e^{-U^{\mathsf T}T_0V}
 \le\det(I-T_0T_0^{\mathsf T})^{-1/2}=1+o(1).
\]
To obtain the upper estimate, condition on \(U\) and apply subgaussianity to
\(V\); then represent the resulting factor \(e^{\|T_0^{\mathsf T}U\|^2/2}\) by
an auxiliary normal linear moment and apply subgaussianity to \(U\). Thus the
mixed joint atom is \(1+o(1)\) times the product of its single-candidate
integrals, uniformly over both dense targets. Summing over both dense balls,
using the untruncated tail weights as a joint upper bound, and applying the
relative single-candidate truncation estimate handles all typical pairs.

The remaining pair classes are all negligible relative to the squared mean.
Atypical fixed-contraction pairs satisfy the same shell bounds as before, now
incurring at most an \(\exp(o(s))\) weight loss per factor against an
\(\exp(-\Omega(n^{3/4}))\) count loss. For distinct near partners, we use the
two-shift forest of Section~\ref{sec:noiseless-exact} and condition on all
landmarks outside its cells. On the support of both weights, the concatenated
folded vector has \(\ell^1\) norm at most \(4E\). Summing along a fundamental
cut therefore confines each independent cell count to an interval of length at
most \(8E+1\). Using this interval bound in place of the single-atom factor
preserves the full two-half leading rate at \(A-2B=C\), and the number of
near-partner pairs keeps the same extensive entropy margin because
\(\ex Y=e^{-a_{r/n}(C)n/2+o(n)}\). Finally, identical pairs contribute at most
\((\sup Y)N\ex Y\), which is negligible because the log mean is positive on the
\(s\) scale and \(\sup Y=e^{o(s)}\). Hence the relative variance tends to zero,
proving failure.
\end{proof}

\section{Which binary faces cause failure above the cutoff?}
A face fixing a positive fraction of coordinates governs the first-moment
cutoff, but the residual failures above it come from faces fixing at most two
coordinates.

\begin{corollary}\label{cor:rare-faces}
Take the binary model and error budget of Theorem~\ref{thm:box}, strictly
above the cutoff, and write \(K=e^{A n^{1/3}}\). Let \(F\) be the event that
some pair of faces is confusable, and let \(Z\) count all unordered confusable
pairs. Then
\[
 \prob(F)\sim\frac{n(n-1)(4E+1)}{\sqrt{\pi K}},
 \qquad \ex Z\sim 3\prob(F).
\]
For each coordinate \(i\), let \(U_i=\sum_{v\in S}(-1)^{v_i}\).
If we discard a landmark sample whenever two entries of the list
\((U_1,-U_1,\ldots,U_n,-U_n)\) are at distance at most \(4E\),
the discard probability is \((1+o(1))\prob(F)\). Conditional on retaining
the sample,
\[
 \lim_{n\to\infty}n^{-1/3}\log\prob(F\mid\text{retained})=-D,
 \qquad D=A-2B.
\]
\end{corollary}

\begin{proof}
We show that strictly above the cutoff, failure is dominated by pairs of faces
whose distance histograms differ along a single nonloop edge, and that
filtering out these one-edge collisions drops the failure rate to the two-edge
scale. We first rule out pairs of unequal, large, or multi-edge codimension
by union bounds, then classify all one-edge pairs and compute their
probabilities and pairwise overlaps by Fourier analysis, and finally prove the
matching conditional lower bound using a fixed codimension-three opposite pair.

\step{Eliminating high codimensions and multi-edge pairs.}
Put \(s=n^{1/3}\) and \(x=(E+1)/\sqrt K=e^{-Ds/2+o(s)}\).
For pairs of faces with different codimensions, the empirical-mean bound of
Section~\ref{sec:sum-codim} makes the summed failure probability
superexponentially small on the \(s\) scale.
For equal codimension \(r\ge\varepsilon n\), the strict first-moment margin
above the cutoff gives an exponentially small summed probability on the \(n\)
scale.
For intermediate equal codimensions \(1\le r<\varepsilon n\), the robust forest
estimate of Section~\ref{sec:robust-forest} bounds the confusability
probability of a single pair by
\[
 \exp[-c\Lambda\min(r,\sqrt{r\Lambda})],\qquad \Lambda=Ds .
\]
Take \(\varepsilon>0\) small. For \(\Lambda<r<\varepsilon n\), the logarithm
of the number of pairs is at most \(2r\log(2en/r)\), which is at most half the
exponent in absolute value once \(\varepsilon\) is small enough; summing over
this range costs at most \(\exp[-\Omega(n^{2/3})]\). For \(r\le \Lambda\), the
same pair count is absorbed by half of the exponent \(c\Lambda r\). Therefore,
for any fixed \(M\), choosing a sufficiently large fixed \(R\) makes the total
contribution from all codimensions above \(R\) equal to \(o(e^{-Ms})\).

Now consider pairs of equal codimension \(r\le R\). Draw an edge between two
distance levels for each distinct attainable unordered nonloop pair of
distances from a landmark to the two faces. Any two distinct edges form a
forest. Each chosen oriented edge is realized by some assignment of bits to the
union of the two face supports, defining a cell of density at least
\(2^{-2R}\). Condition on all landmarks outside these disjoint cells. Whenever
the two faces are confusable, each of the two forest cuts has an imbalance of
absolute value at most \(2E\), because the total positive histogram discrepancy
is at most that value. The two-cut bound therefore gives probability
\(O_R(x^2)\) for each pair with at least two edges. Summing over all such pairs
with \(r\le R\) gives \(O_R(n^{2R}x^2)\), which is negligible against \(n^2x\).

\step{Classifying one-edge pairs.}
It remains to classify pairs of equal codimension that have only one nonloop
edge. A landmark lying in the first face but not in the second produces some
edge \(\{0,k\}\); complementing all bits of that landmark produces the edge
\(\{r,r-k\}\), so having only one edge forces \(k=r\). Let \(b\) be the number
of conflicts between the two faces on their common support, and let \(c\) be
the size of each face's exclusive support. As a landmark varies over the first
face, its distance to the second face ranges over all integers from \(b\) to
\(b+c\). If \(c>0\) and \(r\ge2\), this range includes an interior distance
level, which produces a second edge and contradicts the one-edge assumption.
If \(c=0\), the two faces share the same support and must prescribe opposite
bits on every coordinate of that support. Their attainable nonloop pairs are
then \(\{j,r-j\}\), so there are \(\lceil r/2\rceil\) such edges. Requiring
only one edge thus implies that the one-edge pairs are precisely the pairs of
distinct codimension-one faces and the opposite pairs of codimension two with
common support.

For two codimension-one faces on different coordinates, the imbalance in the
single nonloop edge is, up to sign, \((U_i\pm U_j)/2\). An opposite pair on
the two-coordinate support \(\{i,j\}\) has the exact same imbalance, by the
identity
\(\mathbf1_{\text{both match}}-\mathbf1_{\text{both mismatch}}
=(\varepsilon_i+\varepsilon_j)/2\),
where \(\varepsilon_i,\varepsilon_j\) are the prescribed signs multiplied by
the respective coordinate signs. Thus, for every \(i<j\), each of the two
events
\[
 \big| (U_i\pm U_j)/2\big|\le 2E
\]
corresponds simultaneously to three confusable pairs: two pairs of
single-coordinate faces and one opposite pair of codimension two. Separately,
single-coordinate opposites on the same coordinate give the \(n\) events
\(|U_i|\le2E\), which contribute \(O(nx)\) in total and are negligible
compared to \(n^2x\).

\step{Probabilities and pairwise overlaps.}
For any dominant event, the linear form is the difference of two independent
\(\Bin(2^{n-2},K2^{-n})\) random variables, which has variance \(K(1-p)/2\)
with \(p=K2^{-n}\). On the fundamental frequency interval, its characteristic
function is \(\big(1-2p(1-p)(1-\cos t)\big)^{2^{n-2}}\). After rescaling by
the standard deviation, this characteristic function tends to the standard
Gaussian characteristic function and is dominated by \(e^{-c't^2}\) on the
rescaled interval. Furthermore, the Fourier phase for atom \(k\) tends to one
uniformly for \(|k|\le2E=o(\sqrt K)\) on compact frequency sets. By Fourier
inversion and dominated convergence, each dominant event has probability
\((1+o(1))(4E+1)/\sqrt{\pi K}\).

To bound pairwise overlaps, note that any two distinct dominant forms are
linearly independent functions of at most four landmark bits. We can therefore
choose two bit assignments with linearly independent contribution vectors and
condition on all landmarks outside their two disjoint cells. For each
permissible output pair of the two forms, at most one tuple of cell counts
works. Because the binomial maximal atom on each cell is \(O(K^{-1/2})\), the
intersection of two distinct dominant events costs \(O(x^2)\). Bonferroni's
inequalities then give an aggregate overlap of \(O(n^4x^2)=o(n^2x)\), and the
asymptotic equivalents for \(\prob(F)\) and \(\ex Z\) follow.

\step{The filter and conditional failure rate.}
The filtering rule removes precisely the one-edge events, and it can be checked
by sorting the signed values and comparing adjacent differences. Its rejection
probability follows from the same union estimate as above. Conditional on
retaining the sample, no one-edge pair is confusable, so the bounds for pairs
with two or more edges give a conditional log rate at most \(-D\).

For the matching lower bound, fix a pair of opposite faces on a
three-coordinate support, and write \(W_0,W_1\) for the signed count
differences in the disjoint shell pairs \((0,3),(1,2)\). These are independent
symmetric binomial differences, so the same Fourier estimate shows that
\(|W_0|,|W_1|\le E\) has probability \(\Theta(x^2)\) and implies confusability.
Let \(f_0,f_1\) be the corresponding single-landmark contribution functions.
Then \(3f_0+f_1\) is the sum of the three prescribed coordinate signs, and
\(f_0-f_1\) is their product. By orthogonality of the cube characters, no
rejected form, which is a nonzero linear form supported on at most two
coordinates, belongs to their span. Hence, for every rejected form, we can
choose three bit assignments with linearly independent contribution vectors on
the union of the at most five involved coordinates. Conditioning outside these
three disjoint cells and applying three maximal atom bounds gives an
intersection probability of \(O(x^3)\) with this fixed collision. Summing over
all rejected forms yields a total overlap of \(O(n^2x^3)=o(x^2)\), proving the
matching conditional lower rate.
\end{proof}


\appendix
\section{Entropy lower bound}\label{app:entropy}
In this section we prove a lower bound on the number of landmarks needed for vertex, outer-vertex, or edge resolution. The strategy is to compare the Shannon entropy of a uniformly random query with an upper bound on the entropy of its distance histogram obtained by applying the geometric maximum-entropy bound to each shell count separately.

Write $p_0=(q-1)/q$ for the mismatch probability between two independent uniform letters, and set $\sigma^2=p_0(1-p_0)$. Choose a query $X$ uniformly among all vertices or among all edges, and let $Y_r$ count the landmarks at distance $r$ from $X$. For any fixed set of $k$ landmarks, transitivity of the Hamming graph and linearity of expectation give $\mathbb EY_r=k\pi_r$, where $\pi_r$ is the distance distribution of a single fixed landmark from $X$. For vertex queries, $\pi_r$ is simply the law of $\mathrm{Bin}(n,p_0)$. For edge queries, conditional on the active coordinate of the edge, the inactive letters are uniform and independent of the uniform unordered pair of active letters. Consequently, the distance from a fixed landmark to a uniform edge is the independent sum of $\mathrm{Bin}(n-1,p_0)$ and $\mathrm{Bernoulli}(1-2/q)$.

Any nonnegative integer-valued random variable with mean $u$ has Shannon entropy at most $h(u)=(u+1)\log(u+1)-u\log u$ (with $h(0)=0$), because its relative entropy with respect to the geometric distribution of mean $u$ is nonnegative. Applying subadditivity of entropy across shells and bounding $h(u)$ from above gives
\[
 H(Y)\le\sum_r h(k\pi_r),\qquad h(u)\le (\log u)_++2\quad (u>0).
\]
For small means $0<u\le1$, we also use the sharper bound $h(u)\le u(2+\log(1/u))$. Note that subadditivity requires no independence among the coordinates of the distance histogram.

To bound the landmark count from below, it suffices to pass to subsequences along which $n^{-1/3}\log k\to\lambda<\infty$, which in particular implies $k/q^n\to0$. When the landmark set resolves all vertices, the histogram determines $X$ uniquely, so $H(Y)=n\log q$; similarly, edge resolution gives $H(Y)=\log(n\binom q2 q^{n-1})$. For outer vertex resolution, conditioning on whether $X$ belongs to the landmark set yields $H(Y)\ge(1-k/q^n)\log(q^n-k)$. In all three cases we have $H(Y)\ge n\log q-o(n)$.

We split the sum over shells into a moderate-deviation window around the mean and its tail. For any fixed $B$, Stirling's formula and the Taylor expansion of the binomial relative entropy give, uniformly over $|r-np_0|\le Bn^{2/3}$,
\[
 \log\pi_r=-\frac{(r-np_0)^2}{2\sigma^2n}+O_{q,B}(\log n).
\]
On this window the cubic Taylor remainder in the exponent is bounded. Replacing $n$ by $n-1$ and applying bounded shifts preserve this asymptotic estimate, as does taking a mixture of the resulting atoms with fixed positive weights (for binary edges, only one weight is nonzero). Outside the window, choose $B$ large enough that Hoeffding's inequality yields $\pi_r\le\exp[-(\lambda+1)n^{1/3}]$ for all sufficiently large $n$. Then $k\pi_r\le\exp[-n^{1/3}/2]$ on the tail, so monotonicity of $h$ together with the small-mean bound shows that all tail shells combined contribute $o(n)$ to the entropy.

Inside the window, the additive constant $2$ in the upper bound for $h(u)$ contributes $O(n^{2/3})$ across all shells in the window. Rescaling the shell index via $x=(r-np_0)/n^{2/3}$ converts the remaining sum into a Riemann sum, yielding
\[
 \log q\le\limsup H(Y)/n\le
 \int_{\mathbb R}\!\left(\lambda-\frac{x^2}{2\sigma^2}\right)_+dx
 =\tfrac43\sqrt{2\sigma^2}\,\lambda^{3/2}.
\]
Rearranging this inequality gives $\lambda\ge C_q/4^{2/3}$, as claimed.

\section{Collision witnesses}
We now introduce a general method for bounding the probability that two distinct queries produce the same distance histogram in the independent Bernoulli landmark model with inclusion probability $p=q^{-n}\exp(A n^{1/3})$. The idea is to find disjoint classes of candidate landmarks whose distance pairs form a forest; conditioning on all selections outside these classes then leaves at most one valid tuple of landmark counts inside them, allowing us to multiply single-binomial maximum-atom bounds.

Observe that a class of words at distances $r,s$ from the first and second queries, respectively, contributes its number of selected landmarks multiplied by $e_r-e_s$ to the difference between the two histograms. Suppose we select disjoint classes whose unordered distance pairs form a forest with no repeated edges. By repeatedly deleting leaves, the corresponding oriented incidence vectors are linearly independent. Conditional on all landmark selections outside these classes, a histogram collision therefore admits at most one tuple of counts for the remaining independent binomial variables. This conclusion holds even when each chosen class is only a proper subset of a joint-distance cell.

To bound the probability of hitting that unique tuple, we use the following maximum-atom estimate: for a class of size $L$, the maximum point probability of its selected landmark count is at most
\[
 \min\{1,C[Lp(1-p)]^{-1/2}\}.
\]
Indeed, Fourier inversion bounds the maximum atom by the normalized integral of $|1-p+pe^{it}|^L$ over $[-\pi,\pi]$. Writing the squared modulus as $1-4p(1-p)\sin^2(t/2)$ and applying the inequality $\sin(|t|/2)\ge |t|/\pi$ on $[-\pi,\pi]$ gives the desired bound.

\subsection{Paths and matchings}
We next construct explicit forest families of collision witnesses for two centers on a coordinate block of length $N\in\{n,n-1,n-2\}$, allowing additive distance offsets $t_1,t_2\in\{0,1,2\}$. Assume that either the two centers differ or the two offsets differ. Depending on whether the centers differ in at most half or more than half of the coordinates, we obtain either a collection of disjoint paths or a matching.

\step{Coordinate decomposition and binomial estimates.}
Let $d$ be the number of coordinates where the two centers differ, and let $m=N-d$ be the number of coordinates where they agree. On the agreeing coordinates, we take words of Hamming weight $i$ relative to the common center; there are $\binom mi(q-1)^i$ such words. On the differing coordinates, let $a,b$ count the positions matching the first and second centers, respectively, and let $c=d-a-b$ count the positions matching neither center. There are $d!(q-2)^c/(a!b!c!)$ such choices on the differing block, and each resulting word has unshifted distance pair $(i+b+c,i+a+c)$. When $q=2$, no letter differs from both centers, so we set $c=0$ and use the convention $0^0=1$.

To evaluate the sizes of these classes, we use two basic facts about binomial probabilities on $M=\Theta(n)$ trials with a fixed success parameter. First, any atom within $O(n^{2/3})$ of the mean has probability $\exp[-O(n^{1/3})]$: Stirling's formula expresses its logarithm as $-M D(j/M\Vert\pi)+O(\log n)$, and the relative entropy satisfies $D(\pi+v\Vert\pi)=O(v^2)$ near $0<v+\pi<1$ as $v\to0$. Second, fixing a binomial variable at a mode loses at most a factor of $M+1$, and shifting a bounded number of steps away from a mode (while remaining in the support) loses only an additional constant factor, as seen from ratios of adjacent binomial probabilities.

\step{Path construction for small distance.}
Suppose $m\ge N/2$ and $d\ge1$. We fix the differing-block counts and vary the agreeing-block weight $i$. Specifically, fix $u=a+b$ at a mode of $\mathrm{Bin}(d,2/q)$, taking $u=d$ when $q=2$. If that mode happens to be zero (which occurs only for bounded $d$), replace it by $1$. Next choose $a\ne b$ summing to $u$ as close to equal as possible, and order them so that $b-a+t_1-t_2\ne0$. The number of differing-block choices equals $q^d$ times the point probabilities at $u$ and $a$ of $\mathrm{Bin}(d,2/q)$ and $\mathrm{Bin}(u,1/2)$, respectively, and is therefore at least $q^d/n^{O_q(1)}$. Now let $i$ range over consecutive integers in the moderate-deviation window around $mp_0$. Every resulting edge joins two levels separated by the same nonzero shift, so the edges decompose into vertex-disjoint paths within residue classes and hence form a forest. When $d=0$, the offsets must differ, and the same path construction goes through directly without any choice on differing coordinates.

\step{Matching construction for large distance.}
Suppose instead that $d>N/2$. Here we fix $i$ at a mode of the agreeing-block weight (taking zero if $m=0$) and fix $u$ at the mode of $\mathrm{Bin}(d,2/q)$ as above, which is now of order $n$. Let $a$ vary over consecutive integers below $\lfloor u/2\rfloor-2$ that lie within $O(n^{2/3})$ of $u/2$, and set $b=u-a$. The two shifted distance endpoints then lie on opposite sides of a common fixed midpoint. As $a$ increases, the lower endpoint increases and the upper endpoint decreases, so all endpoints are distinct and the edges form a matching. In this construction the varying binomial atom comes from $\mathrm{Bin}(u,1/2)$, while fixing $i$ and $u$ incurs only polynomial factors.

\subsection{Reduction of edges to shifted distances}
To apply the path and matching constructions to pairs of edge queries, we fix the landmark letters on the active coordinates of both edges. Doing so converts each edge distance into a vertex distance on the remaining coordinates plus an additive offset: fixing a landmark letter on an edge's active coordinate adds offset zero if that letter lies in the edge's active pair, and offset one otherwise.

For two parallel edges, they share the same active coordinate. If their words on the inactive coordinates differ, any fixed letter on the active coordinate leaves distinct centers on the remaining block. If their inactive words coincide, the two edges must have distinct active pairs, so we choose a letter in the symmetric difference of those pairs to ensure that the two offsets differ.

For two nonparallel edges, we fix the landmark letters on both active coordinates, leaving a block of length $N=n-2$. If the two restricted centers differ on these $N$ coordinates, any choice of the two active letters suffices. Otherwise, the restricted centers agree, and we examine the first active coordinate, where the first edge has active pair $P$ and the second edge has a fixed letter $\beta$. The contribution of a landmark letter $\ell$ at this coordinate to the difference between the two distances is $\mathbf{1}_{\ell\notin P}-\mathbf{1}_{\ell\ne\beta}$. If $\beta\in P$, comparing $\ell=\beta$ with the other letter in $P$ yields different values; if $\beta\notin P$, comparing $\ell=\beta$ with either letter in $P$ yields different values. Thus this contribution is nonconstant in every case. Fixing the other active coordinate arbitrarily, we can choose $\ell$ so that the total offset difference is nonzero. In all cases the resulting offsets are at most two, and fixing at most two coordinates reduces shell sizes by at most a factor of $q^2$.

\section{Sparse construction and upper bound}\label{app:sparse-upper}
We now combine the forest witnesses with the maximum-atom bound to upper-bound the expected number of vertex and edge collisions, proving the upper bound in Theorem~1.

\subsection{The shell integral}
First we translate a forest of witness classes covering almost all shells in a moderate-deviation window into an exponential bound on the collision probability. Suppose that every first-query level $r$ in a fixed rescaled window, with at most $o(n^{2/3})$ exceptions, contributes a forest class of size $L_r$ that is at least its marginal shell size divided by a polynomial in $n$. When $N=n-O(1)$ and $x=(r-Np_0)/n^{2/3}$, bounded additive offsets do not affect the uniform lower bound
\[
 \log(pL_r)\ge [A-x^2/(2\sigma^2)]n^{1/3}-O(\log n).
\]
By the maximum-atom estimate, the logarithm of the maximum atom for class $r$ is at most $-\frac12(\log(pL_r))_++O(1)$. Multiplying these bounds across the forest and taking the rescaled window wider than the support of the positive part gives a log collision probability of at most $-a_q(A)n+o(n)$, where $a_q(A)=\frac23\sqrt{2\sigma^2}A^{3/2}$.

Two variations of this integral estimate will be needed for the path and matching constructions. If the varying shell lies on a subblock of length $m=\Theta(n)$ while the fixed block incurs only polynomial factors, the quadratic term in the integrand becomes $x^2/(2\sigma^2 m/n)$, which scales the exponential rate to $\sqrt{m/n}\,a_q(A)$. If we only sum over a one-sided window terminating at the mean, the rate is halved. Both estimates hold uniformly whenever $m/n$ is bounded away from zero.

\subsection{A radially descending forest}
When the Hamming distance $d$ between the two centers is bounded away from $0$ (and, for $q=2$, also bounded away from $N$), we obtain the full shell-integral rate $a_q(A)$ by building a forest in which every edge moves strictly closer to the mean shell radius.

Condition on the unshifted distance $r$ from the first center. The $r$ mismatches occupy a uniformly random $r$-subset of the $N$ coordinates, and their letters are independent and uniform among the $q-1$ alternatives to the first center. First fix a mode $j$ of the number of these mismatches that fall on the $d$ differing coordinates, and then fix a mode $b$ of the number of those $j$ mismatches that match the second center. The distribution of $j$ is hypergeometric, and conditional on $j$, the distribution of $b$ is $\mathrm{Bin}(j,1/(q-1))$. Each modal probability is therefore at least $1/(N+1)$. Moreover, the hypergeometric mode $\lfloor(d+1)(r+1)/(N+2)\rfloor$ differs from $dr/N$ by $O(1)$, and $b=j/(q-1)+O(1)$, which also covers the deterministic binary case. Since the unshifted distance to the second center is $s=r+d-j-b$, setting $z=Np_0$ yields
\[
 s-z=\left(1-\frac{qd}{(q-1)N}\right)(r-z)+O(1).
\]
Adding bounded offsets changes neither the coefficient nor the $O(1)$ error.

Whenever the contraction coefficient stays inside a compact subinterval of $(-1,1)$, we pick one such modal class for every first-query level $r$ outside a bounded central window around $z$, and direct its edge from the endpoint of larger radius about $z$ to the endpoint of smaller radius. Because each level contributes at most one outgoing edge, no edge is repeated and every vertex has outdegree at most one. Any undirected cycle would have to contain a vertex of maximum radius about $z$, from which two edges would be directed outward to strictly smaller radii, contradicting outdegree at most one. Thus the chosen classes form a forest and achieve the full shell integral. For any fixed small $\delta>0$, this argument covers $d\ge\delta N$ when $q\ge3$, and $\delta N\le d\le(1-\delta)N$ when $q=2$.

\subsection{Endpoint ranges}
It remains to handle pairs whose centers are nearly identical or, in the binary case, nearly antipodal, and then sum the resulting collision probabilities over all pairs.

When $d<\delta N$, we apply the fixed-shift path construction on the $m=N-d$ agreeing coordinates across the entire moderate-deviation window. By the subblock shell-integral bound, this achieves an exponential rate of at least $\sqrt{1-\delta}\,a_q(A)+o(1)$. Given $A>C_q$, we choose $\delta$ small enough that this rate exceeds $2\log q$ in the limit, which beats even the crude total pair count $O_q(n^2q^{2n})$.

When $q=2$ and $d>(1-\delta)N$, we apply the differing-block matching with $u=d$ over one half of the moderate-deviation window, obtaining a rate of at least $\frac12\sqrt{1-\delta}\,a_2(A)+o(1)$. To count such pairs, note that once the first binary word is chosen, the second must lie in the Hamming ball of radius $\delta N$ around its complement. For edge queries, specifying the active directions and fixed letters contributes only polynomial factors. Hence the number of pairs in this regime is at most $\exp[(\log2+h_2(\delta)+o(1))n]$, where $h_2(t)=-t\log t-(1-t)\log(1-t)$. Because $a_2(A)/2>\log2$, shrinking $\delta$ if necessary ensures that the matching rate dominates this pair count as well.

Summing the collision probabilities over all vertex pairs and all edge pairs shows that the combined expected number of vertex and edge collisions tends to zero. Furthermore, the total number of landmarks is a binomial random variable whose variance is at most its diverging mean; with probability tending to one, it is therefore nonzero and at most twice its mean. Since full vertex resolution implies outer resolution, letting $A\downarrow C_q$ completes the proof of the upper bound in Theorem~1.

\section{Pairwise collision rates and the first moment}
In this and the next two sections we sharpen the pairwise collision bounds, establish a matching lower bound, and use them to determine the exact exponential decay rates of failure probabilities. Throughout these three sections, set
\[
 L=\log q,\quad \mu=n(1-1/q),\quad \sigma^2=(q-1)/q^2,\quad
 p=q^{-n}e^{A n^{1/3}},\quad a=\tfrac23\sqrt{2\sigma^2}\,A^{3/2},
\]
for fixed $q\ge2$ and $A>0$. Write $\pi_r$ for the distance distribution of a uniformly random word from a fixed vertex, and let $\lambda_r=pq^n\pi_r$ be the expected number of landmarks in the radius-$r$ shell. For a sufficiently large fixed constant $B$, define the moderate-deviation window $W=\{r:|r-\mu|\le Bn^{2/3}\}$.

By a uniform Stirling expansion on $W$, the expected shell sizes satisfy
\[
 \log\lambda_r=A n^{1/3}-\frac{(r-\mu)^2}{2\sigma^2n}+O_{q,B}(\log n),
\]
where the cubic Taylor remainder is bounded on $W$. Approximating the sum over $W$ by a Riemann sum with mesh $n^{-2/3}$ and integrand $\tfrac12(A-u^2/(2\sigma^2))_+$ yields
\[
 \tfrac12\sum_{r\in W}(\log\lambda_r)_+=an+O(E_n),
 \qquad E_n=n^{2/3}\log n.
\]
Outside $W$, increasing $B$ if necessary allows Hoeffding's inequality to give the tail bound $\sum_{r\notin W}\lambda_r=o(1)$.

\subsection{Upper bounds for nonantipodal pairs}
We first show that every nonantipodal pair of distinct vertices has collision probability at most $\exp(-an+o(n))$, uniformly over all such pairs.

For $q\ge3$, we split at query distance $d=n^{5/6}$. When $d\ge n^{5/6}$, we use the radially descending forest from the sparse construction: the conditional mean contracts deviations from $\mu$ by an absolute factor of at most $1-\Omega_q(n^{-1/6})$, so bounded rounding omits only $O_q(n^{1/6})$ central levels, which is negligible on the window scale. When $d<n^{5/6}$, we instead use the full-window fixed-shift paths on the common block of length $n-d=n-o(n)$. This gives an exponential rate of $a-o(1)$ uniformly, while fixing a mode on the differing block incurs only a polynomial factor per chosen cell.

For binary words ($q=2$) with $1\le d<n$, write the first distance as $r=n/2+x$, where $r\in W$ and $x$ is positive and exceeds a fixed constant. Suppose first that $d\le n/2$. Conditional on $r$, the number $j$ of mismatches on the differing block is hypergeometric with support $[\max(0,r-(n-d)),\min(d,r)]\cap\mathbb Z$ and mode $dr/n+O(1)$. We choose $j$ to be the larger of a mode and $\lfloor d/2\rfloor+1$. This value lies in the support (even when $d=1$) and sits within a bounded distance of a mode. Because the ratio of any two adjacent positive hypergeometric probabilities is bounded above and below by powers of $n$, the chosen atom accounts for at least $n^{-O(1)}$ of the shell. Moreover, since
\[
 0<2j-d\le 2dx/n+O(1)\le x+O(1)<2x,
\]
the resulting second distance $s=r+d-2j$ has strictly smaller absolute deviation from $n/2$ than $r$ does.

When $d>n/2$, we apply the same selection to the $m=n-d$ agreeing positions: choosing $i>m/2$ mismatches on the agreeing block produces $s=n-r+2i-m$, which again has strictly smaller absolute deviation from $n/2$. For negative deviations $x<0$, complementing all landmark bits reflects the construction symmetrically. Consequently, every retained level has a directed arc to a level of strictly smaller deviation, forming a forest. Applying the forest maximum-atom bound yields a collision probability of at most $\exp(-an+o(n))$ uniformly over all nonantipodal pairs.

\subsection{An involution lower bound}
We now prove a matching lower bound on the collision probability of any pair of distinct vertices $x,y$. The idea is to pair up words using a distance-swapping involution so that the histogram difference takes the symmetric form $X-X'$, restrict $X$ to a typical integer box, and apply Cauchy--Schwarz.

Consider the involution on the Hamming cube that exchanges the letters of $x$ and $y$ on every coordinate where they differ and fixes all other letters. This map exchanges $x$ and $y$. Any word fixed by the involution is equidistant from $x$ and $y$ and therefore contributes nothing to the histogram discrepancy. From each nontrivial two-element orbit, pick one representative $z$; its contribution vector to the discrepancy is $e_{d(x,z)}-e_{d(y,z)}$, while its partner in the orbit contributes the negative of this vector. The full discrepancy is therefore distributed as $X-X'$, where $X$ and $X'$ are independent identically distributed Bernoulli sums over the chosen representatives.

Each coordinate $r$ of $X$ receives contributions from at most the two radius-$r$ shells around $x$ and $y$, so its variance is at most $2\lambda_r$. By the shell tail bound outside $W$, with probability $1-o(1)$ no selected representative has distance outside $W$ from either query. Inside $W$, Chebyshev's inequality and a union bound show that, with probability $1-o(1)$, every coordinate of $X$ lies within $n^2\max(1,\sqrt{2\lambda_r})$ of its mean. The resulting integer box, with zeros outside $W$, has log cardinality at most $an+O(E_n)$ by the shell sum. Applying Cauchy--Schwarz to the probability distribution of $X$ on this box gives
\[
 \mathbb P(X=X')\ge \exp(-an-O(E_n)).
\]
Note that this argument makes no assumption of independence among the coordinates of $X$. To require in addition that both queries $x,y$ are absent from the landmark set costs only a factor $(1-p)^2$: simply delete their two-element orbit and repeat the exact same argument, whose variance and box bounds only improve.

\subsection{Antipodes and counting}
In the binary cube ($q=2$), antipodal pairs have a different collision rate because their distance shells pair up in reverse order. We compute this rate and combine it with the nonantipodal bounds to obtain the first-moment exponent for the number of vertex collisions.

For a binary antipodal pair, distance $r$ from one vertex is distance $n-r$ from the other. A collision therefore occurs if and only if opposite shell counts coincide, independently for each $r<n/2$; when $n$ is even, the middle shell imposes no constraint. Consider independent counts $U,V\sim\mathrm{Bin}(D,p)$ for a pair of opposite shells. When $\lambda=Dp\ge1$, the maximum-atom estimate gives the upper bound $\mathbb P(U=V)=O(\lambda^{-1/2})$; conversely, Chebyshev's inequality places constant probability mass on an interval of $O(\sqrt\lambda)$ integers, so Cauchy--Schwarz gives the matching lower bound up to a constant factor. When $\lambda<1$, the event that both counts are zero bounds the log probability below by $-O(\lambda)$, since $p=o(1)$. Summing over the half-window $r\in W$ with $r<n/2$ and controlling the tail outside $W$ yields
\[
 \log\mathbb P(\text{antipodal collision})=-an/2+O(E_n)
 =-(\sqrt2/6)A^{3/2}n+O(E_n).
\]
For the outer-vertex version, we additionally fix both endpoint counts at zero, which changes the log probability by at most $O(1)$.

There are $q^n(q^n-1)/2$ unordered pairs of distinct vertices. Removing the $2^{n-1}$ binary antipodal pairs does not affect the leading log pair count $2Ln$. Combining the pair counts with the pairwise collision probabilities shows that the first-moment exponent for collisions is $2L-a$ when $q\ge3$ and $\max(2L-a,L-a/2)$ when $q=2$. This first-moment calculation alone does not control dependence among pairs.

\section{A converse for Bernoulli landmarks}
In the strongly subcritical regime $a<L$, we show that the number of unordered collisions for vertices, outer vertices, and edges all have exponent $2L-a$ in probability. The lower bound comes from a pigeonhole argument on a common box of typical histograms, and the upper bound follows from Markov's inequality.

Fix a vertex or an edge and write $M_r$ for the number of landmarks at distance $r$ from it, using the corresponding shell law $\pi_r$ and mean $\lambda_r=pq^n\pi_r$. For an edge, the distance law is a binomial on $n-1$ coordinates plus an independent bounded Bernoulli offset, so the same Stirling expansion and tail bound hold with the mean shifted by $O(1)$. Applying Markov's inequality outside $W$ and Chebyshev's inequality inside $W$ shows that, with probability $1-o(1)$, the histogram is supported on $W$ and satisfies
\[
 |M_r-\lambda_r|\le n^2\max(1,\sqrt{\lambda_r})\quad(r\in W).
\]

For each query type, the common box of histograms satisfying these bounds has cardinality $K\le\exp(an+O(E_n))$. If there are $T$ queries of that type, then $\log T=Ln+o(n)$, and the expected number of atypical queries is $o(T)$; hence at least $T/2$ queries are typical with high probability. For outer vertices, removing the landmark set discards at most $|S|=o(q^n)$ vertices with high probability by Markov's inequality. A finite union bound handles all three query types simultaneously. When $a<L$, the $T'\ge T/3$ remaining typical queries map into at most $K$ histograms, so Cauchy--Schwarz guarantees at least $(T'^2/K-T')/2$ unordered collisions. This yields the lower collision-count exponent $2L-a$ in probability.

For vertices and outer vertices, the matching upper bound follows directly from Markov's inequality and the first-moment calculation above. For edges, fix $0<\delta<1/8$ and apply the projected-forest bounds: bulk projected pairs have collision probability at most $\exp((-a+o(1))n)$ and small-difference projected pairs have probability at most $\exp((-\sqrt{1-\delta}\,a+o(1))n)$, with at most $O_q(n^2q^{2n})$ pairs of each kind. When $q=2$, the additional near-antipodal regime has collision probability at most $\exp((-\sqrt{1-\delta}\,a/2+o(1))n)$ and at most $\exp((L+h(\delta)+o(1))n)$ pairs, where $h$ denotes binary entropy in nats. Taking the limsup in $n$ and then letting $\delta\downarrow0$ shows that the expected collision count has upper exponent $2L-a$ for $a<L$, and Markov's inequality completes the upper bound. This initial argument covers $a<L$; Section~\ref{sec:classical-birthday} handles the full subcritical range.

\section{Exponential failure probabilities and binary parity}
We now determine the exact exponential decay rate of the failure probability for each resolution notion. Write $F_V,F_O,F_E$ for the events that vertex, outer-vertex, and edge resolution fail, respectively. For each such event $F$, we prove that
\[
 \lim_{n\to\infty} n^{-1}\log\mathbb P(F)=
 \begin{cases}\min(0,2L-a),&q\ge3,\\ \min(0,L-a/2),&q=2.\end{cases}
\]
By sandwiching between a single event and the sum of the three probabilities, the same limit holds for the union of the three failure events, though not necessarily for their intersection. In the subcritical range $a<2L$, the second-moment arguments in this section show that each individual failure probability is at least $\exp(-O_{q,A}(E_n))$; the stronger high-probability failure statement is proved separately in Section~\ref{sec:classical-birthday}.

\subsection{Balanced joint cells}
To control the second moment of the number of colliding pairs, we need a multi-center analogue of the modal cell construction. Given $k\le4$ centers and a chosen anchor among them, we show that for every anchor distance $r\in W$ there is a single joint-distance cell of inverse-polynomial relative size in which the distances to all other centers stay within $O_{q,k}(1)$ of their conditional means.

Condition a uniform word to lie at distance $r\in W$ from the anchor. For another center at distance $d$ from the anchor, an agreeing coordinate mismatches that center with probability $r/n$, whereas a differing coordinate mismatches it with probability $1-r/((q-1)n)$. Summing over coordinates gives the conditional mean distance
\[
 \mu+\rho(d)(r-\mu),\qquad \rho(d)=1-qd/((q-1)n).
\]

We claim that there exists a joint-distance cell at anchor distance $r$ whose size is at least $q^n\pi_r/(n+1)^{O_{q,k}(1)}$ and in which every other center distance lies within $O_{q,k}(1)$ of its conditional mean. To construct it, partition the $n$ coordinates according to their $k$-tuples of letters across the $k$ centers into at most $q^k$ blocks of sizes $n_j$. Round each expected block mismatch count $rn_j/n$ up or down to an integer $l_j$ so that the counts $l_j$ sum to $r$. Each $l_j$ lies within two places of a mode of $\mathrm{Bin}(n_j,r/n)$ and therefore has inverse-polynomial probability, since the modal mass is at least $1/(n+1)$ and adjacent positive binomial masses have polynomial ratios. Conditional on the total number of mismatches being $r$, the joint probability of the block counts $l_j$ is the product of these binomial probabilities divided by $\mathbb P(\mathrm{Bin}(n,r/n)=r)$, which is again inverse polynomial. Finally, within each block distribute the $l_j$ mismatches as evenly as possible (up to rounding) among the $q-1$ nonanchor symbols. This forms a modal multinomial composition and hence has probability at least $(n+1)^{-(q-1)}$. Because the block and symbol counts differ from their conditional means by bounded amounts, all resulting center distances also differ from their conditional means by $O_{q,k}(1)$.

\subsection{Two-pair estimate}
We now apply the second moment method to a subfamily of outer vertex pairs at distance close to $\mu$. This yields the lower bound for vertex and outer-vertex failure when $q\ge3$, and for the even-distance nonantipodal subfamily when $q=2$.

Restrict attention to unordered vertex pairs whose Hamming distance $d_n$ lies within two of $\mu$, taking $d_n$ to be even when $q=2$. The number of such candidate pairs is $q^n\binom n{d_n}(q-1)^{d_n}/2=\exp(2Ln-O(\log n))$. Let $Z$ count the colliding candidate pairs for which both endpoints lie outside the landmark set $S$. By the involution lower bound, each candidate pair contributes at least $\exp(-an-O(E_n))$ to $\mathbb EZ$.

\step{Separating pairs of candidate pairs.}
To bound $\mathbb EZ^2$, measure distance between endpoints using $\Delta=d_H$ for $q\ge3$, and the quotient distance modulo complement $\Delta(u,v)=\min(d_H(u,v),n-d_H(u,v))$ for $q=2$. Fix $0<\delta<\min((q-1)/(4q),1/8)$. Given two candidate pairs, draw a cross edge between an endpoint of the first pair and an endpoint of the second whenever their $\Delta$-distance satisfies $\Delta\le\delta n$. Because the two endpoints within each candidate pair are separated by more than $2\delta n$ in $\Delta$, no endpoint can be incident to two cross edges; thus the cross edges form a matching, including any repeated or complemented endpoints across the two pairs.

\step{Joint upper bound for separated pairs.}
Suppose first that at most one cross edge joins the two candidate pairs. Then each pair has at least one unattached endpoint, which we take as its anchor. Every comparison between an anchor and the other three endpoints satisfies $|\rho|\le1-\eta$ for some fixed $\eta=\eta_{q,\delta}>0$. For each anchor and each $r\in W$ satisfying $|r-\mu|>K_q/\eta$ (where $K_q$ is a sufficiently large rounding constant), select the balanced joint cell constructed in the previous subsection. By our choice of $K_q$, all three other distances in this cell have strictly smaller absolute deviation from $\mu$ than $r$ does.

Orient each pair with its anchor positive and stack the two discrepancy vectors into a single two-block vector. Every chosen cell's column has a unique coordinate of maximum deviation from $\mu$, namely $+1$ in its own anchor's block at shell $r$. These leading coordinates are pairwise distinct across all chosen cells. In any putative linear relation among the columns, examining the largest-deviation layer that appears shows that its coefficient must vanish, proving linear independence. The cells are also pairwise disjoint, because each cell's unique leading anchor and radius identify it. Conditional on all landmark selections outside these cells, a simultaneous collision of both pairs therefore determines at most one vector of independent binomial cell counts. Multiplying their maximum-atom bounds, each block pays half the positive log-shell sum, with $O(\log n)$ loss per level from the polynomial cell factors and bounded losses per omitted central level on the $n^{1/3}$ scale. The joint collision probability for a separated pair of candidate pairs is therefore at most $\exp(-2an+O_\delta(E_n))$. Applying the identical single-anchor argument to a single candidate pair gives an upper bound of $\exp(-an+O(E_n))$.

\step{Close pairs and the second moment.}
Suppose instead that two cross edges join the two candidate pairs. Then the second pair lies inside two specified $\Delta$-balls of radius $\delta n$. Each such ball has cardinality at most $2\sum_{i\le\delta n}\binom ni(q-1)^i\le\exp(b_q(\delta)n+O(\log n))$, where $b_q(\delta)=h(\delta)+\delta\log(q-1)$. Consequently, there are at most $\exp((2L+2b_q(\delta))n+O(\log n))$ ordered pairs of candidate pairs of this kind, including identical pairs. We bound the joint collision probability of each such pair by the single-pair upper bound.

Setting $t=2L-a$, combining the separated and close contributions gives the moment bounds
\begin{align*}
 \mathbb EZ&\ge e^{tn-O(E_n)},\\
 \mathbb EZ^2&\le e^{2tn+O_\delta(E_n)}+e^{(t+2b_q(\delta))n+O_\delta(E_n)}.
\end{align*}
By Cauchy--Schwarz, $\mathbb P(Z>0)$ has lower exponent at least $\min(0,t-2b_q(\delta))$; letting $\delta\downarrow0$ gives the lower bound $\min(0,t)$. Markov's inequality gives the matching upper exponent for this subfamily. When $t>0$, fixing $\delta$ small enough that $2b_q(\delta)<t$ yields the quantitative subexponential lower bound on $\mathbb P(Z>0)$. Combining this with the all-pairs upper bound completes the proof of both vertex failure formulas for $q\ge3$. For binary words, it establishes the exponent $\min(0,t)$ for the even-distance nonantipodal subfamily.

\subsection{The binary antipodal family}
We next apply the same second-moment framework to binary antipodal pairs. Let $Z_*$ count outer antipodal collisions. From the antipodal calculation above, each antipodal pair collides with probability $\exp(-an/2+O(E_n))$.

Consider two antipodal classes whose $\Delta$-distance exceeds $\delta n$. Choose one representative from each class and repeat the balanced-cell construction with each representative as anchor, taking only levels $r<n/2$ outside the bounded central omission around $n/2$. Fold each pair's discrepancy block about $n/2$: a word at distance $s<n/2$ from the representative contributes $+1$ at coordinate $s$, a word at distance $n-s$ contributes $-1$ at coordinate $s$, and the middle shell contributes zero. Because the other representative's folded layer lies strictly closer to $n/2$, the leading coordinates of all chosen columns are distinct, and the half-window joint collision estimate gives at most $\exp(-an+O_\delta(E_n))$.

Close antipodal classes (within $\delta n$ in $\Delta$) contribute at most $\exp((L+h(\delta))n+O(\log n))$ ordered pairs, each bounded by the single-pair probability. Running the preceding second-moment argument with $t'=L-a/2$ and $h(\delta)$ in place of $2b_q(\delta)$ proves that $\mathbb P(Z_*>0)$ has exponential rate $\min(0,t')$ regardless of the parity of $n$. Together with the all-pairs upper bound, this establishes both the vertex and outer-vertex failure formulas for $q=2$.

\subsection{Parity restrictions}
In the binary cube, the parities of the landmark weights impose algebraic constraints on possible collisions. We record these constraints here so we can condition on them in the final subsection.

Partition the landmark set $S$ by Hamming weight parity: set $S_b=\{s\in S:|s|\equiv b\pmod2\}$ and $m_b=|S_b|$. Any odd-distance collision exchanges the two distance-parity classes and therefore requires $m_0=m_1$. Since $m_0$ and $m_1$ are independent $\mathrm{Bin}(2^{n-1},p)$ random variables, the probability that they coincide is $\Theta((2^np)^{-1/2})=\Theta(e^{-A n^{1/3}/2})$. In particular, when $n$ is odd, antipodal pairs have odd distance and hence antipodal collisions are absent with high probability for every fixed $A>0$. When $n$ is even, antipodes have the same parity, so the landmark parity class not appearing at the middle distance must split evenly between opposite shells and therefore must have even population. Thus having both $m_0$ and $m_1$ odd excludes every antipodal collision; by independence and the binomial parity formula, this event has probability $1/4+o(1)$.

For an even-distance collision of distance $d$, let $D$ be the set of differing coordinates and define the parity-class XOR sums $v_b=\bigoplus_{s\in S_b}s$. Because the distance histograms match, the sums of distances to the two queries coincide within each parity class, which yields
\[
 \sum_{i\in D}(-1)^{x_i}\bigl(m_b-2\#\{s\in S_b:s_i=1\}\bigr)=0.
\]
Dividing by two and reducing modulo two gives $v_b\cdot\mathbf1_D=(m_b\bmod2)(d/2+x\cdot\mathbf1_D)$ in $\mathbb F_2$. An even population $m_b$ therefore forces $\mathbf1_D$ to be orthogonal to $v_b$, while two odd populations force orthogonality to $v_0+v_1$. Neither linear restriction excludes all even-distance collisions.

\subsection{Edge failures}
We now prove the matching exponential failure-probability rates for edge resolution. For the lower bound, restrict to parallel edges whose active coordinate is the first coordinate with active pair $\{0,1\}$, indexed by their inactive word $x$ on the remaining $N=n-1$ coordinates. The distance from a landmark $(u,w)$ to this edge is $d_H(w,x)+\mathbf1_{u\notin\{0,1\}}$.

In the second-moment upper bounds on separated pairs, we use only balanced cells with $u=0$, conditioning on all landmark selections with other values of $u$. Both the factor of $q$ in shell sizes and the change from $n$ to $N$ leave the exponential rate and the $O(E_n)$ precision unchanged. For the single-pair lower bound, apply the involution that transposes the two projected queries on differing coordinates while fixing $u$; the true edge-shell law, with its bounded independent offset, produces the same concentration box as for vertices. The candidate-pair moment argument therefore gives lower exponent $\min(0,2L-a)$, together with the quantitative subexponential lower bound for $t>0$. For binary edges, the offset vanishes identically, opposite projected shells have independent $\mathrm{Bin}(2\binom Nr,p)$ counts, and the folded second-moment argument yields lower exponent $\min(0,L-a/2)$.

For the matching upper bound, we apply the all-edge projected forest estimates from Appendix~\ref{app:sparse-upper}: take first the limsup in $n$, then let $\delta\downarrow0$, and cap Markov's probability bound at one. This yields exactly the claimed exponents for edge resolution.

\subsection{Conditioning on binary parity}
Finally, we show that conditioning on the parity event that forbids binary antipodal collisions shifts the exponential failure rate of vertex and outer-vertex resolution to the nonantipodal rate.

Define $G_n=\{m_0\ne m_1\}$ when $n$ is odd, and $G_n=\{m_0,m_1\text{ both odd}\}$ when $n$ is even. As shown above, $\mathbb P(G_n)$ tends to one for odd $n$ and to $1/4$ for even $n$. For $F=F_V,F_O$, we show that
\[
 \lim n^{-1}\log\mathbb P(F\mid G_n)=\min(0,2\log 2-a).
\]
Since $G_n$ excludes all antipodal collisions, applying Markov's inequality to nonantipodal collisions and dividing by $\mathbb P(G_n)$ gives the upper bound immediately.

For the lower bound, consider an even-distance candidate pair from the two-pair estimate. Exactly $H=\binom{d_n}{d_n/2}2^{n-d_n-1}$ neutral words of each parity are equidistant from its two endpoints, since half of the differing positions must mismatch each endpoint and at least one agreeing coordinate is free to determine the parity. Neither endpoint itself is neutral. Conditional on all nonneutral landmark selections, the two neutral populations are independent $\mathrm{Bin}(H,p)$ random variables with mean $Hp\asymp e^{An^{1/3}}/\sqrt n\to\infty$. Uniformly over the nonneutral selections, the probability that the two total populations are equal is $o(1)$ by the maximum-atom bound, while the probability that both total populations are odd is $1/4+o(1)$. Because the outer collision event depends only on the nonneutral selections, it follows that $\mathbb E[Z\mathbf1_{G_n}]\ge c\mathbb EZ$ for some constant $c>0$, and restricting to $G_n$ cannot increase the second moment. The same second-moment argument therefore yields the matching lower bound and quantitative bound.

When $t=2\log2-a<0$, comparing the conditional and unconditional rates shows that, conditional on failure, the probability of even one nonantipodal collision is at most $\exp(tn/2+o(n))$. In odd dimensions, these rare failures therefore occur overwhelmingly at balanced landmark populations, even though balance itself is atypical.

\section{Polynomial corrections at the two endpoints}
Throughout the next three sections we introduce a polynomial factor and a constant prefactor into the landmark density, keeping the exponential scale fixed. Put $L=\log q$, $\sigma^2=(q-1)/q^2$, $R=\sqrt{2\sigma^2 A}$ and $a=(2/3)\sqrt{2\sigma^2}\,A^{3/2}$. The expected number of landmarks is $m_n=q^np=\kappa n^\beta e^{An^{1/3}}$, with fixed $q\ge2$, $A,\kappa>0$ and $\beta\in\mathbb R$. We include the lower endpoint $a=L$ as a benchmark for the typical-set method. Our main objective is to derive sharp first-moment estimates at the upper endpoint $a=2L$; Section~\ref{sec:classical-birthday} then establishes the corresponding probability threshold.

\subsection{Shell volume}
We begin by estimating the expected landmark count in each distance shell and summing its positive log part to obtain the shell-volume functional that governs both our typical-set and entropy bounds.

Write $\pi_r$ for the shell probability (the probability that a uniformly random word lies at distance $r$ from a fixed query) and $\lambda_r=m_n\pi_r$ for the expected number of landmarks in shell $r$. For a vertex query, the distance from a uniform word has distribution $\mathrm{Bin}(n,1-1/q)$; for an edge query, it has distribution $\mathrm{Bin}(n-1,1-1/q)+\mathrm{Bern}(1-2/q)$ with independent summands. Set $\mu=n(1-1/q)$ for the mean vertex distance. Applying Stirling's formula uniformly over all shells satisfying $|r-\mu|\le Dn^{2/3}$, for any fixed $D>R$, yields
\[
 \log\lambda_r=An^{1/3}-\frac{(r-\mu)^2}{2\sigma^2 n}
 +(\beta-\tfrac12)\log n+O(1).
\]
On this window the cubic Taylor remainder of the binomial log-probability is uniformly bounded, and the bounded shifts introduced by the edge mixture do not change the error term.

We call the shells with $|r-\mu|\le Dn^{2/3}$ the retained shells and the remaining shells the outside shells. Choose $D$ large enough that Hoeffding's inequality bounds the sum of the outside shell means by $e^{-c n^{1/3}}$. Writing $x_+=\max(x,0)$ and using the convention that shells with zero mean contribute zero to the sum, we define the shell-volume functional
\[
 V=\tfrac12\sum_r(\log\lambda_r)_+
 =an+R(\beta-\tfrac12)n^{2/3}\log n+O(n^{2/3}).
\]
To see this asymptotic, note that the half-integral of the clipped parabola with intercept $T=An^{1/3}+(\beta-1/2)\log n$ equals $(2/3)\sqrt{2\sigma^2 n}\,T^{3/2}$. Approximating this integral by its discrete Riemann sum incurs an error of $O(n^{1/3})$ by total variation, while summing the bounded shellwise errors over the retained shells costs $O(n^{2/3})$.

\subsection{The lower endpoint}
We next show that when $a=L$ and $\beta<1/2$, a typical-set counting argument forces exponentially many collisions for all three query types simultaneously with high probability.

For any fixed query, the shell counts $M_r$ are independent binomial random variables with means $\lambda_r$. On the $K=O(n^{2/3})$ retained shells, define the fluctuation scale $s_r=\max(1,\sqrt{\lambda_r})$. Each normalized deviation $|M_r-\lambda_r|/s_r$ has mean and variance at most one, so by Chebyshev's inequality their sum across all retained shells is at most $2K$ except on an event of probability at most $1/K$. At the same time, since the outside shell means sum to $e^{-c n^{1/3}}$, Markov's inequality ensures that all outside shell counts vanish with probability $1-o(1)$.

Histograms satisfying both conditions form a common typical set for all queries of a given type. Bounding the indicator of the sum constraint by an exponential weight and evaluating the resulting geometric series in each coordinate, the size of this typical set is at most
\[
 e^{2K}\prod_{r\ {\rm retained}}\sum_{z\in\mathbb Z}
 e^{-|z-\lambda_r|/s_r}
 \le e^{O(K)}\prod_{r\ {\rm retained}}s_r=e^{V+O(n^{2/3})},
\]
by geometric summation.

Applying Markov's inequality to the fraction of atypical queries shows that, with high probability, at least half of the queries of each type fall in the typical set. For outer vertices, removing the randomly selected landmarks eliminates only $o(q^n)$ vertices with high probability. When $a=L$ and $\beta<1/2$, the shell-volume formula above implies that the typical set has size $o(q^n)$. Applying the Cauchy--Schwarz inequality to the bucket occupancies within the typical set then shows that the collision counts for vertices, outer vertices and edges are simultaneously at least
\[
 \exp\{Ln+R(1/2-\beta)n^{2/3}\log n-O(n^{2/3})\}
\]
with high probability. In particular, this covers the unmodified model $\beta=0$, $\kappa=1$. The matching upper rate $L$ on the $n$ scale still follows from the earlier first-moment bound, which is unaffected by an $O(\log n)$ perturbation of $\log m_n$.

\subsection{The upper endpoint}
At the upper endpoint $a=2L$ with $\beta<1/2$, we lower-bound the expected number of collisions for each query type. The strategy is to represent the histogram discrepancy of a pair via an exchanging involution and lower-bound the collision probability using Shannon entropy.

For any two vertices, choose an isometric involution of the Hamming graph that exchanges them by applying coordinatewise transpositions. For two parallel edges sharing the same active letter-pair, choose such an involution that fixes the active coordinate and exchanges the two edges on the remaining coordinates. Decompose the vertex set into orbits of this involution. Fixed words and words equidistant from the two queries contribute zero to the histogram discrepancy. Grouping each two-point orbit into a pair of opposite contributions, the histogram discrepancy takes the form $X-X'$, where $X$ and $X'$ are independent, identically distributed sums over one chosen representative from each two-point orbit, each weighted by an independent Bernoulli coefficient multiplying a signed basis vector $e_r-e_s$. For a pair of outer vertices, we first condition on the two endpoints being absent from the landmark set, which has probability $(1-p)^2$, and then omit their two-point orbit from the sum.

Because $X$ and $X'$ are independent and identically distributed, Jensen's inequality gives $-\log\Pr(X=X')\le H(X)\le\sum_r H(X_r)$, where $H$ denotes Shannon entropy. On each retained coordinate, counting contributions from both shells of each incident orbit yields $\operatorname{Var}(X_r)\le2\lambda_r$. Adding an independent uniform random variable on $[-1/2,1/2]$ to an integer-valued variable converts its discrete Shannon entropy into differential entropy without changing its value. Because the Gaussian distribution maximizes differential entropy at a given variance, we obtain
\[
 H(X_r)\le\tfrac12\log(2\pi e(2\lambda_r+1/12))
 \le\tfrac12(\log\lambda_r)_++O(1).
\]
On each outside shell, the positive and negative counts contributing to $X_r$ each have mean at most $\lambda_r$, so $H(X_r)\le2g(\lambda_r)$, where $g(t)=(1+t)\log(1+t)-t\log t$ is the maximum Shannon entropy of a nonnegative integer-valued random variable with mean $t$. Since the outside means sum to $e^{-c n^{1/3}}$, the outside entropies sum to $o(1)$. Summing over all coordinates, the collision probability is at least $e^{-V-O(n^{2/3})}$, up to the endpoint factor $(1-p)^2$ for outer vertices.

To count pairs, restrict to vertex pairs at a single distance within two of $\mu$, chosen to be even in the binary case; Stirling's formula gives $\exp(2Ln-O(\log n))$ such pairs. For edges, fix the active direction and letter-pair and count parallel pairs in the remaining cube at its own typical distance. Multiplying the pair count by the collision probability bound shows that at $a=2L$ and $\beta<1/2$, each expected collision count is at least
\[
 \exp\{R(1/2-\beta)n^{2/3}\log n-O(n^{2/3})\}.
\]
A diverging expectation does not by itself imply failure with high probability, even though this estimate already holds for the binary even-distance nonantipodal subfamily described above.

\subsection{A sufficient range on the square-root scale}
When $\beta=1/2$, the $n^{2/3}\log n$ term in $V$ vanishes and the leading correction shifts to the $n^{2/3}$ scale, controlled by the prefactor $\kappa$. By refining the shell expansion to $o(1)$ accuracy and replacing the $L^1$ typical set with a quadratic Gaussian constraint, we obtain a sufficient range $\kappa<\kappa_0$ for collisions at $a=L$ and diverging expected collisions at $a=2L$.

Set $\beta=1/2$ and $\kappa_0=\sigma/(e\sqrt{2\pi})$. Introduce the rescaled shell coordinate $u=(r-\mu)/n^{2/3}$, the parabola $F(u)=A-u^2/(2\sigma^2)$ and the constant $c=\log\kappa-\tfrac12\log(2\pi\sigma^2)$. Uniformly for $|u|\le D$, Stirling's formula expands to
\[
 \log\lambda_r=n^{1/3}F(u)+c+\gamma_q u^3+o(1),
\]
where $\gamma_q$ is minus one sixth the third derivative of binary relative entropy in its first argument evaluated at its reference parameter $1-1/q$. On this window the quartic Taylor error is $O(n^{-1/3})$, the Stirling prefactor varies by $o(1)$, and replacing $n$ by $n-1$ or shifting a shell index by one (for the edge law) changes the exponent by $o(1)$.

Partition the retained shells into three zones: call a shell interior if $n^{1/3}F(u)$ exceeds $6\log n$, exterior if it falls below $-6\log n$, and transition otherwise. The transition zone contains $O(n^{1/3}\log n)$ shells, each with $|\log\lambda_r|=O(\log n)$. By Markov's inequality, a fixed query has with high probability zero landmarks in all exterior and far-tail shells, and at most $n^9$ landmarks in every transition shell. Consequently, the transition shells contribute $o(n^{2/3})$ to both the log cardinality of the typical set and the orbit entropy.

There are $K_0=2Rn^{2/3}+O(n^{1/3}\log n)$ interior shells, each satisfying $\lambda_r\gg n^6$. Consider the quadratic deviation statistic $Q=\frac12\sum_{\rm int}(M_r-\lambda_r)^2/\lambda_r$ over the interior shells. By independence of the shell counts and the binomial fourth central moment bound $\lambda+3\lambda^2$, the variable $Q$ has mean at most $K_0/2$ and variance $O(K_0)$. Chebyshev's inequality therefore gives $Q\le K_0/2+n^{1/2}$ with high probability. Weighting each integer histogram by a discrete Gaussian factor bounds the number of admissible interior configurations by
\[
 e^{K_0/2+n^{1/2}}\prod_{\rm int}(\sqrt{2\pi\lambda_r}+O(1)),
\]
where comparing the discrete sum to the Gaussian integral on its two monotone halves supplies the uniform additive error.

In the orbit entropy argument of the previous subsection, the variance bound improves to $\operatorname{Var}(X_r)\le\lambda_r$ because every nonzero orbit incident to row $r$ contains at most one word in shell $r$ of the first query (in fact, exactly one word in shell $r$ of one of the two queries, so each orbit contributes at most one indicator to $X_r$). Each interior coordinate entropy is therefore at most $\tfrac12\log(2\pi e\lambda_r)+o(1)$ uniformly, while the exterior and far-tail entropies still sum to $o(1)$ by the geometric bound. Both the log typical-set size and the orbit entropy are thus bounded above by
\[
 \tfrac12\sum_{\rm int}\log\lambda_r+\tfrac{K_0}{2}\log(2\pi e)+o(n^{2/3})
 =an+R\log(\kappa/\kappa_0)n^{2/3}+o(n^{2/3}).
\]
In evaluating the sum over the interior shells, the cubic term $\gamma_q u^3$ integrates to zero on $[-R,R]$ by oddness, and the clipped parabola integrates to $an$ as before.

It follows that for $\kappa<\kappa_0$, the log collision count at $a=L$ is at least $Ln+R\log(\kappa_0/\kappa)n^{2/3}+o(n^{2/3})$ with high probability, while at $a=2L$ the log expected collision count is at least $R\log(\kappa_0/\kappa)n^{2/3}+o(n^{2/3})$. This sufficient constant $\kappa_0$ does not identify a sharp probability threshold.

\section{The sharp first-moment window}
We now determine the exact $n^{2/3}$-scale asymptotics of the expected collision count at $a=2L$, locating the sharp first-moment constant $\kappa_1$.

Assume $a=2L$. Set $\kappa_1=\sigma/(2\sqrt{2\pi})$ and $B_n=R[(1/2-\beta)\log n+\log(\kappa_1/\kappa)]$. Letting $Z$ denote the collision count for any one of the three query types (vertices, outer vertices, or edges), we prove that
\[
 n^{-2/3}\log\mathbb E Z=
 \begin{cases}B_n+o(1),&q\ge3,\\
 \max(B_n,B_n/2)+o(1),&q=2.
 \end{cases}
\]
At the critical exponent $\beta=1/2$, we write $B=B_n$. When $\kappa>\kappa_1$, the right-hand side is strictly negative, so Markov's inequality and a union bound imply that a single random landmark set resolves all three query types with high probability. When $\kappa<\kappa_1$, all three expectations diverge; the matching probability converse and endpoint are proved in Section~\ref{sec:classical-birthday}. At equality $\kappa=\kappa_1$, the $n^{2/3}$ scale alone does not determine the limiting expectation.

\subsection{The common shell cost}
We first isolate the interior shell set $I$ and compute the common negative-log collision cost $J_n$ that appears in both the lower and upper bounds below.

For arbitrary fixed $\beta$, the refined shell expansion from the previous section acquires the additive term $(\beta-1/2)\log n$. Define the interior shell set $I$ to consist of all retained shells satisfying $n^{1/3}F(u)>(10+|\beta|)\log n$. On $I$ we have $\lambda_r\gg n^8$. The transition zone down to $-(10+|\beta|)\log n$ contributes $O(n^{1/3}\log^2 n)$ to the shell sums below, and the remaining shell means sum to $o(1)$. For both the vertex and edge shell laws,
\[
 J_n=\tfrac12\sum_{r\in I}\log(4\pi\lambda_r)
 =an-B_n n^{2/3}+o(n^{2/3}).
\]
Indeed, the integrated parabola yields $an$, the odd cubic term integrates to zero, and the number of interior shells is $|I|=2Rn^{2/3}+O(n^{1/3}\log n)$. Each shell contributes the constant $\tfrac12\log(4\pi\kappa/\sqrt{2\pi\sigma^2})=\tfrac12\log(\kappa/\kappa_1)$, which combines with the $(\beta-1/2)\log n$ term to produce $-B_n n^{2/3}$.

\subsection{Fourier lower bound}
To match $J_n$ from below, we express the collision probability of any pair admitting an exchanging involution as a Fourier integral, lower-bound the integrand by a Gaussian, and eliminate the cross terms by symmetry and Jensen's inequality.

For any pair of queries with an exchanging involution as constructed above, the characteristic function of the histogram discrepancy $X-X'$ factors over the two-point orbits of nonzero discrepancy, with each orbit contributing one nonnegative factor $1-2p(1-p)(1-\cos(\theta_r-\theta_s))$. Since $4p=o(1)$ and $1-\cos v\le v^2/2$, the product of these factors is bounded below by $\exp[-c_np\sum_{\rm orbits}(\theta_r-\theta_s)^2]$ for a sequence $c_n=1+O(p)>0$.

Integrating over $[-\pi,\pi]$ in each coordinate with normalized Lebesgue measure expresses $\Pr(X=X')$ as the integral of this characteristic function. Expanding the squares $(\theta_r-\theta_s)^2$ in the Gaussian lower bound produces diagonal weights $e^{-c_npN_r\theta_r^2}$, where $N_r$ is the number of nonzero orbits incident to shell $r$, together with a sum of cross terms proportional to $\theta_r\theta_s$. Under the product probability measure obtained by normalizing the diagonal weights, each coordinate $\theta_r$ is symmetric around zero, so every cross product has mean zero. Applying Jensen's inequality to the cross terms and using $pN_r\le\lambda_r$ yields
\[
 \Pr(X=X')\ge \prod_r\frac1{2\pi}\int_{-\pi}^{\pi}
 e^{-c_n\lambda_r\theta^2}\,d\theta .
\]
Evaluating the one-dimensional integrals by regime, each interior factor equals $(4\pi\lambda_r)^{-1/2}(1+o(1))$ uniformly. Each transition shell contributes at most $O(\log n)$ to the negative logarithm, and each exterior shell contributes $O(\lambda_r)$ by Jensen's inequality. Combining all shells gives the uniform lower bound $\exp[-J_n+o(n^{2/3})]$, retaining the harmless endpoint factor $(1-p)^2$ for outer queries.

\subsection{Descending bundles for bulk pairs}
We now prove the matching upper bound $\exp[-J_n+o_\delta(n^{2/3})]$ on the collision probability for all bulk pairs of queries. The strategy is to expose selected landmarks in descending order of distance deviation from the mean $\mu$: for bulk separations, conditional contraction ensures that almost every word in shell $r$ of one query lies strictly closer to $\mu$ from the other query, allowing us to extract a fresh local-limit bound at each layer.

\step{Conditional contraction for bulk pairs.}
To measure the separation $d$ between two query centers, delete the one or two active coordinates in the case of edges, leaving $N=n-O(1)$ free coordinates; for vertices, take $N=n$. Fix a small $\delta>0$. We call a pair bulk if $d\ge\delta N$ when $q\ge3$, and if $\delta N\le d\le(1-\delta)N$ when $q=2$.

Condition on the radius $r$ from the first query and on the letters in any active coordinates. The mismatches against the first query then form a uniformly random subset of size $r-b$, with $0\le b\le2$, inside the $N$ free coordinates. Under this conditioning, the radius from the second query has mean $\mu+\rho(r-\mu)+O_q(1)$, where the contraction factor $\rho=1-qd/((q-1)N)$ satisfies $|\rho|\le1-\eta(q,\delta)$. Its conditional variance is $O_q(n)$, combining the hypergeometric variance for the number of mismatches falling on coordinates where the two queries differ with the binomial variance for the letters chosen at those mismatches. By symmetry, the same mean and variance estimates hold with the roles of the two queries reversed.

\step{Descending exposure of disjoint bundles.}
For each interior shell $r\in I$, define two bundles of words: those at radius $r$ from the first query whose radius from the second query has strictly smaller absolute deviation from $\mu$, and symmetrically those at radius $r$ from the second query whose radius from the first has strictly smaller absolute deviation from $\mu$. Assign opposite signs to the two bundles according to which query attains $r$. Because a word can belong to a bundle only at the strictly larger of its two absolute deviations from $\mu$, all bundles across all shells $r\in I$ are pairwise disjoint.

Condition on the landmark indicators of all words outside these bundles, and expose both bundles completely in descending order of the absolute deviation $|r-\mu|$. At shell $r$, the collision constraint for row $r$ forces a single value for the difference between the two bundle counts. Later bundles cannot affect row $r$, nor can other rows at the same absolute deviation, because every word in a later bundle has strictly smaller deviation at both queries. Moreover, because we expose the entire landmark selections in each bundle rather than only their counts, all contributions from earlier bundles to later rows are completely known when those rows are reached. Sequential conditioning therefore bounds the collision probability by the product of the maximum point probabilities of the signed bundle-count differences.

\step{Bounding the layer atoms.}
For shells far from the center, namely $|r-\mu|\ge\sqrt n\log n$, conditional Chebyshev's inequality implies that each of the two bundles at shell $r$ has size $(1-O_\delta(1/\log^2 n))q^n\pi_r$. Given two disjoint bundles of sizes $s,t$ of this magnitude, Fourier inversion bounds the maximum point probability of their signed binomial difference by the normalized integral over $[-\pi,\pi]$ of
\[
 [1-2p(1-p)(1-\cos\theta)]^{(s+t)/2}.
\]
Rescaling the integration variable by $\sqrt{(s+t)p(1-p)}$, applying Taylor expansion near zero, and using the global envelope $1-\cos\theta\ge2\theta^2/\pi^2$ gives the uniform upper bound $(4\pi\lambda_r)^{-1/2}(1+o(1))$.

For intermediate shells satisfying $C_\delta<|r-\mu|<\sqrt n\log n$, we restrict each bundle to a single balanced subcell: round the within-block mismatch allocations to integers while preserving their total, and distribute the mismatch letters as evenly as possible across the alphabet. By the balanced-cell estimate proved earlier, this subcell captures a conditional fraction $n^{-O_q(1)}$ of the shell, and every word in it has both radii within $O_q(1)$ of their conditional means. For edges, we also fix a compatible pattern on the active coordinates that carries at least a polynomial conditional mass. Choosing the constant $C_\delta$ large enough guarantees that the other radius contracts strictly inside this subcell. Conditioning on all words outside the subcell, the maximum binomial atom of the subcell loses at most $O(\log n)$ in the negative logarithm relative to $\frac12\log(4\pi\lambda_r)$. Across all intermediate shells, this loss accumulates to $O(\sqrt n\log^2 n)=o(n^{2/3})$.

Finally, omitting the remaining bounded number of rows closest to $\mu$ costs $O(n^{1/3})$ in the negative logarithm. Multiplying over all exposed layers shows that every bulk collision, including a pair of nonparallel edges, has probability at most $\exp[-J_n+o_\delta(n^{2/3})]$.

\subsection{Boundary separations and binary reflections}
To complete the sharp first-moment formula, we combine the bulk estimate with exponential upper bounds on nonbulk separations and, in the binary case $q=2$, evaluate the antipodal branch that produces the $B_n/2$ term.

The total number of query pairs is $\exp(2Ln+O(\log n))$, and this many pairs already occur at a single fixed typical separation (taking parallel edges with a fixed active letter-pair for the edge lower bound). Combining this count with the matching lower and upper bounds on bulk collision probabilities gives a bulk contribution of $\exp[B_n n^{2/3}+o(n^{2/3})]$.

Write $h$ for the binary entropy function and set $b_q(\delta)=h(\delta)+\delta\log(q-1)$. For $q\ge3$, the number of nonbulk pairs is at most $\exp[(L+b_q(\delta))n+O(\log n)]$, while the shifted-block forest proved earlier bounds each nonbulk collision probability by $\exp[-\sqrt{1-\delta}\,an+o(n)]$. Choosing $\delta$ sufficiently small makes the combined exponent strictly negative, so nonbulk pairs are negligible.

In the binary case $q=2$, nonantipodal vertex pairs satisfy the earlier uniform rate $-a$, while the number of nonbulk pairs is at most $\exp[(L+h(\delta))n+O(\log n)]$. Parallel nonantipodal edges satisfy the same bound after fixing the active bit. For nonparallel binary edges, delete both active coordinates so that $N=n-2$ free coordinates remain, and treat the three separation regimes by choosing offsets on the two active coordinates:
when $0<d<N$, set both offsets to zero and apply the nonantipodal forest;
when $d=0$, set the two offsets to zero and one to obtain a fixed-shift path;
when $d=N$, take the midpoint $t=(N+1)/2$, set both offsets to zero below $t-1$ so the other radius is $N-r$, and set both offsets to one above $t+1$ so the other radius is $N+2-r$.
In either half of the $d=N$ case, the other radius lies strictly closer to $t$. Because the offsets use independent active bits and the chosen classes have free radius $r$ or $r-1$, they retain at least a polynomial fraction of each shell in the retained range. This descending forest again achieves rate $-a$, and the polynomial number of active-coordinate choices does not close the exponential gap.

It remains to evaluate binary antipodal pairs, including parallel projected antipodal edges. There are $\exp(Ln+O(\log n))$ such pairs. Because distances to antipodal queries are complementary, the collision event factors across lower-half radii into independent coincidences of equal-size binomials of mean $\lambda_r$; in the edge case, each projected shell has $2\binom{n-1}{r}$ words. For equal bundle sizes, the Fourier integral from the previous subsection is exact at zero and has a nonnegative integrand. Each interior coincidence factor is therefore $(4\pi\lambda_r)^{-1/2}(1+o(1))$. Transition shells contribute $O(\log n)$ to the negative logarithm by concentration and Cauchy--Schwarz, and exterior shells contribute $O(\lambda_r)$ by considering the event of selecting no landmarks. By reflection symmetry, the log collision probability is $-J_n/2+o(n^{2/3})$, with a bounded shift of the center or the omission of a midpoint row having negligible effect. For outer vertices, deleting the endpoint orbit costs only $(1-p)^2$. Multiplying by the pair count yields a branch contribution of $\exp[(B_n/2)n^{2/3}+o(n^{2/3})]$, proving the asserted maximum for $q=2$.

\section{Failure probability in the first-moment window}
We now convert the sharp first-moment asymptotics at $\beta=1/2$ into matching logarithmic asymptotics for the nonresolution probability, using the second moment method on restricted collision families.

Fix $a=2L$ and $\beta=1/2$, so that $J_n=2Ln-Bn^{2/3}+o(n^{2/3})$. Let $F$ denote the nonresolution event for any one of the three query types, or for the union of all three. We prove that
\[
 \lim_{n\to\infty} n^{-2/3}\log\Pr(F)=
 \begin{cases}\min(0,B),&q\ge3,\\ \min(0,B/2),&q=2.\end{cases}
\]
A zero limit on this scale is not a claim of high-probability nonresolution.

\subsection{Two typical pairs}
To prove the lower rate $\min(0,B)$, we restrict the collision count $Z$ to pairs at a single typical separation and upper-bound $\mathbb E[Z^2]$ by splitting ordered pairs of query pairs into three cross-distance regimes.

Count only query pairs at one fixed typical separation $d_N=(1-1/q)N+O(1)$ (chosen even in the binary case), requiring absent endpoints for outer vertices and restricting to parallel edges of one fixed direction and letter-pair for edges. By the sharp first-moment bounds proved above, the resulting count $Z$ satisfies $\mathbb EZ=\exp[Bn^{2/3}+o(n^{2/3})]$. For edges, distance measurements still run over the full word space, and we condition on the boundedly many active letters whenever applying contraction.

For two such pairs, compare their four query centers by Hamming distance, taken modulo complements when $q=2$. Choose $\delta>0$ small enough that each within-pair distance exceeds $2\delta n$. By the triangle inequality, cross edges (between one pair and the other) of length at most $\delta n$ cannot share an endpoint and therefore form a matching. Call a cross edge micro-close if its distance is at most $\lfloor\sqrt n\rfloor$.

\step{No cross edge within $\delta n$.}
In this regime every pair among the four queries is well separated, so all pairwise comparisons contract uniformly. For each row of each discrepancy block, form two signed bundles consisting of words for which one query of that pair attains the unique largest absolute radius deviation among all four queries; words with ties stay outside the bundles. Conditional Chebyshev's inequality and a union bound over the three other queries retain a $1-o(1)$ fraction of each oriented shell outside the central $O(\sqrt n\log n)$ rows. On those central rows, use balanced cells rounded simultaneously for all four queries, omitting a bounded number of rows closest to $\mu$. Complete descending exposure works unchanged because no bundled word touches another row at its own layer. The joint collision probability is therefore at most $\exp[-2J_n+o(n^{2/3})]$. Summing over at most $\exp(4Ln+O(\log n))$ pairs of pairs contributes at most $\exp[2Bn^{2/3}+o(n^{2/3})]$.

\step{A cross edge within $\delta n$, at most one micro-close.}
Choosing the first pair freely, one endpoint of the second pair inside a ball of radius $\delta n$ around an endpoint of the first pair, and the remaining endpoint freely gives at most $\exp[(3L+b_q(\delta))n+O(\log n)]$ configurations; in the binary case, including complementary balls only doubles this bound. Because the short cross edges form a matching and at most one is micro-close, each pair contains an anchor query not incident to any micro-close edge. Its distance from all three other endpoints in the chosen metric exceeds $\lfloor\sqrt n\rfloor$, so each conditional contraction slack is at least $c_q/\sqrt n$. Inside the four-query balanced subcell, all other radii lie within $O_q(1)$ of their conditional means and therefore contract strictly outside the central $O_q(\sqrt n)$ shells. Exposing these rigid anchor columns in both discrepancy blocks in descending order, the omitted central shells cost $O(n^{5/6})=o(n)$ and the polynomial cell losses cost $O(n^{2/3}\log n)=o(n)$. The joint collision probability is thus at most $\exp[-2an+o(n)]$ uniformly. Since $b_q(\delta)<L$ for small $\delta$, the total contribution of this regime is exponentially small.

\step{Two micro-close cross edges.}
Because the two micro-close cross edges form a matching, each endpoint of the second pair lies in a micro-close ball around a distinct endpoint of the first pair. A micro-close ball, together with its complementary copy in the binary case, has size $\exp[O_q(\sqrt n\log n)]$, so only $\exp[2Ln+o(n^{2/3})]$ configurations occur. Bounding the joint collision event by a single pair's collision at cost $\exp[-J_n+o(n^{2/3})]$ yields a total contribution of at most $\exp[Bn^{2/3}+o(n^{2/3})]$.

In all three upper bounds on $\mathbb E[Z^2]$, we may drop the endpoint-absence constraint for outer vertices. Combining the three regimes with the Cauchy--Schwarz second-moment inequality gives
\[
 \Pr(Z>0)\ge (\mathbb E Z)^2/\mathbb E[Z^2]
 \ge\exp[(\min(0,B)+o(1))n^{2/3}].
\]
For $q\ge3$, Markov's inequality applied to the full collision count supplies the matching upper bound.

\subsection{Binary antipodal families}
In the binary case $q=2$, the antipodal branch dominates whenever $B<0$. We repeat the second-moment analysis on antipodal pairs to obtain the lower rate $\min(0,B/2)$ and complete the proof.

Restrict attention to antipodal vertex pairs, or to parallel projected antipodal edge pairs of one fixed active direction. There are $\exp(Ln+O(\log n))$ such pairs, and by the previous section their expected collision count is $\exp[(B/2)n^{2/3}+o(n^{2/3})]$. For two such pairs, measure the distance between their representatives modulo complements, and divide into the same three regimes.

When the representative distance exceeds $\delta n$, form two folded discrepancy blocks in which each row below the midpoint records the difference of opposing ordinary shell counts. Within each opposing shell, keep those words whose distance to the other pair's representative has strictly smaller absolute deviation from the midpoint. Conditional Chebyshev's inequality retains a $1-o(1)$ fraction of each shell outside the central range, while inside the central range we use rigid balanced cells as before. Descending exposure yields a negative-log cost of $J_n/2+o(n^{2/3})$ per block. Counting pairs of antipodal pairs then bounds the joint contribution by $\exp[Bn^{2/3}+o(n^{2/3})]$.

When the representative distance lies between $\lfloor\sqrt n\rfloor$ and $\delta n$, the folded anchor rows still contract with slack $c/\sqrt n$; omitting $O(\sqrt n)$ central rows gives a joint collision probability of at most $\exp[-an+o(n)]$. Since there are at most $\exp[(L+h(\delta))n+O(\log n)]$ configurations in this regime, an exponential gap remains. At micro-close distance, bounding the joint event by one collision of probability $\exp[-J_n/2+o(n^{2/3})]$ across at most $\exp[Ln+o(n^{2/3})]$ configurations contributes at most $\exp[(B/2)n^{2/3}+o(n^{2/3})]$.

Applying the same second-moment inequality yields the lower rate $\min(0,B/2)$, while the full first-moment bound and the trivial upper bound of one supply the matching upper rate. For outer vertices, the endpoint-absence factor does not affect this scale. For parallel edges, the folded shell center is $(n-1)/2$ and the projected shell sizes are doubled, as already accounted for in their full shell law. Finally, taking maxima and applying a finite union bound transfers the result to the union of all three failure events.

\section{The classical high-probability birthday threshold}
\label{sec:classical-birthday}

In this section we prove that at \(A=C_q\) the classical high-probability threshold has the sharp prefactor \(\kappa_1=\sigma/(2\sqrt{2\pi})\) at scale \(\sqrt n\exp(C_q n^{1/3})\), with equality belonging to the failure regime. We treat three target types: ordinary vertices, outer vertices, and ordinary edges.

The strategy for ordinary vertices is a second-moment argument on a large family of vertex pairs at nearly typical distance. We first show that for typical pairs of candidate pairs, the joint collision probability factorizes asymptotically via Fourier localization and a Gaussian covariance comparison. We then bound the contribution of atypical pairs of pairs using a descending-radius argument, verify that the expected number of collisions diverges even at the critical prefactor, and finally adapt the proof to outer vertices and ordinary edges.

\step{Candidate pairs and their joint probability.}
We begin by defining our family of candidate pairs and the subset of typical configurations, with the goal of proving that any two typical pairs collide almost independently. Write \(s=n^{1/3}\) and \(K=pq^n\), and consider unordered vertex pairs at distance \(h\) nearest to the mean distance \(\mu=(1-1/q)n\), choosing \(h\) even when \(q=2\). The number of such candidate pairs is \(M=\exp(2Ln-O(\log n))\), and by symmetry they share a common ordinary collision probability \(P\).

We say that two candidate pairs \(x,y\) and \(u,v\) form a typical configuration if every ordered quadruple of letters from the alphabet appears among the coordinate tuples \((x_i,y_i,u_i,v_i)\) with frequency \(n/q^4+O(n^{11/12})\), for a fixed tolerance multiplier. Conditioning on the two within-pair distances being equal to \(h\) costs only a polynomial factor in probability, so standard concentration implies that the fraction of atypical configurations is \(\exp(-\Omega(n^{5/6}))\). Whenever a configuration is typical, all six pairwise distances among the four endpoints lie at \(\mu+O_q(n^{11/12})\).

\step{Disjoint singleton columns.}
To localize the characteristic function of the joint collision event, we first remove the linear dependencies among the shell discrepancy rows and construct disjoint families of words that perturb a single row at a time. In each of the two blocks (corresponding to the pairs \(x,y\) and \(u,v\)), the discrepancy rows sum to zero, so we delete one central discrepancy row per block. In the binary case with \(h\) even, the row sums vanish separately within each parity class, so we delete one central row of each parity instead. In the Fourier inversion integral, we set the Fourier frequency corresponding to each deleted row to zero.

We claim that within a central window of width \(O(n^{2/3})\) around \(\mu\), every remaining radius \(r\) admits disjoint sets of words that produce singleton columns in the two blocks, each with Poisson intensity at least \(\lambda_r\exp(-O(n^{1/6}+\log n))\). To construct the family for the first block, we prescribe words \(z\) satisfying
\[
 d(x,z)=r,\quad d(y,z)=c,\quad d(u,z)=d(v,z),
\]
where \(c\) is the deleted central row of the first block (of matching parity in the binary case).

To lower-bound the number of such words, consider the product law on words in which each coordinate independently disagrees with \(x\) with probability \(r/n\), choosing uniformly among the nonmatching letters. Conditional on the total distance from \(x\) being \(r\), this tilted product law is uniform on the Hamming sphere of radius \(r\) around \(x\). Partition the \(n\) coordinates into finitely many cells according to the coordinate tuple of the four endpoints, and round the expected symbol counts in each cell to integers while preserving the cell sizes. Under this tilted law, the conditional mean distance from \(z\) to any center at distance \(d\) from \(x\) is \(\mu+(1-qd/((q-1)n))(r-\mu)\). Because all six pairwise distances are \(\mu+O_q(n^{11/12})\) and \(r-\mu\) is at most \(O(n^{2/3})\), the rounded allocations satisfy the three constraints up to residuals of size \(O(1),O(1),O(n^{7/12})\), respectively.

We now adjust the symbol counts to eliminate these residuals. When \(q\ge3\), coordinates with endpoint tuples \((0,0,0,0),(0,1,0,0),(0,0,1,0)\) allow single-coordinate moves that change the triple of quantities \((d_x,d_y,d_u-d_v)\) by the vectors \((1,1,0),(0,1,0),(0,0,1)\), respectively, by changing the landmark letters \(0\to1,1\to2,1\to2\) in those cells. In the binary case, we translate coordinates so that \(x\) is the all-zero word and express the three constraints as linear equations on the counts of ones, with coefficient columns \((1,y_i,v_i-u_i)\) and right-hand sides \((r,(r+h-c)/2,(|v|-|u|)/2)\), all of which are integers. The cells with \((y_i,u_i,v_i)=(0,0,0),(1,0,0),(0,0,1)\) form a unimodular minor.

In both cases, we can cancel the residuals by making a bounded number of prescribed adjustments of size \(O(n^{7/12})\) inside cells of linear size, so every symbol count stays strictly interior to its cell. By Stirling's formula, the multinomial probability of the adjusted allocation under the tilted product law is at least \(\exp(-O(\log n+(n^{7/12})^2/n))\), and conditioning on the total distance \(r\) only increases this probability. Multiplying by the shell intensity \(\lambda_r\) gives the claimed Poisson intensity lower bound. Exchanging the roles of the two blocks produces the corresponding family for the second block, and by construction the singleton events for distinct radii and blocks are mutually disjoint.

\step{Joint intensity bound.}
Next we bound the Poisson intensity \(J_{r,t}\) of words lying simultaneously at radius \(r\) from one center and radius \(t\) from another distinct center in a typical configuration. For any two distinct typical centers and any \(r,t\) in the central window, we have the uniform bound
\[
 J_{r,t}\le \lambda_r\lambda_t K^{-1}\exp(o(s)).
\]
To prove this, apply the exponential Markov inequality to the exact probability of the atom \((r,t)\) using the signed tilts \((r-\mu)/(n\sigma^2),(t-\mu)/(n\sigma^2)\). Because the two centers are at distance \(\mu+O_q(n^{11/12})\) and both radii lie within \(O(n^{2/3})\) of \(\mu\), the mixed quadratic term in the cumulant generating function has size \(O(n^{1/4})=o(s)\), while the cubic remainder is \(O(1)\). By Stirling's approximation, the marginal shell exponents match the corresponding quadratic exponents up to \(O(\log n)\), which yields the stated upper bound.

\step{Fourier localization and Gaussian reduction.}
We now pass from the Bernoulli model to the Poisson approximation and localize the joint characteristic function to a small box on the dense rows, where the integrand simplifies to a product of small-row factors and a multivariate Gaussian. Replacing the Bernoulli log characteristic function by its Poisson counterpart introduces a uniform error of \(O(q^np^2)\), and discarding the tail frequencies outside the central window adds a further uniform error of \(o(1)\) to the log characteristic function.

Call a row dense if its shell intensity satisfies \(\lambda_r\ge e^{n^{1/4}}\), and set the box parameter \(H=e^{n^{1/5}}\). Whenever a dense coordinate lies outside the box \(|\theta_r|\le H/\sqrt{\lambda_r}\), the presence of a singleton column for that row bounds the modulus of the characteristic function by \(\exp(-\exp(\Omega(n^{1/5})))\). The same tail bound applies to the single-block Fourier integral, so in both integrals the discarded regions are negligible compared to the lower bound \(P^2\ge\exp(-2an-O(n^{2/3}\log n))\).

Inside the box, group the phase variables into three sets: the small rows of the first block, the small rows of the second block, and the combined dense rows of both blocks. When we split the Poisson exponent into three separate sums corresponding to these groups, the error is bounded by sums of products of individual phases drawn from different groups. Every such product involves distinct endpoints, so the joint intensity bound proved above controls its contribution. Because the \(\lambda_r\)-weighted sums of the absolute small phases and dense phases are at most \(n\exp(o(s))\) and \(O(nH\sqrt K)\), respectively, the total splitting error is \(o(1)\). Expanding the combined dense exponent in Taylor series, the linear term cancels by symmetry and the cubic remainder is \(O(nH^3e^{-n^{1/4}/2})\).

Consequently, on the box the joint integrand is uniformly \(1+o(1)\) times the product of the two single-block small-row characteristic functions and the dense Gaussian \(\exp(-\theta^TQ\theta/2)\), where \(Q\) denotes the dense Poisson covariance matrix. Each small-row factor is real and positive by the involution that swaps the two endpoints of its block. Having a positive proxy justifies converting uniform pointwise relative errors into relative errors for the integral, even though the joint integrand itself need not be real. Moreover, the disjoint singleton columns imply the lower matrix bound \(Q\succeq\exp(-O(n^{1/6}+\log n))\operatorname{diag}(\lambda_r)\), which allows us to extend the Gaussian integral from the box to all of Euclidean space via a union bound on Gaussian tails. By the same analysis, each single-block integral decomposes into the same small-row factor and the marginal Gaussian integral with covariance matrix \(Q_{ii}\).

\step{Radial basis and cross-Gram estimate.}
To show that the joint Gaussian integral factorizes into the product of the two marginal Gaussian integrals, we analyze the covariance matrices in a radial orthonormal basis and bound the cross-block correlation. In the uniform \(L^2\) space on a single coordinate, normalize the centered indicator of a given letter to have unit norm; the inner product between the normalized indicators of two different letters is \(-1/(q-1)\). For any center and each degree \(k\), define the unit degree-\(k\) radial vector as the normalized sum of coordinatewise products of these indicators over all \(k\)-subsets of coordinates. Vectors of distinct degrees are orthogonal, so these \(n+1\) vectors form an orthonormal radial basis around the center.

For two centers at distance \(d\), the inner product \(f_k(d)\) between their degree-\(k\) radial vectors equals the average over all \(k\)-subsets of the product of \(k\) entries drawn without replacement from a multiset containing \(d\) copies of \(-1/(q-1)\) and \(n-d\) copies of \(1\). For any fixed \(k\ge1\), comparing sampling without replacement to sampling with replacement incurs an error of \(O(k^2/n)\), so at typical distances \(f_k(d)\) tends uniformly to zero.

To control larger degrees, we separate the cases \(q\ge3\) and \(q=2\). When \(q\ge3\), applying the elementary symmetric mean inequality to the absolute values of the entries gives \(|f_k(d)|\le [1-(1-1/(q-1))d/n]^k\). When \(q=2\), the symmetry \(f_{n-k}(d)=(-1)^df_k(d)\) allows us to restrict to degrees up to \(n/2\). For \(k\le C\log n\), we use the sampling-with-replacement estimate. For \(C\log n\le k\le n/2\), we pair opposite signs into at least \(n/3\) disjoint pairs; only subsets that meet each such pair in zero or two coordinates can avoid cancellation in the average. Under independent Bernoulli sampling with parameter \(k/n\), the probability of avoiding all single-element intersections is at most \(\exp(-k/3)\), and conditioning on the subset having exact size \(k\) costs a factor of \(O(\sqrt k)\). Combining these bounds, we conclude that in both cases the sum of squared inner products is \(o(1)\), omitting degree 0 and, in the binary case, degree \(n\) (which vanishes identically on signed differences of even-distance pairs).

Now return to the covariance matrices. In each block, the signed shell difference vectors lie in the linear span of the differences between the radial vectors of the two endpoints. For every nonzero degree, the squared norm of this difference vector is \(2(1-f_k(h))\), which is bounded away from zero. Orthonormalizing the full block spans, the cross-Gram matrix between the two blocks is diagonal in the degree index and has Frobenius norm \(o(1)\), by the four cross-distance bounds on \(f_k\) just established. Restricting to the actual dense row subspaces corresponds to compressing this cross-Gram matrix by isometries, which cannot increase its Frobenius norm. Hence the whitened cross-covariance matrix \(T=Q_{11}^{-1/2}Q_{12}Q_{22}^{-1/2}\) has Frobenius norm \(o(1)\), and the ratio of the joint Gaussian integral to the product of the two marginal Gaussian integrals is \(\det(I-T^*T)^{-1/2}=1+o(1)\). Since the small-row factors cancel between the numerator and denominator, we obtain \(\Pr(\text{both collide})=(1+o(1))P^2\) uniformly over all typical configurations.

\step{Atypical pairs and exact cutoff.}
We now control the contribution of atypical configurations to the second moment, showing that they are negligible compared to \((MP)^2\), and then verify that \(MP\to\infty\) even at the exact critical prefactor \(\kappa=\kappa_1\).

First consider atypical pairs of candidate pairs. Whenever there is at most one cross-pair match within distance \(\delta n\), the balanced-cell two-pair bound from earlier already suffices, because the \(\exp(-\Omega(n^{5/6}))\) fraction of atypical configurations overwhelms the \(\exp(O(n^{2/3}\log n))\) loss in that bound.

Next suppose there are two cross-pair matches within distance \(\delta n\), and at least one of them has nonzero quotient distance (for \(q\ge3\), this simply means nonzero ordinary distance). Choose the two endpoints of that match as the anchors, one in each block. By the triangle inequality, the other two endpoint comparisons for each anchor have quotient distance bounded away from zero. We modify the rounded balanced allocations by \(O_q(1)\) coordinates so as to force the close comparison strictly inward toward \(\mu\). Specifically, on the positions where the anchor and its close comparator differ, the difference between their distances to a landmark is the count of anchor letters minus the count of comparator letters, which under the tilted shell law has mean \(-qd(r-\mu)/((q-1)n)\). When the separation \(d\) is sufficiently large, both letter counts are \(\Omega_q(d)\), so swapping a bounded number of letters from one to the other gives the difference the strict inward sign, and a compensating bounded adjustment on agreeing positions restores the total radius \(r\). If instead \(d=O_q(1)\), we place only comparator letters on the differing positions when \(r>\mu\), and only anchor letters when \(r<\mu\). Because all modifications are bounded, omitting a fixed central strip wide enough that these adjustments cannot cross \(\mu\) or disrupt the two far comparisons ensures that every comparison moves inward. In the binary case, we complement the close comparator first if needed, which preserves absolute deviation from the mean.

With these modifications, the balanced tilted cell probabilities remain polynomially large, and the same two-block descending-radius argument gives a joint collision probability of \(\exp(-2an+O(n^{2/3}\log n))\). The number of pairs of pairs with two close cross-matches is at most \(\exp((2L+2b_q(\delta))n+O(\log n))\), which becomes exponentially small after dividing by \(M^2P^2\) once \(\delta\) is chosen sufficiently small. Finally, if both cross-matches have zero quotient distance, there are only \(O(M)\) such pairs of pairs, and each has joint probability bounded by \(P\). Putting all cases together, whenever \(MP\to\infty\) the second moment of the collision count is \((1+o(1))(MP)^2\), which yields concentration and failure with high probability.

\step{Divergence of the expected count at equality.}
Our earlier lower bounds already imply \(MP\to\infty\) whenever \(a<2L\), as well as at \(a=2L\) when \(\beta<1/2\) or when \(\beta=1/2\) and \(\kappa<\kappa_1\). It remains to prove that \(MP\to\infty\) also holds at equality, \(\kappa=\kappa_1\). We do this by computing the leading terms of the Fourier Jensen lower bound for \(\log P\) through order \(s=n^{1/3}\) and showing that the savings from deleting one central row exceed the total excess cost across all remaining rows.

Set \(v=\sigma^2\), \(p_0=1-1/q\), \(R=\sqrt{2vA}\), and rescale the radius by \(x=(r-\mu)/n^{2/3}\). On the retained window at equality, Stirling's expansion gives the log shell means
\[
 h_r=\log(4\pi\lambda_r)=s(A-x^2/(2v))-c_3x^3
       +s^{-1}(b x-c_4x^4)+O(s^{-2}),
 \quad c_3=\frac{2p_0-1}{6v^2},\quad
 c_4=\frac{1-3v}{12v^3},
\]
where \(b\) is the linear Stirling correction. We claim that the half-sum of the positive parts \((h_r)_+\) satisfies
\[
 an+sR^5\left(\frac{c_3^2v}{2}-\frac{c_4}{5}\right)+o(s)
 =an-sR^5\frac{1+2v}{360v^3}+o(s).
\]
To obtain this expansion, replace the discrete sum by the integral of its clipped smooth approximation; the trapezoidal rule error, including the two corner intervals where the argument crosses zero, is \(o(s)\). On the symmetric interval \([-R,R]\), the odd terms integrate to zero. Near each endpoint \(x=\pm R\), the shift of the zero-crossing caused by the \(-c_3x^3\) term contributes a triangular region of area \((c_3R^3)^2/(2s(R/v))\) before multiplying by the half-density of shells, while the \(O(s^{-2})\) shell error sums to \(O(1)\) across the window.

Now apply the single-block Fourier Jensen bound with the central frequency already fixed to zero. Each remaining factor in the product is the average of \(e^{-c_n\lambda_r\theta^2}\) over \([-\pi,\pi]\), where \(c_n=1+O(p)\) has negligible effect. Compare the negative logarithm of each factor with \((h_r)_+/2\). When \(h_r<0\), Jensen's inequality bounds the excess by \(\pi e^{h_r}/12\); when \(h_r\ge0\), the excess equals \(-\log\operatorname{erf}(\sqrt{\pi e^{h_r}/4})\). Near each boundary \(h_r=0\), consecutive values of \(h_r\) change by \((R/v+o(1))/s\). Summing these two excess bounds over all rows in the window therefore costs at most \((2v/R)sJ+o(s)\), where
\[
 J\le\frac{\pi}{12}+
     \frac{e^{-\pi/4}}{(\pi/4)(1-e^{-\pi/4})}<1.34 .
\]
For the positive part \(h_r\ge0\), this integral bound follows from \(\operatorname{erfc}(z)\le e^{-z^2}\) and the elementary inequality \(e^h\ge 1+h\). In the rescaled Riemann sum, both bounds are integrable and dominate the tails because of the linear slopes at the two boundaries. Outside the retained window, Jensen's inequality and the shell tails bound the total cost by \(o(1)\).

On the other hand, deleting one central row already saves \(As/2+o(s)\) in the negative log probability. This savings strictly exceeds the total excess bound \((2v/R)sJ+o(s)\), because \(AR/(4v)=3L/(4v)\ge3\log2>1.34\). It follows that \(\log(MP)\ge c s\) for some constant \(c>0\), so the second-moment argument applies at equality as well.

\step{Outer vertices and ordinary edges.}
We finish the section by extending the threshold proof to outer vertices and ordinary edges.

For outer vertices, any ordinary collision between two vertices forces their endpoint selectors to agree, because the radius-zero shell must match. The contributions of the two endpoints to the discrepancy cancel both when both endpoints are selected and when neither is selected. By independence of the endpoint selectors from the remaining landmarks, the ratio of outer to ordinary collision probabilities is \(P_{\rm out}/P=(1-p)^2/((1-p)^2+p^2)=1-o(1)\). Since the ordinary joint collision bounds dominate the outer joint collision bounds, the second-moment concentration carries over directly to outer vertices.

For ordinary edges, we restrict to the parallel family of edges whose differing coordinate is the first coordinate, with letters \(\{0,1\}\). For any landmark, the distances to the two endpoints of an edge separate according to the common offset \(b\in\{0,1\}\) contributed by the first coordinate. Offset zero occurs with probability \(2/q\) and supplies both the singleton column constructions and the balanced-cell bounds. In the Gaussian comparison, the cross-Gram bound holds in the direct sum of the two orthogonal offset fibers, which increases the squared Frobenius norm bound by at most a factor of two and remains valid under compression to the actual dense rows. The joint intensity bounds follow by conditioning on the offset and shifting each projected shell radius by at most one. In the binary case, the offset is always zero and both parity gauges continue to apply.

Compared with the \(n\)-coordinate vertex shell probability, the exact edge shell probability differs by the ratio \(1-(r/n-p_0)/p_0\). Taking its logarithm adds only an odd linear term of order \(s^{-1}\) to the expansion of \(h_r\) at equality, which integrates to zero on \([-R,R]\) and does not affect the leading calculation. Passing to the projected dimension changes \(M\) by at most a constant factor up to its polynomial prefactor. This establishes the failure converse for all three target types, and hence simultaneously by a union bound. The matching high-probability success assertions follow immediately from the first moment.

\section{The narrow vertex window for alphabets of size at least three}
\label{sec:classical-poisson}

For fixed alphabet size \(q\ge3\), the sparse Bernoulli transition for ordinary and outer vertices exhibits a finer critical window around the threshold of Section~\ref{sec:classical-birthday}. This section establishes a Poisson limit for the number of collisions in that window; it does not treat the binary antipodal case, ordinary edges, or fixed-size sampling.

Recall the notation \(s=n^{1/3}\), \(v=\sigma^2\), \(A=C_q\), and \(R=\sqrt{2vA}\), and define the auxiliary function and constants
\[
 \begin{split}
 G(h)&=-\log\!\left[e^{-e^h/(2\pi)}
                   I_0(e^h/(2\pi))\right]-\tfrac12(h)_+,\\
 J_j&=\int_{\mathbb R} h^j G(h)\,dh\quad(j=0,1),\\
 t_*&=\frac1R\left(\frac A2+
    \frac{R^5(1+2v)}{360v^3}-\frac{2v}{R}J_0\right),
 \qquad u_*=\frac{1}{4R},
 \qquad D=\frac{R^3}{6v}+\frac{2v^2}{R^3}J_1 .
 \end{split}
\]
Here \(I_0\) is the modified Bessel function of the first kind; see its \href{https://dlmf.nist.gov/10.32.E1}{integral representation} and \href{https://dlmf.nist.gov/10.40.E1}{large-argument expansion}. From these properties, \(G\) and its piecewise derivatives decay exponentially at both ends of the real line, so the integrals \(J_0\) and \(J_1\) converge absolutely.

We prove that for any fixed \(y\in\mathbb R\) and expected landmark count
\[
 K_n=\kappa_1\sqrt n\,\exp\!\left(
 As+\frac{t_*}{s}+\frac{u_*\log n+y}{s^2}\right),
\]
the number of unordered ordinary vertex collisions converges in distribution to a Poisson random variable with mean
\[
 \Lambda_q(y)=\tfrac12(2\pi v)^{1/4}e^{-D-Ry}.
\]
The same Poisson limit holds for outer vertices, so in both cases the probability of successful resolution tends to \(e^{-\Lambda_q(y)}\). As a consequence, if we replace \(t_*/s\) by \(t/s\) and omit the \(s^{-2}\) term, success holds with high probability if and only if \(t>t_*\). Likewise, at \(t=t_*\), adding a term \(u\log n/s^2\) yields high-probability success for \(u>u_*\) and high-probability failure for \(u<u_*\).

The proof proceeds in three steps. First, we refine the Fourier localization of Section~\ref{sec:classical-birthday} to obtain a sharp asymptotic formula for the single-pair collision probability \(P_d\) across the entire typical distance band. Second, we evaluate the resulting scalar sums through constant order to compute the expected collision count. Third, we prove by the method of factorial moments that the collision count converges to a Poisson limit.

\step{The single collision atom and its gauge.}
Our first goal is to derive an exact-constant asymptotic expression for \(\log P_d\), uniform over all vertex pairs at typical distance \(|d-\mu|\le n^{11/12}\).

As in Section~\ref{sec:classical-birthday}, we localize the Fourier integral for a single collision after deleting one central row \(c\) nearest to \(\mu\). The two constraints \(d_x=r,d_y=c\) leave residuals of size \(O(n^{7/12})\) after the shell tilt and rounding, and the same two unimodular coordinate moves produce the required singleton column intensities. On the small rows, we split the phases of the two endpoints into separate factors; the error introduced by this splitting is bounded by a constant times the sum of joint intensities over pairs of small radii, which is \(o(1)\) by the two-tilt intensity bound. Once separated, each independent small-row factor integrates exactly to \(e^{-2\lambda_r}I_0(2\lambda_r)\).

For the dense rows, we whiten the covariance matrix by the marginal standard deviations of the rows. On the full radial shell space, write \(w_r=\sqrt{\pi_r}\). In these coordinates, the joint-radii operator takes the form \(ww^T+F\), where \(Fw=0\) and the operator \(F\) has eigenvalues \(f_k(d)\), \(k\ge1\), in the centered radial basis. Combining the geometric bound on \(f_k\) for large degrees with the fixed-degree decay bound shows that the trace norm satisfies \(\|F\|_1=o(1)\). Let \(E\) denote the inclusion matrix of the dense undeleted rows, and set \(N=E^T(I-ww^T)E\). Then the matrix \((I-ww^T)EN^{-1/2}\) is isometric, which implies that the compressed dense covariance matrix, normalized by its diagonal scale \(2\lambda_r\), has determinant \((1+o(1))\det N=(1+o(1))\pi_c\). The second equality follows from the matrix determinant identity \(\det N=1-\sum_{\mathrm{dense}\ r\ne c}\pi_r\), together with the fact that the total shell mass of all omitted rows other than \(c\) is \(\exp(-As+o(s))\).

In the resulting product over dense rows, each independent Gaussian factor \((4\pi\lambda_r)^{-1/2}\) can be replaced by the exact Bessel factor \(e^{-2\lambda_r}I_0(2\lambda_r)\) at a total relative error of \(1+o(1)\), by the large-argument expansion of \(I_0\). Setting \(h_r=\log(4\pi\lambda_r)\) and combining the small and dense contributions, we obtain the single-atom formula
\[
 \log P_d=
 -\sum_r\{\tfrac12(h_r)_++G(h_r)\}
 +\tfrac12 h_c-\tfrac12\log\pi_c+o(1).                 \tag{*}
\]
Every error term in this derivation, including the relative Fourier localization error, is uniform across the entire typical distance band.

\step{Scalar sums through constant order.}
We now evaluate the two sums in (*) up to \(o(1)\) error, allowing general parameters \(t,u\) in the formula for \(K_n\) so that we also recover the coarser transitions.

Write \(p_0=1-1/q\) and \(x=(r-\mu)/s^2\). Uniformly on a fixed moderate window around \(x=0\), Stirling's expansion of the binomial shell mass (whose factorial remainder is \(O(s^{-3})\)) yields
\[
 h_r=sF_0(x)+g(x)+s^{-1}j(x)+s^{-2}k(x)+O(s^{-3}),
\]
where the coefficient polynomials are
\[
 \begin{split}
 F_0(x)&=A-x^2/(2v),\quad g(x)=-c_3x^3,\\
 j(x)&=t+b x-c_4x^4,\quad b=(2p_0-1)/(2v),\\
 k(x)&=u\log n+y+(1-2v)x^2/(4v^2)
                      +(\text{an odd quintic monomial}).
 \end{split}
\]
Here \(c_3,c_4\) are the same constants defined in the preceding section.

We first evaluate the clipped half-sum \(\sum_r \tfrac12(h_r)_+\). Extending the calculation from Section~\ref{sec:classical-birthday}, the contribution through order \(s\) is \(2Ln+s[Rt-R^5(1+2v)/(360v^3)]\). We claim that the remaining constant-order and logarithmic contribution is
\[
 R(u\log n+y)+\frac{(1-2v)R^3}{12v^2}-c_3 b v R^3
   =R(u\log n+y)+\frac{R^3}{6v},
\]
up to a remainder of \(o(1)\). To derive the boundary correction near each endpoint \(x=\pm R\), introduce the local coordinate \(z=F_0(x)\) with positive Jacobian \(a(z)=|dx/dz|\), and set \(\epsilon=s^{-1}\). Integrating between \(z=0\) and the shifted root \(z=-\epsilon g+O(\epsilon^2)\) on each local branch shows that the boundary correction to the normalized integral of \((F_0+\epsilon g+\epsilon^2j+\epsilon^3k)_+\) is \(\epsilon^2 a g^2/2\) at second order and \(\epsilon^3(a g j-a g^2 g_z/2-a_z g^3/6)\) at third order, with all functions evaluated at \(z=0\). When we sum the two opposite branches, the last two third-order terms cancel by parity, leaving only the product of \(g\) with the odd part of \(j\). The trapezoidal rule error for replacing the discrete sum by this integral is \(O(s^{-1})\), because the derivative with respect to \(x\) of the clipped polynomial has total variation \(O(s)\) (including the jumps at the two corners) on a mesh of spacing \(s^{-2}\). Finally, the \(O(s^{-3})\) remainder in the shell approximation contributes \(o(1)\) after summing over the window.

We next evaluate the Bessel correction sum \(\sum_r G(h_r)\). The same trapezoidal estimate applies here, since the exponential decay of the piecewise derivatives of \(G\) and the bounded number of monotone pieces imply that the derivative in \(x\) of \(G(h_r)\) has total variation \(O(s)\). In fixed neighborhoods of the two roots \(x_0=\pm R\), we change variables to \(h\). Inverting the expansion of \(h_r\) gives \(x-x_0=(h-g(x_0))/(sF_0'(x_0))+O((1+|h|^2)/s^2)\), and expanding the derivative at \(x_0\) shows that uniformly for \(|h|\le C\log n\), the density of shells is
\[
 s^2|dx/dh|
 =\frac{sv}{R}
  +\frac{v^2}{R^3}(h-g(x_0))
  -\frac{v^2}{R^2}\operatorname{sgn}(F_0'(x_0))g'(x_0)
  +O((1+|h|^2)/s).
\]
When we add the contributions from the two branches \(x_0=\pm R\), both terms involving \(g\) cancel by symmetry. The integrable exponential tails of \(G\) justify truncating the integral at \(|h|\le C\log n\), while outside the two neighborhoods of \(\pm R\) the sum is exponentially small (with the remote shell tail bounded as before). Integrating term by term gives \(\sum_rG(h_r)=2svJ_0/R+2v^2J_1/R^3+o(1)\).

Lastly, for the deleted central row \(c\), we have \(h_c=As+o(1)\) and \(\pi_c\sim(2\pi vn)^{-1/2}\). Substituting all three pieces into (*) yields
\[
 \log P_d=-2Ln-sR(t-t_*)+(1/4-Ru)\log n
             +\tfrac14\log(2\pi v)-D-Ry+o(1).
\]

To pass from typical pairs to the full expected collision count, we rule out atypical distances. For pairs outside the typical band with \(d\ge\delta n\), the entropy loss of \(\Omega(n^{5/6})\) outweighs the \(O(n^{2/3}\log n)\) loss from the balanced descending forest bound. For close pairs with \(1\le d<\delta n\), choosing \(\delta\) small enough ensures that the uniform pair bound \(e^{-2Ln+o(n)}\) beats the total number of such pairs \(e^{(L+b_q(\delta))n+O(\log n)}\). Hence all atypical single-pair collisions have total probability \(o(1)\). Since the number of typical unordered pairs is \((1+o(1))q^{2n}/2\), the expected number of collisions tends to \(\Lambda_q(y)\) in the centered window. Away from this window, the claimed high-probability success and failure regimes follow from the first and second moments of the preceding section, now applied over the full central distance band.

\step{Fixed factorial moments.}
To complete the proof of the Poisson limit, we compute the \(k\)-th factorial moment of the collision count for each fixed integer \(k\). We restrict to the count \(M_n\) of unordered collisions whose pair lies in the typical internal-distance band; as just shown, the expected number of omitted collisions is \(o(1)\).

Fix a factorial order \(k\) and consider an ordered \(k\)-tuple of distinct unordered candidate pairs, viewed as candidate edges on their endpoints. Choose a sufficiently small constant \(\delta>0\), join any two endpoints whose Hamming distance is at most \(\delta n\), and partition the endpoints into the resulting connected clusters. Each cluster has diameter at most \(2k\delta n\), and because every candidate pair lies in the typical distance band, no candidate edge has both endpoints in the same cluster. Contracting each cluster to a single vertex produces a multigraph of candidate edges with \(b\) vertices, \(c\) connected components, and no isolated vertices. For any specified cluster partition of the endpoints, the number of vertex configurations is at most
\[
 \exp\{bLn+2k b_q(2k\delta)n+O_k(\log n)\}.
\]
Furthermore, by a union bound on the distances from each endpoint in a cluster to a fixed cluster representative, every cluster has a linear number of coordinates on which all of its endpoints agree.

\step{Spanning-tree construction and non-disjoint configurations.}
We first show that configurations with overlapping clusters or parallel edges contribute negligibly to the \(k\)-th factorial moment, by constructing an independent system of shell-discrepancy columns along a spanning forest.

In each connected component of the contracted graph, choose a rooted spanning tree and retain one candidate edge per tree edge. For each nonroot cluster, anchor one word construction at the endpoint incident to its parent tree edge. Fix a shell of radius \(r\) around this anchor in the main window, excluding a fixed central interval around \(\mu\). Under the tilted shell law centered at the anchor, partition the coordinates into cells according to the full tuple of all endpoint letters.

We round the expected letter counts in each cell so that every comparison with another endpoint moves inward toward \(\mu\). When \(r>\mu\), in each large cell we round the counts of all nonanchor letters upward and assign the remaining count to the anchor letter. On any cell where a same-cluster comparator differs from the anchor, its distance to the landmark minus the anchor's distance equals the count of anchor letters minus the count of comparator letters; our rounding ensures this difference is at most its tilted mean, which is strictly negative. In small cells we place no anchor letters at all, so the difference is nonpositive for every same-cluster comparator. If a designated same-cluster comparator is required to move strictly inward, we place its letter on every small-cell coordinate where it differs from the anchor. When \(r<\mu\), we instead round all nonanchor letter counts downward, making the signed distance difference nonnegative for all same-cluster comparators and strictly positive where a designated comparator differs. Every cell undergoes only a bounded adjustment, and we restore the exact total radius \(r\) inside a single linear-size cell where all endpoints of the anchor cluster agree.

Under this allocation, all same-cluster comparisons remain weakly inward on the same side of \(\mu\), and any designated same-cluster comparator is strictly inward. Meanwhile, all comparisons with endpoints outside the anchor cluster are uniform contractions, because their distance from the anchor exceeds \(\delta n\) and the negative coordinatewise correlation is bounded in modulus by \(1/(q-1)<1\). Once the excluded central strip around \(\mu\) is wide enough to absorb the bounded rounding errors, every external comparison is strictly inward as well. By the multinomial modal-count estimate, the prescribed allocations have conditional probability at least an inverse polynomial in \(n\) under the shell law.

Now examine the incidence matrix of these word classes on the discrepancy rows of the retained tree edges, orienting each tree edge by its child endpoint. At the layer of maximal absolute deviation from \(\mu\), the column associated with a nonroot anchor has entry \(+1\) in its parent-edge row and possibly nonzero entries in the child-edge rows of its own cluster at the same radius. Ordered from parent to child, these leading incidence columns form a triangular matrix with diagonal entries equal to one. Radii on opposite sides of \(\mu\) at the same absolute level cannot interfere, because near comparisons stay on the same side of the mean and all other comparisons move strictly inward. Descending level by level toward \(\mu\) proves that all of these columns are linearly independent, and therefore the corresponding exact column classes are disjoint.

Conditional on all other landmark selections, the zero-discrepancy system admits at most one vector of counts for these independent binomial column classes. Multiplying the maximal point probabilities of those binomial variables gives a joint collision bound of \(\exp[-(b-c)an+O_k(n^{2/3}\log n)]\), since each class loses at most a polynomial factor relative to its shell intensity. Because \(a=2L\), whenever \(b>2c\) the product of this probability bound with the configuration count is exponentially small once \(\delta\) is chosen small enough.

The remaining non-disjoint case is \(b=2c\), in which every connected component consists of two clusters. If any component contains more than one candidate edge, pick two distinct edges in that component and anchor them on the side where their endpoints differ. In the two resulting anchor constructions, designate the other anchor as strictly inward using the modification described above. Their leading columns are then diagonal on the two chosen edge rows. Selecting one anchor per remaining component as before, all external-cluster contributions move strictly inward, yielding \(c+1\) independent columns per radius instead of \(c\). This extra factor makes all such multi-edge configurations exponentially negligible as well.

\step{Separated typical configurations and conclusion.}
The only configurations that remain have \(2k\) distinct endpoints lying in pairwise separated clusters, with a single candidate edge in each of the \(k\) components. If any \(2k\)-tuple cell count deviates from its mean \(n/q^{2k}\) by more than \(n^{11/12}\), Hoeffding's inequality bounds the proportion of such endpoint tuples by \(\exp(-\Omega_k(n^{5/6}))\), which defeats the \(O_k(n^{2/3}\log n)\) loss in the joint atom bound.

For all remaining typical \(k\)-tuples of pairs, we apply the Fourier comparison of Section~\ref{sec:classical-birthday} simultaneously across all \(k\) blocks. To build a singleton column in a given block, we prescribe its two radii as before and impose equality between the two endpoint distances in each of the other \(k-1\) blocks. Each additional equality constraint is adjusted by an independent unit move, using the cell whose endpoint tuple has a single letter \(1\) at the positive endpoint of that block and changing the landmark letter \(1\to2\). Because every tuple cell has linear size, we recover the same uniform singleton intensity lower bound, and there are only \(O_k(1)\) cross-endpoint pairs to control when splitting the small and dense phases. After whitening by the marginal block covariances, the joint dense Gaussian covariance matrix has identity blocks on the diagonal and off-diagonal blocks of Frobenius norm \(o(1)\); its trace-zero perturbation therefore has log determinant \(o(1)\). As in the two-block case, the small-row factors cancel, giving uniform relative factorization of the joint collision probability into the product of the \(k\) individual collision atoms.

Since there are \((1+o(1))q^{2kn}/2^k\) typical ordered \(k\)-tuples of unordered pairs, combining our single-atom asymptotic with the bounds above shows that \(\mathbb E[(M_n)_k]\to\Lambda_q(y)^k\) for every fixed \(k\). Because the factorial moments uniquely determine the Poisson distribution, \(M_n\) converges in distribution to a Poisson random variable with mean \(\Lambda_q(y)\). The omitted ordinary collisions vanish in probability, so the full ordinary collision count has the same asymptotic distribution. Finally, the difference between the ordinary and outer collision counts also vanishes in probability: by the exact selector ratio from Section~\ref{sec:classical-birthday}, the expected difference is \(O(p^2)\) times the bounded expectation of the full ordinary count. This establishes both claimed Poisson limits without needing higher moments of the discarded pairs.

\section{Binary vertex collisions in the narrow window}
\label{sec:binary-orbit-poisson}

In this section we determine the limiting distribution of vertex collisions in the binary hypercube (\(q=2\)) inside the narrow scaling window. Throughout, evaluate the constants of Section~\ref{sec:classical-poisson} at \(q=2\), with \(v=1/4\), \(s=n^{1/3}\), and \(L=\log 2\), and set
\[
 t_2=\frac{1}{R}\left(A+\frac{R^5(1+2v)}{360v^3}
                      -\frac{2v}{R}J_0\right),
 \qquad u_2=\frac{1}{2R}.
\]
Under independent Bernoulli landmark selection with mean
\[
 K_n=\kappa_1\sqrt n\exp\left(As+\frac{t_2}{s}
                         +\frac{u_2\log n+y}{s^2}\right)
\]
for a fixed parameter \(y\), the numbers of ordinary and outer colliding unordered vertex pairs each converge in distribution to twice a Poisson random variable with mean
\[
 \Lambda_2(y)=\frac{\sqrt{2\pi v}}{16}e^{-D-Ry}.
\]
In particular, the resolution probability converges to \(e^{-\Lambda_2(y)}\). The same limit law holds when the landmark set is a uniform random subset of size \(K_n\) rounded to an integer.

The factor of two in the limit reflects bitwise complement symmetry: complementing a vertex reverses its distance histogram, so nonantipodal colliding pairs occur in complement orbits of size two. Our proof proceeds in four steps. We first compute the collision probability for a single pair of vertices at a typical even distance, then show that odd, antipodal, and atypical distances make negligible contributions, next compute the fixed-order factorial moments of typical even-distance orbits via spanning-forest bounds and Fourier factorization, and finally transfer the result from independent Bernoulli selection to uniform subsets of fixed cardinality.

\step{The even-distance atom.}
We begin by computing the collision probability \(P_d\) for a pair of vertices at a typical even distance \(d\). Set \(h_r=\log(4\pi\lambda_r)\) and \(S=\sum_r(\tfrac12(h_r)_++G(h_r))\), and call an internal distance typical when \(|d-n/2|\le n^{11/12}\).

When two vertices lie at an even Hamming distance, the distances from any landmark to the two vertices always have the same parity. Consequently, the discrepancy between their distance histograms sums to zero separately over the even-radius shells and over the odd-radius shells. To remove these two linear dependencies, we delete one nearest central radius of each parity, denoted \(c_0,c_1\). The Fourier singleton construction of Section~\ref{sec:classical-birthday} then uses the deleted radius of the prescribed parity, and the errors incurred in splitting the two small-row endpoint phases are bounded by the same two-tilt estimate as in Section~\ref{sec:classical-poisson}.

In shell coordinates, the resulting whitened dense covariance is the compression of \(I-P-F\), where \(P\) projects onto the constant and alternating radial directions, and \(F\) has radial eigenvalues \(f_j(d)\) for \(1\le j<n\), and zero otherwise. Uniformly over the typical distance band, the trace norm satisfies \(\|F\|_1=o(1)\). To see this, use the symmetry of the spectrum to reflect the degrees to \(j\le n/2\), then sum the replacement bound \(|1-2d/n|^j+O(j^2/n)\) for degrees up to \(C\log n\) and the opposite-sign matching bound \(O(\sqrt j\,e^{-j/3})\) for the remaining degrees.

Let \(E\) be the inclusion of the dense undeleted rows, and set \(N=E^*(I-P)E\). Then \((I-P)EN^{-1/2}\) is an isometry, so the normalized determinant is \((1+o(1))\det N\). Since each parity class has total uniform mass \(1/2\), the quantity \(\det N\) equals the product of the omitted shell masses conditional on each parity, which is \((1+o(1))4\pi_{c_0}\pi_{c_1}\), where \(\pi_r\) denotes the uniform shell probability. Combining the small-shell factors with the diagonal dense factors produces the same scalar \(S\) as in Section~\ref{sec:classical-poisson}, with corrections coming only from the two deleted rows and the determinant:
\[
 \log P_d=-S+\tfrac12(h_{c_0}+h_{c_1})
             -\tfrac12\log(4\pi_{c_0}\pi_{c_1})+o(1)
         =-S+As+\tfrac12\log n+\tfrac12\log(2\pi v)-\log 2+o(1).
\]
Because the scalar expansion in Section~\ref{sec:classical-poisson} does not depend on the choice of gauges, evaluating it at the displayed mean \(K_n\) yields \(S=2Ln+As+\tfrac12\log n+D+Ry+o(1)\).

\step{Excluded events.}
We next show that odd-distance collisions, antipodal collisions, and collisions at nontypical even distances all have negligible effect on the limiting distribution.

An odd-distance collision forces the numbers of selected landmarks in the two global parity classes to be equal. By the maximal binomial atom bound, this equality occurs with probability \(O(K_n^{-1/2})\). This excludes all odd-distance collisions in probability, though not necessarily in expectation.

For an antipodal pair, having identical distance histograms is equivalent to having equal landmark counts in each unordered complementary shell pair \(\{r,n-r\}\) with \(r\ne n-r\). Because distinct shell pairs are disjoint, these conditions are independent. Within each shell pair, the logarithm of the binomial difference characteristic function differs uniformly from the logarithm of the symmetric Poisson difference characteristic function by \(O(p^2\binom nr)\), and both characteristic functions are positive for all sufficiently large \(n\). Summing this discrepancy over all shells gives a total error of \(o(1)\). Since \(\lambda_r=\lambda_{n-r}\), the log collision probability of an antipodal pair is \(-S/2+\mathbf 1_{2\mid n}h_{n/2}/4+o(1)\), where the central term is included only when \(n\) is even. Multiplying this probability by the \(2^{n-1}\) antipodal pairs gives a quantity tending to zero.

For all other nonzero, nonantipodal distances outside the typical band, we use the one-anchor balanced-column construction established below. It bounds each collision probability uniformly by \(\exp(-2Ln+O(n^{2/3}\log n))\). Since restricting to the excluded distance band incurs an entropy loss of \(\Omega(n^{5/6})\), the expected number of collisions in this band is \(o(1)\).

\step{Fixed factorial moments on complement orbits.}
We now show that the fixed-order factorial moments of the colliding complement orbits converge to those of a Poisson random variable with mean \(\Lambda_2(y)\). Because complementing a vertex reverses its distance histogram, a nonantipodal unordered pair and its simultaneous complement define the exact same collision event. Define the retained count to be the number of colliding complement orbits of typical even-distance pairs.

Fix a factorial order \(k\), and consider an ordered tuple of \(k\) candidate pairs with \(2k\) labeled endpoints. Cluster the \(2k\) endpoints by taking the transitive closure of the quotient-distance relation \(\min(d,n-d)\le\delta n\). In each cluster, select a representative vertex and orient all other words in the cluster (by complementing them when needed) so that each lies within distance \(2k\delta n\) of the representative; the raw orientation choices cost only a constant factor. Contracting each cluster to a single vertex produces a loopless multigraph of candidate edges with \(b\) clusters and \(c\) nontrivial connected components. The number of configurations with these parameters is bounded by \(\exp(bLn+2k b_2(2k\delta)n+O_k(\log n))\).

\step{Spanning-forest column bounds.}
To bound the joint collision probability of a clustered configuration, choose a rooted spanning forest of the contracted multigraph, keeping one edge per parent-child link and anchoring that edge at its child endpoint. For each absolute-deviation layer \(\{r,n-r\}\) outside a fixed central strip, we build an exact column at each of the two radii using the shell tilt and rounding procedure of Section~\ref{sec:classical-poisson}.

First orient every comparator vertex inside the anchor's cluster to lie close to the anchor. The comparison radii then remain weakly inward on the same side of the central strip, and by the same cell rounding we can make the radius strictly inward against one nominated distinct aligned comparator. Restoring the total count in an agreement cell does not disturb these inequalities. Undoing the orientations may reflect some radii across the center, but it leaves all absolute deviations unchanged. Meanwhile, comparators in other clusters contract these deviations by a fixed factor, and removing the fixed central strip absorbs rounding errors so that their radii are strictly inward as well.

Each exact column class retains at least a polynomial fraction of its anchor shell. At the leading absolute-deviation layer, arranging the columns in parent-before-child order makes the forest matrix block triangular: each diagonal block is a signed permutation of the two radii in the layer, and the only possible extra entries occur in child-edge blocks. Descending layer by layer proves linear independence of the columns. Conditioning on the landmarks outside these disjoint classes bounds the collision probability by the product of their maximal binomial atoms, contributing a cost of \(2Ln+O_k(n^{2/3}\log n)\) for each complete two-radius family. This same construction also proves the one-anchor bound used earlier for nontypical distances: any quotient-separated comparator contracts, while a close nonidentical aligned comparator is made strictly inward at both radii. Whenever \(b>2c\), the \(b-c\) complete families beat the configuration entropy bound once \(\delta\) is chosen sufficiently small.

\step{Parallel edges and folded constraints.}
When \(b=2c\), every nontrivial component has two clusters, and any component containing parallel edges supplies extra columns. If two parallel edges in a component have distinct quotient endpoints on at least one side, we anchor both edges on that side and, in each column construction, nominate the other aligned anchor to be strictly inward. Both leading two-radius blocks are then diagonal up to signed permutations, providing two complete families.

If instead both endpoint quotient distances are zero, then distinct complement orbits must have opposite relative complement signs. Choosing representatives \(X,Y\) for the two clusters, the two collision conditions imply \(H_X=H_Y\) and \(H_Y=JH_Y\), where \(J\) reverses the radius order. In each layer with \(r>n/2\), we retain the two rows of \(H_X-H_Y\) together with the single row \(H_Y(r)-H_Y(n-r)\). Because the opposite cluster is strictly inward, the two \(X\)-anchor columns and the \(Y\)-anchor column at radius \(r\) form a triangular leading matrix. Having three columns per layer rather than two adds half the cost of a complete family. In both cases, retaining one complete family in each of the remaining components yields a total probability bound that defeats the configuration entropy.

\step{Typical tuples and factorial moments.}
The only configurations not yet eliminated consist of single-edge components with distinct quotient endpoints. The spanning-forest bound above also beats the \(\exp(-\Omega_k(n^{5/6}))\) fraction of such tuples in which some coordinate-type count differs from \(n/2^{2k}\) by more than \(n^{11/12}\).

For the remaining typical tuples, the fixed-order Fourier factorization from the preceding sections extends with two gauges per block. Specifically, translate one anchor to zero and impose three conditions: its own radius \(r\), the deleted radius \(c\) of matching parity for its partner, and zero discrepancy in every other pair. Writing \(y\) for the translated partner and \((u_j,v_j)\) for the other translated pairs, the coefficient columns for the landmark bits are \((1,y_i,(v_{ji}-u_{ji})_j)\), and the corresponding right-hand sides are \(r,(r+|y|-c)/2,((|v_j|-|u_j|)/2)_j\), all of which are integers. Coordinate types with all zero bits, or with a single one at the partner or at the positive endpoint of another pair, yield a unimodular minor supported on linear-size cells. This gives the same singleton adjustment and intensity estimates as in the single-pair analysis.

In the cross-block radial Gram bounds, degrees 0 and \(n\) vanish identically in every even-distance block and are therefore omitted. The off-diagonal whitened blocks have Hilbert--Schmidt norm \(o(1)\), and because the off-diagonal matrix has zero trace, the ratio of the joint determinant to the product of the individual block determinants is \(1+o(1)\). All small-shell and dense-shell splitting errors are summed over only a bounded number of block pairs, which yields uniform relative factorization of the collision atoms.

After all typicality restrictions, the fraction of ordered tuples with all internal parities even is \(2^{-k}+o(1)\). Each ordered list of \(k\) complement orbits is represented \(4^k\) times by labeled tuples, accounting for swapping the endpoints and taking their simultaneous complement in each pair. The fixed factorial moments of the retained orbit count therefore converge to the corresponding powers of the stated Poisson mean \(\Lambda_2(y)\). Since the retained ordinary vertex collision count is exactly twice the orbit count, the exclusions established earlier transfer this limit law to the full ordinary count. For outer vertices, the radius-zero independence identity shows that the expected number of retained physical pairs lost because an endpoint is selected is \(o(1)\), as their ordinary mean is bounded. Since the excluded outer collision events are subevents of the excluded ordinary events, the limit law transfers to outer vertices as well.

\step{Fixed cardinality.}
Finally, we transfer the limit law from independent Bernoulli selection to a uniform landmark subset of fixed size \(K_n\).

For typical even-distance tuples, omitting the anchor-radius constraint from the unimodular construction above and imposing equal distances on every pair yields a neutral class of polynomial mass. The relative iid and no-repeat transfer of Section~\ref{sec:fixed-cardinality} therefore applies unchanged to both single and joint typical atoms. In the remaining factorial-moment bounds and in the excluded nonantipodal even-distance bounds, conditioning Bernoulli selection on its total size incurs a factor of only \(O(\sqrt{K_n})\), which is absorbed by their entropy margins. Odd-distance collisions are impossible when the sample size is odd; when it is even, global parity balance occurs with probability
\[
 \frac{\binom{2^{n-1}}{K_n/2}^2}{\binom{2^n}{K_n}}
 =(1+o(1))2^{-K_n}\binom{K_n}{K_n/2}
 =O(K_n^{-1/2}),
\]
where the relative approximation follows by expanding the falling factorials using \(K_n^2/2^n=o(1)\).

For antipodal pairs with even \(n\), we apply the relative transfer directly to a folded antipodal block indexed by coordinates \(r<n/2\). A word in shell \(r\) or shell \(n-r\) supplies the two signed unit columns of that coordinate, each with probability \(\pi_r\), while the central shell forms a neutral class of polynomial mass. The independent Poisson proxy is the exact positive product of the symmetric shell characteristic functions. The same dense-box estimate, centered quadratic bound, and shifted-atom estimate therefore give the folded atom with relative error \(o(1)\), preserving its vanishing total mean. For odd \(n\), parity already eliminates all antipodal collisions.

To complete the transfer for outer vertices, observe that for any retained pair, the probability of both colliding and having a specified endpoint selected is at most \((1+o(1))K_n2^{-n}\) times its Poisson atom. Indeed, we sum over the \(K_n\) iid positions that could take that endpoint, bound the required shifted atom of the remaining \(K_n-1\) positions by the same absolute Fourier integral estimate, and divide by the probability of having no repetitions, which tends to one. Summing over both endpoints and all retained pairs proves the outer transfer.

\section{Binary edge collisions in the narrow window}
\label{sec:binary-edge-poisson}

We now establish the analogous Poisson limit law for edge collisions in the binary hypercube. Let \(q=2\), \(N=n-1\), and \(s=N^{1/3}\). Evaluate the constants of Section~\ref{sec:binary-orbit-poisson} at \(q=2\), and set
\[
 u_{\rm edge}=\frac{3}{2R},\qquad
 K=\kappa_1\sqrt N\exp\left(As+\frac{t_2}{s}
                    +\frac{u_{\rm edge}\log N+y}{s^2}\right).
\]
For fixed \(y\), the number of colliding unordered pairs of distinct edges converges in distribution to twice a Poisson random variable of mean \(\Lambda_2(y)=\sqrt{2\pi v}\,e^{-D-Ry}/16\), with \(v=1/4\). In particular, the edge resolution probability tends to \(e^{-\Lambda_2(y)}\). This holds both under independent landmark selection with probability \(K/2^n\) and for a uniform random subset of size \(K\) rounded to an integer.

The argument reduces edge collisions to vertex collisions in dimension \(N\) by projecting each edge along its axis. Parity eliminates all nonparallel collisions and all odd-separation parallel collisions. On a single axis, the even-distance atom follows from Section~\ref{sec:binary-orbit-poisson}, and for tuples spanning multiple axes we project onto the common core of unshared coordinates to carry out the spanning-forest and Fourier factorization arguments.

\step{Parity and single atoms.}
We first show that nonparallel collisions and odd-separation parallel collisions occur with vanishing probability, and then compute the single-pair atom for parallel edges at a typical even distance.

Represent an edge by its axis \(i\) and its projected center in the \(N\)-cube obtained by deleting coordinate \(i\). Up to a sign depending on its center, the alternating sum of its shell counts equals the sum over all selected landmark words of \((-1)^{\sum_{h\ne i}z_h}\). Consequently, a collision between two nonparallel edges forces two such signed sums on distinct axes to agree up to sign. For either choice of sign, half of the cube contributes zero to the difference, and each nonzero sign is represented by a quarter of the cube. Under independent selection, the maximal binomial atom bound gives probability \(O(K^{-1/2})\) per comparison. In the uniform subset model, we use the iid draw model conditioned on no repetitions: condition first on the number of nonzero contributions, which is distributed as \(\mathrm{Bin}(K,1/2)\), and then apply the maximal atom bound to the independent fair signs. Since the probability of no repetition is \(1-o(1)\), the same bound holds, and a union bound over all pairs of axes and signs gives \(o(1)\). By the same argument, an odd-separation parallel collision forces the signed sum on its shared axis to equal zero, which occurs with probability \(O(K^{-1/2})\) per axis in either model.

On a single axis, each projected word has two preimages in the full cube, so in the Poisson occupancy model their counts simply add and the projected shell intensities are exactly \(K2^{-N}\binom Nr\). The single even-distance typical atom of Section~\ref{sec:binary-orbit-poisson} therefore applies with \(n\) replaced by \(N\). Compared to the vertex scaling, the new centering adds \(\log N\) to the scalar \(S\). This follows from the expansion in Section~\ref{sec:binary-orbit-poisson}, since the derivative in the displayed shift parameter is \(R+o(1)\), with an error that remains \(o(1)\) after multiplication by \(\log N\) by the clipped-sum and boundary expansions; equivalently, one can substitute the new coefficient of \(\log N\) directly into the scalar expansion. The resulting collision atom is
\[
 (1+o(1))\,N^{-1}2^{-2N}\frac{\sqrt{2\pi v}}2e^{-D-Ry}.
\]
The one-anchor bound from Section~\ref{sec:binary-orbit-poisson} removes all other even nonantipodal distances outside the typical distance band, even after multiplying by the \(n\) possible axes and, in the fixed-size model, by the conditioning penalty \(O(\sqrt K)\). Even projected antipodes have an expected count per axis bounded by \(\exp(-As/4+O(\log N))\), by the folded-shell atom of Section~\ref{sec:binary-orbit-poisson}. Its fixed-size relative transfer uses iid draws on the full cube: projection has the same one-draw discrepancy law, and excluding repetitions on full words changes the probability by only a relative error of \(o(1)\) by the shifted-atom estimate. Odd projected antipodes have already been excluded by parity.

\step{Cluster bounds across axes.}
Next we bound the contribution of clustered or atypical tuples of parallel edge pairs by restricting to the common core of coordinates outside their active axes.

Simultaneously complementing the two projected centers of a parallel edge pair reverses both histograms and leaves the collision event unchanged. We therefore count retained parallel collisions on these complement orbits, keeping the axis as part of the orbit label. For a fixed factorial tuple of order \(k\), let \(I\) be its set of active axes and put \(M=n-|I|\). Once the landmark letters on \(I\) are frozen, the distance to every edge in the tuple equals its common-core distance in dimension \(M\) plus a fixed offset in \([0,|I|-1]\). We cluster the endpoints on the common core using the same quotient-distance rule as in Section~\ref{sec:binary-orbit-poisson}. Choosing the active axes and the frozen endpoint letters costs only a polynomial factor. For \(b\) clusters in \(c\) connected components, the same entropy bound applies in dimension \(M\), and the contracted multigraph has no loops by the typicality of internal distances.

\step{Exact columns on the common core.}
To construct the exact columns needed for a rooted spanning forest, first align each close comparator to the anchor by complementing it when needed; absolute layers are measured from \(N/2\) and are unaffected by this operation. For the upper radius of a layer, choose all frozen letters seen by the anchor to mismatch it, giving the maximal offset; for the lower radius, choose them to match, giving offset zero. Every comparator offset is therefore weakly inward relative to the anchor's offset. If a nominated comparator uses a different axis, we can choose the anchor's free letter to match it at the upper radius or mismatch it at the lower radius, making its offset strictly inward. For a same-axis comparator with different aligned frozen letters, the offset is automatically strictly inward.

On the common core, tilt the mismatch probability to the requested core radius divided by \(M\), round the mismatch counts upward at the upper radius and downward at the lower radius in each tuple cell, and restore the total in a linear-size cell where all close aligned comparators agree with the anchor; such a cell exists once the fixed cluster parameter is taken sufficiently small. On the differing set of each close comparison, the mismatch count then lies on the inward side of half the set size. If a nominated comparison differs from the anchor on the core and still achieves equality, adjust one count on its differing set in the inward direction and compensate again in the common agreement cell. This preserves all weak inequalities.

Combining the core and offset comparisons gives weak inwardness for all close comparators and strict inwardness against any nominated distinct aligned edge. After excluding a sufficiently large fixed central strip, the contraction of the tilted mean together with bounded rounding and offsets prevents overshoot, and far comparisons are strictly inward by the same contraction bound. Applying Stirling's approximation to finitely many cells (including bounded-size cells and the fixed frozen pattern) gives polynomial relative shell mass. Consequently, the same leading-layer triangular argument and maximal binomial atom bound yield one complete family cost of \(2LN+O_k(N^{2/3}\log N)\) per forest link.

\step{Entropy comparison.}
For \(b>2c\), these \(b-c\) complete families defeat the configuration entropy bound. When \(b=2c\), every component has two clusters. If such a component has two candidate links whose aligned edges on some side are distinct, we anchor both links on that side and nominate the other anchor for strict inwardness. This works even for identical common-core words on different axes, and it supplies two complete families. If neither side has distinct aligned edges, the two links use the same axis and the same projected quotient endpoints. Distinct complement orbits then have opposite relative complement signs, and the third column per layer from the folded histogram constraint in Section~\ref{sec:binary-orbit-poisson} provides the extra half-family. In both cases the total cost defeats the entropy bound.

For the remaining single-link components, the spanning-forest bound also absorbs the \(\exp(-\Omega_k(N^{5/6}))\) fraction of common-core tuples whose coordinate-type frequencies differ from uniform by more than \(O_k(N^{11/12})\). All of these negligibility estimates absorb the fixed-size conditioning penalty.

\step{Relative typical joint atoms and counting.}
Finally, we show that the joint collision atoms of typical tuples factorize relatively into single-pair atoms, and we deduce the limiting Poisson distribution from the resulting factorial moments.

For typical tuples, carry out the Fourier singleton construction conditioned on a fixed frozen landmark pattern on \(I\). For any even projected pair, the difference between its frozen offsets is congruent modulo two to the distance between its common-core endpoints. Thus the full anchor radius and the deleted partner radius of matching parity impose compatible core radius constraints, and equality for every other pair likewise imposes compatible core difference constraints. The unimodular minor on linear-size common-core tuple cells is unchanged. Because the fixed offset shifts are bounded, the previous singleton intensity bound continues to hold, and omitting the anchor-radius constraint again gives a polynomial-mass neutral column. Moreover, conditioning on the frozen pattern preserves the two-tilt estimate for any two distinct typical endpoints: removing finitely many coordinates preserves the type and distance estimates and shifts radii only by bounded amounts. Thus the small-row splitting and dense-box estimates remain unchanged.

In the dense Gaussian comparison, the radial Walsh inner products vanish across different degrees. On the same axis, they are identical to the previous degree-wise inner products in dimension \(N\). Across different axes, retaining only equal Walsh supports from the two radial sums shows that the degree-\(j\) inner product equals the binary radial inner product on the \(N-1\) shared coordinates multiplied by \(\binom{N-1}{j}/\binom Nj=(N-j)/N\). Each even-distance block omits degrees zero and \(N\). On the shared core, the reflected degree zero at \(j=N-1\) contributes only \(1/N\), and the remaining squared bounds sum to \(o(1)\) by the same matching and replacement estimates. The cross-blocks therefore have whitened Hilbert--Schmidt norm \(o(1)\). Just as in the vertex proof, the joint determinant has relative error \(o(1)\), and the Fourier integrals factorize relatively. The neutral and singleton classes also verify the hypotheses of the relative iid and no-repeat transfer of Section~\ref{sec:fixed-cardinality}.

For each ordered list of \(k\) axes, before discarding the negligible counting fraction of atypical tuples, the parity restrictions select an exact fraction \(2^{-k}\) of all independent labeled projected endpoints. Every retained ordered orbit tuple has \(4^k\) labeled representatives. Consequently, the factorial moments are asymptotic to the corresponding powers of \(n2^{2N}/8\) times the single atom, and that product tends to \(\Lambda_2(y)\) because \(n/N\to1\). Since the retained physical count is exactly twice the orbit count, the exclusions above transfer the Poisson limit law to the full count.

\section{Edge collisions for alphabets of size at least three}
\label{sec:qary-edge-poisson}

We now establish the Poisson limit for colliding pairs of edges when the alphabet has size at least three. Fix \(q\ge3\). As in Section~\ref{sec:classical-poisson}, set \(s=n^{1/3}\) and \(p_0=1-1/q\), and retain the constants \(A,v,R,t_*,D,\kappa_1\) defined there. Define
\[
 \Delta=\frac{R^3(3p_0-2)}{6vp_0}=\frac{q-3}{q-1}\log q,\qquad
 K=\kappa_1\sqrt n\exp\left(As+\frac{t_*}{s}
              +\frac{(9/(4R))\log n+y}{s^2}\right).
\]
For any fixed \(y\), the number of colliding unordered pairs of distinct edges converges in distribution to a Poisson random variable with mean
\[
 \frac{(q-1)^2}{8}(2\pi v)^{1/4}\exp(-D-\Delta-Ry).
\]
Consequently, the resolution probability converges to the exponential of minus this mean. Both landmark models yield the same limit: independent selection of each vertex with probability \(K/q^n\), and uniform selection of a subset of size equal to the nearest integer to \(K\).

To prove this limit, we adapt the vertex argument of Section~\ref{sec:classical-poisson} to edges. We first determine how edge distances alter the shell intensities and the scalar expansion, next verify the Fourier and determinant estimates on typical configurations, and finally bound the atypical configurations and pass to factorial moments.

\step{Shell intensities and the scalar expansion.}
We begin by comparing the distance distribution from a uniformly random landmark to a fixed edge with the corresponding distribution for a vertex, and computing the resulting shift in the scalar expansion. The distance from a random landmark to an edge (the edge radius) is the sum of a \(\mathrm{Bin}(n-1,p_0)\) variable, arising from the coordinates where the two endpoints of the edge agree, and an independent \(\mathrm{Bern}((q-2)/q)\) variable, arising from the single active coordinate where the landmark mismatches both endpoint letters. Dividing this convolution by the vertex shell mass yields the exact ratio \(2-r/(np_0)\).

In the rescaled radial coordinate \(x=(r-np_0)/s^2\), the resulting change to the log-intensity profile \(h_r\) is \(-x/(p_0s)-x^2/(2p_0^2s^2)+O(s^{-3})\) on the moderate window. In the clipped half-sum, the odd linear term integrates to zero over the interior, while the quadratic term contributes \(-R^3/(6p_0^2)\). At the boundary, the shift in \(j\) contributes a correction of \(c_3vR^3/p_0\), as given by the two-branch formula in Section~\ref{sec:classical-poisson}. Adding the interior and boundary contributions yields \(\Delta\).

The remaining parts of the scalar expansion are unchanged to leading order: the boundary density for the \(G\) sum through constant order still depends only on \(F_0,g\), which are unaffected, and the central deletion terms change by \(o(1)\). Replacing the vertex coefficient \(u_*=1/(4R)\) by \(9/(4R)\) in the scalar expansion therefore yields the single-pair collision atom
\[
 (1+o(1))\,q^{-2n}n^{-2}(2\pi v)^{1/4}
   \exp(-D-\Delta-Ry),
\]
once the relative Fourier and determinant estimates below are established.

\step{Typical Fourier estimates.}
Next, we verify that the Fourier localization, dense-box approximation, and small-row splitting from the vertex analysis carry over to typical collections of edge pairs. The strategy is to freeze the landmark letters on the active coordinates of the edges, reducing all edge distances to ordinary vertex distances on a common core up to bounded offsets.

Fix a list of edge pairs, and freeze the landmark letters on the set \(I\) of their active axes. Conditioned on these frozen letters, every edge distance is an ordinary Hamming distance on a common core of dimension \(n-|I|\) plus a bounded frozen offset. Call the list typical when all common-core tuple frequencies of its listed endpoints lie within \(O(n^{11/12})\) of their uniform means.

To adapt the singleton construction of Section~\ref{sec:classical-birthday}, fix any frozen pattern on \(I\), subtract its offsets from the radius constraints, and subtract the corresponding offset differences from the equality constraints. Because the extra residuals introduced by the offsets are bounded, the same unimodular moves inside linear-size core cells still solve all constraints. Fixing the frozen pattern on \(I\) costs only a constant factor in probability. Omitting the anchor radius constraint and imposing equality in every block likewise yields the polynomial-mass neutral word class, obtained by rounding the untilted allocation on the core and correcting the bounded residuals. The two-tilt bounds also hold on each frozen fiber, because bounded shifts leave their error orders unchanged. Thus singleton localization, dense-box approximation, and small-row splitting all apply with the edge shell intensities. These singleton and neutral classes also supply the hypotheses required for the relative fixed-size transfer of Section~\ref{sec:fixed-cardinality}, including the step that excludes repeated draws on full words.

\step{Determinant estimates and factorization.}
We now compute the Gaussian determinant for a typical edge pair and show that the joint determinant for multiple typical pairs factors asymptotically into the product of the individual block determinants. Write \(\pi\) for the common edge distance marginal and \(w=(\sqrt{\pi_r})_r\) for the unit vector of square-rooted shell probabilities. For a frozen pattern \(z\), the shift map \(T_{e,z}f(t)=f(t+a_e(z))\) from the edge radial space to the core radial space has operator norm at most \(q^{|I|/2}\), because \(\pi\) is a mixture over the frozen patterns. For two typical distinct core centers, the centered core joint operator has trace norm \(o(1)\) by the radial Krawtchouk estimate of Section~\ref{sec:classical-poisson}; conjugating by the bounded maps \(T_{e,z}\) preserves this \(o(1)\) trace-norm bound.

We next control the conditional mean density \(g_{e,z}\) of the edge distance given \(z\), taken relative to \(\pi\). Because this density arises from conditioning on a frozen pattern on \(I\), it is bounded by \(q^{|I|}\). Moreover, bounded shifts of the core binomial distribution have total variation distance \(O(n^{-1/2})\) from the unshifted distribution by unimodality and telescoping. Combining this total variation bound with the uniform bound on \(g_{e,z}\) gives \(\|g_{e,z}-1\|_{L^2(\pi)}^2=O(n^{-1/2})\). Since the densities \(g_{e,z}\) average to one over \(z\), the conditional mean outer products average to the constant projection plus an operator of trace norm \(o(1)\).

In whitened coordinates, the full joint operator therefore takes the form \(ww^T+F\), where \(Fw=F^Tw=0\) and \(\|F\|_1=o(1)\). Here \(F\) need not be symmetric. For a discrepancy between two edges, dividing the covariance matrix by the marginal diagonal scale \(2\lambda_r\) gives \(I-ww^T-(F+F^T)/2\). Applying the isometric compression of Section~\ref{sec:classical-poisson} yields the determinant \((1+o(1))\pi_c\) after omitting one central row and restricting to dense rows. For distinct discrepancy blocks, the constant projections cancel in their cross-covariance, and the same isometries show that the cross-covariance has trace norm \(o(1)\) after whitening the individual covariance blocks. Hence the ratio of the joint determinant to the product of the individual determinants tends to one. The small-row factors likewise factor relatively, by the two-tilt and singleton estimates.

\step{Atypical configurations and counting.}
It remains to bound the atypical configurations and sum the typical contributions to obtain the factorial moments. We adapt the forest estimates of Section~\ref{sec:classical-poisson} to the common core by adding one rule for choosing the frozen landmark letters on the active coordinates.

For an upper-layer anchor, choose the frozen letters to mismatch each anchor factor; for a lower-layer anchor, choose them to match each anchor factor. Because every active factor of an edge is either a singleton or a two-letter set, both a matching choice and a mismatching choice exist for \(q\ge3\). Under this rule, every comparator's frozen offset is weakly inward. To obtain strictness against a specified close comparator, we consider two cases. If the comparator has a different core word from the anchor, the core rounding construction already gives strictness. If its core word equals the anchor's, the two edges are still distinct, so their frozen product sets are distinct subsets of equal cardinality. Thus some comparator factor contains a letter outside the corresponding anchor factor, and some anchor factor contains a letter outside the corresponding comparator factor; otherwise one product set would contain the other. Choosing the first letter for an upper-layer anchor and the second for a lower-layer anchor makes the frozen offset strictly inward. Core rounding, restoration in an agreement cell, and a sufficiently wide fixed central exclusion strip then give the same leading-layer independence and inverse-polynomial relative shell mass as for vertices.

We now apply these forest bounds to eliminate atypical pairs and tuples. For a single close pair, nominate its other edge as the close comparator; the resulting single complete family gives an atom bound of \(\exp[-2Ln+O(n^{2/3}\log n)]\). At a fixed, sufficiently small cluster threshold \(\delta\), the number of such close pairs is at most \(\exp[(L+b_q(\delta))n+O(\log n)]\), so their total contribution vanishes. For a well-separated atypical pair, the forest gives the same atom bound, while the nonuniform core type frequencies reduce the count by a fraction \(\exp[-\Omega(n^{5/6})]\). Together these bounds remove all atypical pairs, including pairs with identical cores and differing frozen factors.

For factorial tuples with \(b\) clusters in \(c\) connected components, the same forest produces \(b-c\) complete families. In the case \(b=2c\) where some component has two distinct links, the anchors of those links differ on at least one side, so the strict comparison rule produces an extra complete family. The unchanged entropy comparison eliminates both of these clustered cases, and any remaining nontypical tuples incur the same type-frequency penalty. Choosing the active axes and frozen endpoint letters costs only polynomial factors, and all of these negligibility bounds absorb the \(O(\sqrt K)\) penalty from conditioning on a fixed landmark count.

\step{Counting and factorial moments.}
There are \(n(q-1)q^n/2\) edges in total. Typical unordered pairs have asymptotically the full unordered pair count, and typical ordered factorial tuples have asymptotically its corresponding powers: for any choice of active axes, the free core endpoints have uniform independent coordinate tuples before atypical or repeated pairs are discarded. Parallel pairs (sharing the same active axis) account for an \(O(1/n)\) proportion of all pairs and do not require a separate typical atom. Combining the relative factorization with the single-pair atom therefore shows that every fixed factorial moment converges to the corresponding power of the displayed Poisson mean. Since atypical single pairs vanish in probability, this proves the asserted Poisson and resolution probability limits.

\section{Transfer to exactly \(K\) landmarks}
\label{sec:fixed-cardinality}

In this section we show that the classical sparse probability thresholds also hold when the landmarks form a uniformly sampled \(K\)-element subset, with \(K\) rounded to an integer at the prescribed scales. For ordinary and outer vertices over alphabets of size \(q\ge3\), the Poisson law of Section~\ref{sec:classical-poisson} and its smaller probability windows likewise hold in this model. This transfer does not rely on monotonicity of resolution. The all-face results and the classical exponential failure rates are not transferred in this section.

The argument proceeds in two steps. On typical configurations, we represent the uniform subset by conditioning \(K\) independent draws with replacement on having no repetitions, and compare the resulting Fourier integral directly to the Poisson model with relative error \(o(1)\). On atypical configurations, we condition independent Bernoulli landmarks on having total count \(K\), transferring all earlier upper bounds at a multiplicative cost of \(O(\sqrt K)\).

\step{Relative control for typical configurations.}
We first draw \(K\) independent words with replacement and compare the zero-discrepancy probability to the Poisson model. Let \(D\) be the stack of one-word signed discrepancy columns in gauged coordinates, and write \(\phi(\theta)=\mathbb E e^{i\theta\cdot D}\) for its characteristic function. The zero-discrepancy atom is then the integral of \(\phi^K\) over the torus with normalized Lebesgue measure.

To control the integral outside the dense frequency box, we verify the neutral and singleton lower bounds for a single draw. For the typical configurations used in Sections~\ref{sec:classical-birthday} and \ref{sec:classical-poisson}, the neutral word class \(D=0\) has probability at least \(n^{-O_k(1)}\). To see this, round the allocations of uniformly distributed landmark letters across the finitely many typical endpoint tuple cells, and correct the bounded residual in each pair's signed distance difference. For \(q\ge3\), a unit correction in a single block is given by choosing a single exceptional center letter and changing the landmark letter to a third letter, just as in the singleton construction. For \(q=2\), the internal candidate separations are even, so every residual is even, and a lone-bit tuple allows a change of two in that block alone. When parallel edges are present, we restrict to the common zero-offset fiber, which has positive probability. Because all correction cells have linear size, the multinomial bounds imply that the rounded allocations carry inverse-polynomial mass. By the same construction, the singleton class of each retained central row still has probability at least \(\lambda_r K^{-1}\exp[-O_k(n^{1/6}+\log n)]\).

\step{Fourier comparison with the Poisson model.}
Let \(D'\) be an independent copy of \(D\). The symmetrization identity
\[
 1-|\phi|^2=\mathbb E[1-\cos(\theta\cdot(D-D'))]
\]
together with the neutral and singleton classes gives the same superexponential off-box bound as in the Poisson model, with only an additional polynomial loss in the witness mass inside the exponent. In particular, the off-box integral is negligible relative to the corresponding Poisson zero atom, and this remains true with \(K-2\) draws in place of \(K\).

On the dense frequency box, we have \(\mathbb E(\theta\cdot D)=0\), because the two members of each pair have equal marginal shell laws. Applying the bound \(|e^{ix}-1-ix|\le x^2/2\) shows that
\(K|\phi-1|\le \tfrac12K\mathbb E(\theta\cdot D)^2
=\exp(o(n^{1/3}))\).
Indeed, each landmark word has at most \(2k\) nonzero terms: the dense terms cost at most \(O_k(nH^2)\), the small terms cost at most \(O_k(n e^{n^{1/4}})\), and the previous tail cutoff has negligible total intensity, where \(H=e^{n^{1/5}}\) as before. Since this implies \(K|\phi-1|^2=o(1)\), we obtain
\[
 \phi^K
  =\exp(K(\phi-1))(1+o(1))
\]
uniformly on the dense frequency box. The same estimate holds for \(K-2\) powers relative to the original exponential. Because the Poisson exponent has the positive proxy established previously, integrating \(|\phi|^{K-2}\) in absolute value also gives \(1+o(1)\) times the Poisson zero atom. In particular, every shifted atom of the \(K-2\) sum is bounded by that same quantity.

\step{Excluding repeated draws.}
We now condition the independent draws on having no repetitions to pass to a uniformly chosen subset. If two specified positions draw the same word, the other \(K-2\) draws remain independent and identically distributed, and for the total discrepancy to vanish, their discrepancy sum must equal the negative of that shared word's column counted twice, which is a shifted atom of the \(K-2\) sum. Summing over all pairs of positions and applying the shifted-atom bound, the joint probability of a repetition and zero discrepancy is at most
\[
 \binom K2 q^{-n}(1+o(1))\,P_{\rm Pois},
\]
where \(P_{\rm Pois}\) denotes the zero atom in the independent Poisson occupancy model of total mean \(K\). Since \(K^2q^{-n}=o(1)\), conditioning on no repetitions changes the zero atom by only \(o(P_{\rm Pois})\). This conditioning produces a uniformly chosen \(K\)-subset. Because the Bernoulli atom at parameter \(K/q^n\) already agrees relatively with the Poisson atom on typical configurations, we have transferred every typical single and joint atom with a uniform relative error \(o(1)\).

\step{Atypicals and conclusion.}
We now transfer the atypical bounds and deduce the threshold and Poisson limits. Conditioning independent Bernoulli landmarks with inclusion probability \(K/q^n\) on having total count \(K\) produces the uniform subset directly. By Stirling's formula, this conditioning event has probability \(\Theta(K^{-1/2})\), so every prior atypical upper bound transfers after multiplication by \(O(\sqrt K)\). This factor is negligible both in the exponential comparison for nonzero close matches and against the \(\exp(-\Omega(n^{5/6}))\) frequency loss for atypical well-separated tuples.

In the classical second moment method, exact repetitions and the binary complement variants are bounded directly by the actual single atom for the uniform subset rather than by the conditioning penalty. Because that single atom was transferred relatively, the remainder is still \(O(1/(MP))\) as before. This proves all high-probability failure assertions at the endpoint and below. High-probability success away from the narrow vertex window follows from the preceding first-moment bounds together with the \(O(\sqrt K)\) conditioning penalty, since the savings there are on the \(n\) or \(n^{2/3}\) scales, or larger. In the narrower \(q\ge3\) vertex window, we use the relative typical atoms instead; the discarded internal distances still satisfy the same negligible total bound after conditioning.

For outer vertices in the uniform subset model, an ordinary collision involving a landmark forces both queried vertices to be landmarks because of the radius-zero entry. On typical configurations, the corresponding event in the independent-draws representation requires two draws to hit those two vertices and the remaining \(K-2\) draws to achieve a shifted discrepancy, which costs at most \(K(K-1)q^{-2n}\) times the upper bound for a shifted \(K-2\) atom. Restricting to no repetitions therefore changes the typical collision atom by a relative \(o(1)\). For atypical configurations, the ordinary upper bounds also dominate the outer events. In the Poisson window, the same relative single estimate and the negligible bound on discarded pairs show that the ordinary and outer collision counts for a uniform subset agree with probability tending to one.

Finally, in the fixed-order factorial moment proof, identical edges are excluded from each factorial moment, leaving only the two kinds of nonzero atypical configurations handled above. Because the typical terms transfer relatively and the atypical terms remain negligible, the factorial moments converge to the same limits as before, yielding the same Poisson limits for the uniform subset model.


\section*{Statements and Declarations}

\paragraph{Use of AI tools.} During the preparation of this manuscript, the author used Claude
Opus~5.5 (Anthropic) to support developing the proofs, writing the
verification code and preparing the manuscript. The author reviewed,
validated and critically assessed all AI-assisted outputs and takes full
responsibility for the content of the manuscript.

\paragraph{Competing interests.} The author has no relevant financial or non-financial
interests to disclose.

\paragraph{Funding.} No funding was received for this work.

\paragraph{Data availability.} The verification code is provided as Online
Resource~1 and at \url{https://github.com/brandonlign/hamming}. No other data
were used.

\begin{thebibliography}{99}
\bibitem{Abdurakhmanov} J. K. Abdurakhmanov, \emph{Distribution of codewords on the faces of a hypercube and new combinatorial identities}, preprint (2025), \href{https://arxiv.org/abs/2506.18494}{arXiv:2506.18494}.
\bibitem{Allikvere} J. Allikvere, \emph{The edge multiset dimension of hypercubes}, preprint (2026), \href{https://arxiv.org/abs/2608.09983}{arXiv:2608.09983}.
\bibitem{BMP} B. Bollob\'as, D. Mitsche and P. Pra{\l}at, \emph{Metric dimension for random graphs}, Electron. J. Combin. 20(4) (2013), \#P1, \href{https://doi.org/10.37236/2639}{doi:10.37236/2639}.
\bibitem{BLM} S. Boucheron, G. Lugosi and P. Massart, \emph{Concentration Inequalities: A Nonasymptotic Theory of Independence}, Oxford University Press, 2013.
\bibitem{CCM} P. Camion, B. Courteau and A. Montpetit, \emph{Metric coset schemes revisited}, Ann. Inst. Fourier (Grenoble) 49 (1999), 829--859, \href{https://doi.org/10.5802/aif.1695}{doi:10.5802/aif.1695}.
\bibitem{CM} D. G. Cantor and W. H. Mills, \emph{Determination of a subset from certain combinatorial properties}, Canad. J. Math. 18 (1966), 42--48, \href{https://doi.org/10.4153/CJM-1966-007-2}{doi:10.4153/CJM-1966-007-2}.
\bibitem{Chvatal} V. Chv\'atal, \emph{Mastermind}, Combinatorica 3 (1983), 325--329, \href{https://doi.org/10.1007/BF02579188}{doi:10.1007/BF02579188}.
\bibitem{EP} A. Eide and P. Pra{\l}at, \emph{Multiset metric dimension of binomial random graphs}, preprint (2025, revised 2026), \href{https://arxiv.org/abs/2507.11686}{arXiv:2507.11686}.
\bibitem{GH} GitHub user \texttt{eliottcassidy2000}, \emph{The edge multiset dimension of the hypercube grows like \(\exp(\Theta(d^{1/3}))\)}, repository \texttt{math}, file \texttt{01-canon/theorems/THM-4534}, commit 0e27abd2aa, 2 October 2026, \href{https://github.com/eliottcassidy2000/math}{github.com/eliottcassidy2000/math}.
\bibitem{ER} P. Erd\H{o}s and A. R\'enyi, \emph{On two problems of information theory}, Magyar Tud. Akad. Mat. Kutat\'o Int. K\"ozl. 8 (1963), 229--243.
\bibitem{FKKY} M. Farhan, S. Klav\v{z}ar, D. Kuziak and I. G. Yero, \emph{Multiset resolvability parameters in graphs: A survey with new results and open problems}, preprint (2026), \href{https://arxiv.org/abs/2607.10311}{arXiv:2607.10311}.
\bibitem{GRTY} R. Gil-Pons, Y. Ram\'irez-Cruz, R. Trujillo-Rasua and I. G. Yero, \emph{Distance-based vertex identification in graphs: The outer multiset dimension}, Appl. Math. Comput. 363 (2019), 124612, \href{https://doi.org/10.1016/j.amc.2019.124612}{doi:10.1016/j.amc.2019.124612}.
\bibitem{ID} A. Hakanen and I. G. Yero, \emph{Complexity and equivalency of multiset dimension and ID-colorings}, Fund. Inform. 191 (2024), 315--330, \href{https://doi.org/10.3233/FI-242185}{doi:10.3233/FI-242185}.
\bibitem{HM} F. Harary and R. A. Melter, \emph{On the metric dimension of a graph}, Ars Combin. 2 (1976), 191--195.
\bibitem{Hoeffding} W. Hoeffding, \emph{Probability inequalities for sums of bounded random variables}, J. Amer. Statist. Assoc. 58 (1963), 13--30, \href{https://doi.org/10.1080/01621459.1963.10500830}{doi:10.1080/01621459.1963.10500830}.
\bibitem{IISN} H. M. Ikhlaq, R. Ismail, H. M. A. Siddiqui and M. F. Nadeem, \emph{A new technique to uniquely identify the edges of a graph}, Symmetry 15 (2023), 762, \href{https://doi.org/10.3390/sym15030762}{doi:10.3390/sym15030762}.
\bibitem{JLR} S. Janson, T. {\L}uczak and A. Ruci\'nski, \emph{Random Graphs}, Wiley, 2000, \href{https://doi.org/10.1002/9781118032718}{doi:10.1002/9781118032718}.
\bibitem{KMSY} A. Kelenc, A. T. Masa Toshi, R. \v{S}krekovski and I. G. Yero, \emph{On metric dimensions of hypercubes}, Ars Math. Contemp. 23 (2023), \#P2.08, \href{https://doi.org/10.26493/1855-3974.2568.55c}{doi:10.26493/1855-3974.2568.55c}.
\bibitem{KTY} A. Kelenc, N. Tratnik and I. G. Yero, \emph{Uniquely identifying the edges of a graph: The edge metric dimension}, Discrete Appl. Math. 251 (2018), 204--220, \href{https://doi.org/10.1016/j.dam.2018.05.052}{doi:10.1016/j.dam.2018.05.052}.
\bibitem{KRR} S. Khuller, B. Raghavachari and A. Rosenfeld, \emph{Landmarks in graphs}, Discrete Appl. Math. 70 (1996), 217--229, \href{https://doi.org/10.1016/0166-218X(95)00106-2}{doi:10.1016/0166-218X(95)00106-2}.
\bibitem{KKY} S. Klav\v{z}ar, D. Kuziak and I. G. Yero, \emph{Further contributions on the outer multiset dimension of graphs}, Results Math. 78 (2023), 50, \href{https://doi.org/10.1007/s00025-022-01829-8}{doi:10.1007/s00025-022-01829-8}.
\bibitem{Several} T. Laihonen, \emph{The metric dimension for resolving several objects}, Inform. Process. Lett. 116 (2016), 694--700, \href{https://doi.org/10.1016/j.ipl.2016.06.002}{doi:10.1016/j.ipl.2016.06.002}.
\bibitem{Laird} L. Laird, R. C. Tillquist, S. Becker and M. E. Lladser, \emph{Resolvability of Hamming graphs}, SIAM J. Discrete Math. 34 (2020), 2063--2081, \href{https://doi.org/10.1137/19M1274511}{doi:10.1137/19M1274511}.
\bibitem{Lindstrom} B. Lindstr\"om, \emph{On a combinatorial problem in number theory}, Canad. Math. Bull. 8 (1965), 477--490, \href{https://doi.org/10.4153/CMB-1965-034-2}{doi:10.4153/CMB-1965-034-2}.
\bibitem{ODonnell} R. O'Donnell, \emph{Analysis of Boolean Functions}, Cambridge University Press, 2014, \href{https://doi.org/10.1017/CBO9781139814782}{doi:10.1017/CBO9781139814782}.
\bibitem{Saenpholphat} V. Saenpholphat, \emph{On multiset dimension in graphs}, Academic SWU 1 (2009), 193--202.
\bibitem{SSV} R. Simanjuntak, P. Siagian and T. Vetr\'ik, \emph{The multiset dimension of graphs}, preprint (2017, revised 2019), \href{https://arxiv.org/abs/1711.00225}{arXiv:1711.00225}.
\bibitem{Slater} P. J. Slater, \emph{Leaves of trees}, Congr. Numer. 14 (1975), 549--559.
\bibitem{TFL} R. C. Tillquist, R. M. Frongillo and M. E. Lladser, \emph{Getting the lay of the land in discrete space: A survey of metric dimension and its applications}, SIAM Rev. 65 (2023), 919--962, \href{https://doi.org/10.1137/21M1409512}{doi:10.1137/21M1409512}.
\end{thebibliography}
\end{document}
